\chapter{椭圆曲线的算术}\label{ch:4}

本章证明的基本定理是有限基定理.

\begin{theorem}[有限基定理] \label{thm:4.0.1} 
对数域 (number field) $K$ 上的椭圆曲线 $E$, 群 $E(K)$ 是有限生成的.
\end{theorem}

换言之, 从 $E$ 上的有限个点出发, 通过作切线与弦可构造出所有点. 该定理在 $K = \mathbb{Q}$ 时由 \cite{mordell1922} 证明, 对一切数域由 Weil 在他的学位论文 (1928) 中证明. 在 $\mathbb{Q}$ 上它也称为 Mordell 定理, 在一般数域上称为 Mordell--Weil 定理. 群 $E(K)$ 称为 $E$ 的 Mordell--Weil 群\index{Mordell-Weil-qun@Mordell-Weil 群(Mordell-Weil group)}, 其秩 (rank) 称为\index{zhi@秩(rank)} $E$ 的秩.

证明该定理的第一步是证明较弱的结果.

\begin{theorem}[弱有限基定理] \label{thm:4.0.2} 
对数域 $K$ 上的每条椭圆曲线 $E$ 与整数 $n\ge 1$, $E(K)/nE(K)$ 是有限的.
\end{theorem}

显然, 对 Abel 群 $M$,
\[
M \text{ 有限生成} \implies M/nM \text{ 对所有 } n\ge 1 \text{ 有限},
\]
但逆不真. 例如 $(\mathbb{Q},+)$ 有性质 $\mathbb{Q} = n\mathbb{Q}$ (对所有 $n\ge 1$), 但它不有限生成, 因为任何有限生成子群中元素有公共分母.

我们在 $K = \mathbb{Q}$ 情形概述弱定理蕴含强定理的证明. 定义 $E(\mathbb{Q})$ 中某点 $P$ 的高度为它作为 $\mathbb{P}^2(\mathbb{Q})$ 中元素的高度 (见I~\ref{thm:1.2.3}). 假设 $E(\mathbb{Q})/2E(\mathbb{Q})$ 有限, 令 $P_1, \dots, P_s\in E(\mathbb{Q})$ 为 $E(\mathbb{Q})/2E(\mathbb{Q})$ 中元素的一组代表元. 则每个点 $Q\in E(\mathbb{Q})$ 可写成
\[
Q = P_i + 2Q'
\]
其中 $Q'\in E(\mathbb{Q})$. 高度 $h$ 具有性质: 存在常数 $C$ 使 $h(Q) > C$ 蕴含 $h(Q') < h(Q)$.

现在考虑 $Q\in E(\mathbb{Q})$. 若 $h(Q) > C$, 则 $h(Q') < h(Q)$. 若还有 $h(Q') > C$, 则对 $Q'$ 重复该论证. 如此继续, 得到
\[
Q = P_i + 2Q' = P_i + 2(P_{i'} + 2Q'') = \cdots, \qquad
h(Q) > h(Q') > h(Q'') > \cdots.
\]
最终 $h(Q^{(n)})\le C$, 而 $Q$ 将是 $P_i$ 的线性组合与 $2^nQ^{(n)}$ 的和. 故 $P_i$ 连同高度 $\le C$ 的有限点集生成 $E(\mathbb{Q})$.

\begin{remarkstar}
(a) 上一段的论证称为``下降证明'' (proof by descent). 费马用它回答了 Pisano 的一个古老问题: 求所有有理平方 $x^2, y^2$ 使 $x^2 + y^2$ 与 $x^2 - y^2$ 都是平方. 若 $(x,y)$ 是满足 $y\neq 0$ 的一个解, 费马证明存在另一个解 $(x', y')$ 满足 $0 < |y'| < |y|$, 若 $(x,y)$ 选得使 $|y|$ 最小则矛盾 (\cite[p.~207]{cassels1966}). 费马通常被归功于开创下降证明, 但在某种意义上这可追溯到古希腊人. 考虑 $Y^2 = 2X^2$ 无整数解的证明. 定义整数对 $(m,n)$ 的高度为 $\max(|m|, |n|)$. 通常的论证表明, 若 $(m,n)$ 是方程的一个解, 则存在另一个高度更小的解, 又得到矛盾.

(b) Weil 在其论文中证明: 对任意亏格 $g$ 曲线 $C$ 于数域 $K$ 上, 次数为 $0$ 的除子模线性等价的群 $\operatorname{Pic}^0(C)$ 有限生成.\footnote{Weil 的叙述有时粗略, 且不用高度. 在他写``il est ais\'e de combler cette lacune'' (容易填补这个缺口) 并随后提出不充分论证的地方, 今天最自然的是用高度理论. 见 Weil 对他的 Collected Papers 中文章的评论. Hadamard 著名地劝阻 Weil 不要发表他的论文, 直到他``完成''了 Mordell 猜想的证明.} 对椭圆曲线, $\operatorname{Pic}^0(C) = C(K)$ (见I, \ref{thm:1.4.10}); 对一般曲线 $\operatorname{Pic}^0(C) = J_C(K)$, 其中 $J_C$ 是 $C$ 的 Jacobi 簇. 因此, 一旦数域上曲线 $C$ 的 Jacobi 簇 $J_C$ 被定义, Weil 的结果就表明 $J_C(K)$ 有限生成. 对数域 $K$ 上任意 阿贝尔簇 $A$, $A$ 上的每条曲线 $C$ 定义同态 $J_C\to A$, 当 $C$ 足够一般时它是满射 (例如 \cite[10.1]{milne1986b}). 由此 $A$ 是 $J_C$ 的直接因子 (例如 \cite[12.1]{milne1986a}), 故 $A(K)$ 有限生成. 这也称为 Mordell--Weil 定理.
\end{remarkstar}

\section{群上同调}\label{sec:4.1}

在证明弱有限基定理, 以及稍后研究 Tate--Shafarevich 群\index{Tate-Shafarevich-qun@Tate-Shafarevich 群(Tate-Shafarevich group)}时, 我们将用到一些群上同调.

\subsection*{群的上同调}

设 $G$ 为群, $M$ 为 Abel 群. $G$ 在 $M$ 上的作用 (action) 是映射 $G\times M\to M$ 满足:
\begin{enumerate}[label=\textup{(\alph*)}]
\item $\sigma(m + m') = \sigma m + \sigma m'$, 对所有 $\sigma\in G$, $m, m'\in M$;
\item $(\sigma\tau)(m) = \sigma(\tau m)$, 对所有 $\sigma, \tau\in G$, $m\in M$;
\item $1_G m = m$, 对所有 $m\in M$.
\end{enumerate}
因此, 给出 $G$ 在 $M$ 上的作用等价于给出同态 $G\to\Aut(M)$. $G$-模\index{G-mo@G-模(G-module)} (G-module) 是带 $G$ 作用的 Abel 群.

\begin{example} \label{thm:4.1.1} 
设 $L$ 为域 $K$ 的有限 Galois 扩张, Galois 群为 $G$, $E$ 为 $K$ 上的椭圆曲线. 则 $L$, $L^{\times}$ 与 $E(L)$ 都有显然的 $G$-作用.
\end{example}

对 $G$-模 $M$, 定义
\[
H^0(G, M) = M^G = \{m\in M \mid \sigma m = m, \ \forall\sigma\in G\}.
\]
在上述例子中,
\[
H^0(G, L) = K, \qquad H^0(G, L^{\times}) = K^{\times}, \qquad H^0(G, E(L)) = E(K).
\]
交叉同态 (crossed homomorphism)\index{jiaochatongtai@交叉同态(crossed homomorphism)} 是满足
\[
f(\sigma\tau) = f(\sigma) + \sigma f(\tau), \qquad \forall\sigma,\tau\in G
\]
的映射 $f\colon G\to M$. 注意该条件蕴含 $f(1) = f(1\cdot 1) = f(1) + f(1)$, 故 $f(1) = 0$. 对每个 $m\in M$, 令
\[
f(\sigma) = \sigma m - m, \qquad \forall\sigma\in G,
\]
得到交叉同态. 这样的交叉同态称为主的 (principal).\index{jiaochatongtai@交叉同态(crossed homomorphism)!zhu@主(dominant)} 两个交叉同态的和与差仍是交叉同态, 两个主交叉同态的和与差仍是主的. 因此可定义
\[
H^1(G, M) = \frac{\{\text{交叉同态}\}}{\{\text{主交叉同态}\}}
\]
(商 Abel 群). 对 $n > 1$ 还有上同调群 $H^n(G, M)$, 但我们将不需要它们.

\begin{example} \label{thm:4.1.2} 
当 $G$ 平凡作用在 $M$ 上 (即对所有 $\sigma\in G$, $m\in M$ 有 $\sigma m = m$) 时, 交叉同态就是同态, 而每个主交叉同态为零. 故 $H^1(G, M) = \operatorname{Hom}(G, M)$.
\end{example}

\begin{proposition} \label{thm:4.1.3} 
设 $L$ 为 $K$ 的有限 Galois 扩张, Galois 群为 $G$, 则 $H^1(G, L^{\times}) = 0$, 即每个交叉同态 $G\to L^{\times}$ 都是主的.
\end{proposition}

\begin{proof}
设 $f$ 为交叉同态 $G\to L^{\times}$. 用乘法记号, 这意味着
\[
f(\sigma\tau) = f(\sigma)\cdot\sigma(f(\tau)), \qquad \sigma, \tau\in G,
\]
且我们要找 $\lambda\in L^{\times}$ 使 $f(\sigma) = \sigma\lambda/\lambda$ 对所有 $\sigma\in G$. 因 $f(\tau)$ 非零, Dedekind 特征标独立性定理 (\cite[5.14]{milne-ft}) 表明
\[
\sum_{\tau\in G}f(\tau)\tau\colon L\to L
\]
不是零映射, 即存在 $\alpha\in L$ 使
\[
\beta \overset{\text{def}}{=} \sum_{\tau\in G}f(\tau)\tau\alpha\neq 0.
\]
但对 $\sigma\in G$,
\begin{align*}
\sigma\beta &= \sum_{\tau\in G}\sigma(f(\tau))\cdot\sigma\tau(\alpha) = \sum_{\tau\in G}f(\sigma)^{-1}\cdot f(\sigma\tau)\cdot\sigma\tau(\alpha)\\
&= f(\sigma)^{-1}\sum_{\tau\in G}f(\sigma\tau)\cdot\sigma\tau(\alpha) = f(\sigma)^{-1}\beta,
\end{align*}
故 $f(\sigma) = \beta/\sigma\beta = \sigma(\beta^{-1})/\beta^{-1}$.
\end{proof}

\begin{corollary} \label{thm:4.1.4} 
点 $P = (x_0:\dots:x_n)\in\mathbb{P}^n(L)$ 被 $G$ 固定, 当且仅当它由 $K$ 中的 $n+1$ 元组表示.
\end{corollary}

\begin{proof}
若 $\sigma P = P$ ($\sigma\in G$), 则
\[
\sigma(x_0, \dots, x_n) = c(\sigma)(x_0, \dots, x_n)
\]
对某个 $c(\sigma)\in L^{\times}$. 映射 $\sigma\mapsto c(\sigma)$ 是交叉同态, 故 $c(\sigma) = \sigma c^{-1}/c^{-1}$ 对某个 $c\in L^{\times}$. 现在
\[
\sigma(cx_0, \dots, cx_n) = (cx_0, \dots, cx_n),
\]
故诸 $cx_i$ 位于 $K$ 中.
\end{proof}

\begin{proposition} \label{thm:4.1.5} 
对 $G$-模的任何正合列
\[
0\to M\to N\to P\to 0,
\]
有典范正合列
\[
0\to H^0(G, M)\to H^0(G, N)\to H^0(G, P)
\xrightarrow{\delta} H^1(G, M)\to H^1(G, N)\to H^1(G, P).
\]
\end{proposition}

\begin{proof}
映射 $\delta$ 定义如下: 设 $p\in P^G$; 存在映射到 $p$ 的 $n\in N$, 且对所有 $\sigma\in G$ 有 $\sigma n - n\in M$; 映射 $\sigma\mapsto\sigma n - n\colon G\to M$ 是交叉同态, 其类定义为 $\delta(p)$. 映射到 $p$ 的另一个 $n'$ 产生与第一个相差主交叉同态 $\sigma\mapsto\sigma(n'-n) - (n'-n)$ 的交叉同态, 故 $\delta(p)$ 良定义. 证明的其余部分按惯例 (所有不熟悉群上同调的读者都应写出来).
\end{proof}

设 $H$ 为 $G$ 的子群, $M$ 为 $G$-模. $G$ 的交叉同态限制为 $H$ 的交叉同态, 由此得到``限制'' (restriction) 同态
\[
\operatorname{Res}\colon H^1(G, M)\to H^1(H, M).
\]
现在假设 $H$ 在 $G$ 中指标有限, 设 $S$ 为 $G/H$ 的代表元集, 故 $G = \bigsqcup_{s\in S}sH$. 设 $f$ 为 $H$ 的交叉同态. 则 $\sigma\mapsto\sum_{s\in S}t\cdot f(t^{-1}\cdot\sigma s)$ (其中 $t$ 是满足 $tH = \sigma sH$ 的 $S$ 中唯一元素) 是 $G$ 的交叉同态, 由此得到``余限制'' (corestriction) 同态
\[
\operatorname{Cor}\colon H^1(H, M)\to H^1(G, M).
\]
复合 $\operatorname{Cor}\circ\operatorname{Res} = (G:H)$. 见 \cite[VIII, \S 2]{serre1979}.

\begin{proposition} \label{thm:4.1.6} 
若 $G$ 的阶为 $m$, 则 $mH^1(G, M) = 0$.
\end{proposition}

\begin{proof}
在上面讨论中取 $H = 1$, 我们看到 $H^1(G, M)$ 上乘以 $m$ 穿过 $H^1(1, M) = 0$.
\end{proof}

\begin{remark} \label{thm:4.1.7} 
设 $H$ 为群 $G$ 的正规子群, $M$ 为 $G$-模. 则 $M^H$ 是 $G/H$-模, 交叉同态 $f\colon G/H\to M^H$ 由复合定义 $G$ 的交叉同态:
\[\begin{tikzcd}
	G & M \\
	{G/H} & {M^H.}
	\arrow[dashed, from=1-1, to=1-2]
	\arrow[from=1-1, to=2-1]
	\arrow["f", from=2-1, to=2-2]
	\arrow[hook, from=2-2, to=1-2]
\end{tikzcd}\]
由此得到``膨胀'' (inflation) 同态
\[
\operatorname{Inf}\colon H^1(G/H, M^H)\to H^1(G, M),
\]
且容易验证序列
\[
0\to H^1(G/H, M^H)\xrightarrow{\operatorname{Inf}}H^1(G, M)\xrightarrow{\operatorname{Res}}H^1(H, M)
\]
正合.
\end{remark}

\subsection*{无限 Galois 群的上同调}

设 $k$ 为完美域, $k^{\mathrm{al}}$ 为 $k$ 的代数闭包. 固定 $k$ 中元素的 $k^{\mathrm{al}}$ 的自同构群 $G$ 有自然拓扑, 称为 Krull 拓扑 (Krull topology)\index{Krull@Krull 拓扑(Krull topology)}, 其中子群开当且仅当它由固定 $k^{\mathrm{al}}$ 中 $k$ 的某个有限扩张元素的那些自同构组成. 配 Krull 拓扑时, $G$ 称为 $k^{\mathrm{al}}/k$ 的 Galois 群 (Galois group). $G$ 的开子群构成恒等元的邻域基. 开子群总是闭的, 故开子群的任何交闭; 反之, 每个闭子群是开子群的交. 群 $G$ 紧致, 故 $G$ 的开子群指标有限, 且 $G$ 的连续映射一致连续. 通常的 Galois 理论延拓给出中间域 $K$ ($k\subset K\subset k^{\mathrm{al}}$) 与 $G$ 的闭子群之间的一一对应, 其中在 $k$ 上有限次的域对应 $G$ 的开子群. 对 $k^{\mathrm{al}}$ 中 $k$ 的每个有限 Galois 扩张 $K$, 映射 $\Gal(k^{\mathrm{al}}/k)\to\Gal(K/k)$ 是带开正规核的满射, 且这些映射把 $\Gal(k^{\mathrm{al}}/k)$ 实现为逆极限
\[
\Gal(k^{\mathrm{al}}/k) = \varprojlim_K\Gal(K/k).
\]
细节见 FT, \S7.

$G$-模 $M$ 称为离散的 (discrete)\index{lisanmo@离散模(discrete module)}, 若映射 $G\times M\to M$ 相对于 $M$ 的离散拓扑与 $G$ 的 Krull 拓扑连续. 这等价于要求
\[
M = \bigcup_H M^H, \qquad H \text{ 在 } G \text{ 中开},
\]
即要求 $M$ 的每个元素被固定 $k$ 的某个有限扩张的 $G$ 的子群固定. 例如 $k^{\mathrm{al}}$, $k^{\mathrm{al}\times}$ 与 $E(k^{\mathrm{al}})$ 都是离散 $G$-模, 因为
\[
k^{\mathrm{al}} = \bigcup_K K, \qquad k^{\mathrm{al}\times} = \bigcup_K K^{\times}, \qquad E(k^{\mathrm{al}}) = \bigcup_K E(K),
\]
其中每个情形中并集遍历 $k^{\mathrm{al}}$ 中 $k$ 的有限扩张 $K$. 当 $M$ 离散时, 交叉同态 $f\colon G\to M$ 连续当且仅当它在 $G$ 的某个开正规子群 $H$ 的陪集上为常数, 从而由 $G/H\to M^H$ 的交叉同态经膨胀产生. 每个主交叉同态连续, 因为 $M$ 的每个元素被 $G$ 的开正规子群固定.

对无限 Galois 群 $G$ 与离散 $G$-模 $M$, 定义 $H^1(G, M)$ 为连续交叉同态 $f\colon G\to M$ 模主交叉同态子群. 用此定义,
\[
H^1(G, M) = \varinjlim_H H^1(G/H, M^H),
\]
其中 $H$ 遍历 $G$ 的开正规子群. 明确地, 这意味着:
\begin{enumerate}[label=\textup{(\alph*)}]
\item $H^1(G, M)$ 是膨胀映射 $\operatorname{Inf}\colon H^1(G/H, M^H)\to H^1(G, M)$ ($H$ 遍历 $G$ 的开正规子群) 的像之并;
\item 若 $\xi\in H^1(G, M)$ 且 $H$ 使 $\xi$ 来自 $H^1(G/H, M^H)$, 则 $\xi$ 在 $\operatorname{Res}\colon H^1(G, M)\to H^1(H, M)$ 下的像为零.
\end{enumerate}

\begin{example} \label{thm:4.1.8} 
\begin{enumerate}[label=\textup{(\alph*)}]
\item 命题 \ref{thm:4.1.3} 表明
\[
H^1(G, k^{\mathrm{al}\times}) = \varinjlim_K H^1(\Gal(K/k), K^{\times}) = 0.
\]
\item 对域 $L$ 与整数 $n > 1$, 令
\[
\mu_n(L) = \{\zeta\in L^{\times} \mid \zeta^n = 1\}.
\]
由正合列
\[
1 \to \mu_n(k^{\mathrm{al}}) \to k^{\mathrm{al}\times} \xrightarrow{n} k^{\mathrm{al}\times} \to 1,
\]
得上同调群的正合列
\[
1 \to \mu_n(k) \to k^{\times} \xrightarrow{n} k^{\times} \to H^1(G, \mu_n(k^{\mathrm{al}})) \to 1,
\]
故有典范同构
\[
H^1(G, \mu_n(k^{\mathrm{al}})) \simeq k^{\times}/k^{\times n}.
\]
当 $k$ 是数域时, 该群无限. 例如, 数
\[
(-1)^{\varepsilon(\infty)}\prod_{p\text{ 素数}}p^{\varepsilon(p)},
\]
(其中每个指数为 $0$ 或 $1$ 且除有限多个外全为零) 构成 $\mathbb{Q}^{\times}/\mathbb{Q}^{\times2}$ 的元素的一组代表元, 故它是 $\mathbb{F}_2$ 上无限维向量空间.
\item 若 $G$ 平凡作用在 $M$ 上, 则 $H^1(G, M)$ 是连续同态 $\alpha\colon G\to M$ 的集合. 因 $M$ 离散, 这样的 $\alpha$ 的核 $H$ 开, 且 $\alpha$ 定义单同态 $\Gal(K/k)\to M$, 其中 $K$ 是 $H$ 的不动域.
\end{enumerate}
\end{example}


对 $k$ 上的椭圆曲线 $E$, 我们把 $H^i(\Gal(k^{\mathrm{al}}/k), E(k^{\mathrm{al}}))$ 简写为 $H^i(k, E)$.

\begin{aside} \label{thm:4.1.9}
设 $E$ 为 $\mathbb{Q}$ 上的椭圆曲线, $\mathbb{Q}^{\mathrm{al}}$ 为 $\mathbb{Q}$ 在 $\mathbb{C}$ 中的代数闭包. 选取 $\mathbb{Q}$ 的 $p$-进赋值到 $\mathbb{Q}^{\mathrm{al}}$ 的一个延拓, 令 $\mathbb{Q}_p^{\mathrm{al}}$ 为 $\mathbb{Q}_p$ 在由该赋值定义的 $\mathbb{Q}^{\mathrm{al}}$ 的完备化中的代数闭包. 则 $\mathbb{Q}_p^{\mathrm{al}}$ 是 $\mathbb{Q}_p$ 的代数闭包, 且有交换图
\begin{equation} \label{eq:27}
\begin{tikzcd}
{\mathbb{Q}^{\mathrm{al}}} & {\mathbb{Q}_p^{\mathrm{al}}} \\
{\mathbb{Q}} & {\mathbb{Q}_p}
\arrow[hook, from=2-1, to=1-1]
\arrow[hook, from=2-2, to=1-2]
\arrow[hook, from=2-1, to=2-2]
\arrow[hook, from=1-1, to=1-2]
\end{tikzcd}
\tag{27}
\end{equation}
$\Gal(\mathbb{Q}_p^{\mathrm{al}}/\mathbb{Q}_p)$ 在 $\mathbb{Q}^{\mathrm{al}}\subset\mathbb{Q}_p^{\mathrm{al}}$ 上的作用定义同态
\[
\Gal\left(\mathbb{Q}_p^{\mathrm{al}}/\mathbb{Q}_p\right)\to\Gal\left(\mathbb{Q}^{\mathrm{al}}/\mathbb{Q}\right).
\]
于是每个交叉同态 $\Gal(\mathbb{Q}^{\mathrm{al}}/\mathbb{Q})\to E(\mathbb{Q}^{\mathrm{al}})$ (经复合) 定义交叉同态 $\Gal(\mathbb{Q}_p^{\mathrm{al}}/\mathbb{Q}_p)\to E(\mathbb{Q}_p^{\mathrm{al}})$. 这样得到同态
\[
H^1(\mathbb{Q}, E)\to H^1(\mathbb{Q}_p, E),
\]
它与图~\eqref{eq:27} 的选取无关. 对 $\mu_n$ 与 $E_n$ 的上同调群有类似注记. 见 \cite[VII, \S 5, Proposition 3]{serre1979}. 在下面的 \S\ref{sec:4.7} 中, 我们将给出这些``局部化''同态更自然的几何解释.
\end{aside}

\begin{example} \label{thm:4.1.10}
设 $k$ 为完美域, $\Gamma$ 为其 Galois 群. 设 $C$ 为 $k$ 上曲线, $\bar{C}$ 为 $k^{\mathrm{al}}$ 上相应的曲线. 我们考察 $\operatorname{Pic}(C) \mathrel{=_{\mathrm{def}}} \operatorname{Pic}(\bar{C})^{\Gamma}$ 是否是 $C$ 上除子模主除子得到的群. 由 $\Gamma$-模的正合列
\[
0\to k^{\mathrm{al}}(\bar{C})^{\times}/k^{\mathrm{al}\times}\to\operatorname{Div}(\bar{C})\to\operatorname{Pic}(\bar{C})\to 0
\]
得上同调正合列
\[
\left(k^{\mathrm{al}}(\bar{C})^{\times}/k^{\mathrm{al}\times}\right)^{\Gamma}\to\operatorname{Div}(\bar{C})^{\Gamma}\to\left(\operatorname{Pic}(\bar{C})\right)^{\Gamma}\to H^1(\Gamma, k^{\mathrm{al}}(\bar{C})^{\times}/k^{\mathrm{al}\times}),
\]
而由正合列
\[
0\to k^{\mathrm{al}\times}\to k^{\mathrm{al}}(\bar{C})^{\times}\to k^{\mathrm{al}}(\bar{C})^{\times}/k^{\mathrm{al}\times}\to 0,
\]
(用~\ref{thm:4.1.8}) 得上同调正合列
\[
0\to k^{\times}\to k(C)^{\times}\to\left(k^{\mathrm{al}}(\bar{C})^{\times}/k^{\mathrm{al}\times}\right)^{\Gamma}\to 0,\qquad
0\to H^1(\Gamma, k^{\mathrm{al}}(\bar{C})^{\times}/k^{\mathrm{al}\times})\to H^2(\Gamma, k^{\mathrm{al}\times}).
\]
当用后两个序列替换前一上同调序列中的项, 并用 $\operatorname{Div}(\bar{C})^{\Gamma} = \operatorname{Div}(C)$ (见 I, \S\ref{sec:1.5}) 时, 得正合列
\[
0\to k^{\times}\to k(C)^{\times}\to\operatorname{Div}(C)\to\operatorname{Pic}(C)\to H^2(\Gamma, k^{\mathrm{al}\times}).
\]
现 $H^2(\Gamma, k^{\mathrm{al}\times})$ 是 $k$ 的 Brauer 群\index{Brauer@Brauer 群(Brauer group)} (\cite[X, \S 5]{serre1979}). 于是, 若 $k$ 的 Brauer 群为零, 例如 $k$ 有限时 (Wedderburn 定理; \cite[X, \S 7]{serre1979}), 则 $\operatorname{Pic}(C)$ 是 $C$ 上除子类的群.
\end{example}

\begin{remark} \label{thm:4.1.11}
若 $k$ 不完美, 则 $k^{\mathrm{al}}$ 不是 $k$ 的 Galois 扩张; 特别地, $k$ 不是 $k^{\mathrm{al}}/k$ 的自同构群的不动域. 但当把 $k^{\mathrm{al}}$ 换成 $k$ 的可分闭包 $k^{\mathrm{sep}}$ 时, 上面的讨论本质不变, 只是在~\ref{thm:4.1.8}b 中, 必须取 $n$ 与 $k$ 的特征互素, 才能使 $k^{\mathrm{sep}\times}\xrightarrow{n}k^{\mathrm{sep}\times}$ 是满射.
\end{remark}

\section{Selmer 群\index{Selmer-qun@Selmer 群(Selmer group)}与 Tate--Shafarevich 群}\label{sec:4.2}

令 $\mathbb{Q}_\infty = \mathbb{R}$.

\begin{lemma} \label{thm:4.2.1} 
对代数闭域 $k$ 上的每条椭圆曲线 $E$ 与整数 $n\ge 1$, 映射 $P\mapsto nP\colon E(k)\to E(k)$ 是满射.
\end{lemma}

\begin{proof}
当 $k = \mathbb{C}$ 时, 这由III, \ref{thm:3.3.10} 推出. 若 $k$ 特征为零, 则如 III, \ref{thm:3.3.12} 的证明, 可设 $k\subset\mathbb{C}$. 给定点 $P\in E(k)$, 证存在点 $Q$ 使 $nQ = P$ 等价于证某些 $X$ 与 $Y$ 的多项式有公共零点 (见III, \ref{thm:3.3.14}). 因这些方程在 $\mathbb{C}$ 中有解, 多项式生成 $k[X, Y]$ 的真理想, 由零点定理在 $k$ 中有解.

对任意代数闭 $k$, 利用对每个 $m\neq\operatorname{char}(k)$, 存在 $E(k)$ 中阶为 $m$ 的非零点 $P$ (见II, \ref{thm:2.6.3}). 特别地, 存在阶与 $n$ 互素的非零点 $P\in E(k)$. 则 $nP\neq 0$, 故正则映射 $n\colon E\to E$ 非常数, 从而是支配映射 (I \ref{thm:1.4.23}a), 因而是满射 (I\ref{thm:1.4.23}b).
\end{proof}

设 $E$ 为 $\mathbb{Q}$ 上的椭圆曲线. 对整数 $n > 1$, 令 $E_n$ 为满足
\[
E_n(\mathbb{Q}^{\mathrm{al}}) = E(\mathbb{Q}^{\mathrm{al}})_n \overset{\text{def}}{=} \operatorname{Ker}(E(\mathbb{Q}^{\mathrm{al}})\xrightarrow{n}E(\mathbb{Q}^{\mathrm{al}}))
\]
的 $E$ 的代数子簇. 由引理得正合列
\[
0\to E_n(\mathbb{Q}^{\mathrm{al}})\to E(\mathbb{Q}^{\mathrm{al}})\xrightarrow{n}E(\mathbb{Q}^{\mathrm{al}})\to 0
\]
与上同调正合列
\[
0\to E_n(\mathbb{Q})\to E(\mathbb{Q})\xrightarrow{n}E(\mathbb{Q})\to H^1(\mathbb{Q}, E_n)\to H^1(\mathbb{Q}, E)\xrightarrow{n}H^1(\mathbb{Q}, E),
\]
从中取出序列
\begin{equation} \label{eq:28}
\boxed{0\to E(\mathbb{Q})/nE(\mathbb{Q})\to H^1(\mathbb{Q}, E_n)\to H^1(\mathbb{Q}, E)_n\to 0.}
\tag{28}
\end{equation}
{\emergencystretch=3em 这里, 如常, $H^1(\mathbb{Q}, E)_n$ 是被 $n$ 杀掉的 $H^1(\mathbb{Q}, E)$ 的子群. 若 $H^1(\mathbb{Q}, E_n)$ 有限, 则能推出 $E(\mathbb{Q})/nE(\mathbb{Q})$ 有限, 但它未必. 例如, 若 $E$ 上所有 $2$ 阶点的坐标都在 $\mathbb{Q}$ 中, 使 $\Gal(\mathbb{Q}^{\mathrm{al}}/\mathbb{Q})$ 平凡作用在 $E_2(\mathbb{Q}^{\mathrm{al}})\cong(\mathbb{Z}/2\mathbb{Z})^2$ 上, 则\par}
\[
H^1(\mathbb{Q}, E_2)\cong H^1(\mathbb{Q}, \mu_2\times\mu_2)\cong\mathbb{Q}^{\times}/\mathbb{Q}^{\times2}\times\mathbb{Q}^{\times}/\mathbb{Q}^{\times2},
\]
我们在 \ref{thm:4.1.8}b 中看到这是无限的. 相反, 我们用 $H^1(\mathbb{Q}, E_n)$ 的某个子群 (Selmer 群 $S$) 替换它, 它足够大以包含 $E(\mathbb{Q})/nE(\mathbb{Q})$ 的像, 但我们可以证明它有限.

把 $E$ 视为 $\mathbb{Q}_p$ 上的椭圆曲线, 得到 \eqref{eq:28} 的类似正合列, 以及交换图 (见 \ref{thm:4.1.9}):
\[
\begin{tikzcd}[column sep=large]
0 \ar[r] &
E(\mathbb{Q})/nE(\mathbb{Q}) \ar[d] \ar[r] &
H^1(\mathbb{Q}, E_n) \ar[d] \ar[r] &
H^1(\mathbb{Q}, E)_n \ar[d] \ar[r] &
0 \\
0 \ar[r] &
E(\mathbb{Q}_p)/nE(\mathbb{Q}_p) \ar[r] &
H^1(\mathbb{Q}_p, E_n) \ar[r] &
H^1(\mathbb{Q}_p, E)_n \ar[r] &
0
\end{tikzcd}
\]
若 $\xi\in H^1(\mathbb{Q}, E_n)$ 来自 $E(\mathbb{Q})$ 的元素, 则当然它的像 $\xi_p\in H^1(\mathbb{Q}_p, E_n)$ 来自 $E(\mathbb{Q}_p)$ 的元素. 这提示定义
\begin{align*}
S^{(n)}(E/\mathbb{Q}) &= \{\xi\in H^1(\mathbb{Q}, E_n) \mid \text{对所有 } p, \ \xi_p \text{ 来自 } E(\mathbb{Q}_p)\}\\
&= \operatorname{Ker}\Bigl(H^1(\mathbb{Q}, E_n)\to\prod_{p = 2, 3, 5, \dots, \infty}H^1(\mathbb{Q}_p, E)\Bigr).
\end{align*}
群 $S^{(n)}(E/\mathbb{Q})$ 称为 Selmer 群. 本着同样的精神, 定义 Tate--Shafarevich 群为
\[
\Sha(E/\mathbb{Q}) = \operatorname{Ker}\Bigl(H^1(\mathbb{Q}, E)\to\prod_{p = 2, 3, 5, \dots, \infty}H^1(\mathbb{Q}_p, E)\Bigr).
\]
我们需要一个初等引理.

\begin{lemma} \label{thm:4.2.2} 
任何 Abel 群同态对
\[
A\xrightarrow{\alpha}B\xrightarrow{\beta}C
\]
给出正合 (核--余核) 序列\index{heyuhexulie@核--余核序列(kernel-cokernel sequence)}
\[
0\to\operatorname{Ker}(\alpha)\to\operatorname{Ker}(\beta\circ\alpha)\xrightarrow{\alpha}\operatorname{Ker}(\beta)
\xrightarrow{}\operatorname{Coker}(\alpha)\xrightarrow{\beta}\operatorname{Coker}(\beta\circ\alpha)\to\operatorname{Coker}(\beta)\to 0.
\]
\end{lemma}

\begin{proof}
直接证明, 或把蛇引理应用于图
\[\begin{tikzcd}
	& A & B & {\Coker(\alpha)} & 0 \\
	0 & C & C & 0.
	\arrow["\alpha"', from=1-2, to=1-3]
	\arrow["{\beta\circ\alpha}", from=1-2, to=2-2]
	\arrow[from=1-3, to=1-4]
	\arrow["\beta", from=1-3, to=2-3]
	\arrow[from=1-4, to=1-5]
	\arrow[from=1-4, to=2-4]
	\arrow[from=2-1, to=2-2]
	\arrow["\id", from=2-2, to=2-3]
	\arrow[from=2-3, to=2-4]
\end{tikzcd}\]
\end{proof}

把引理应用于映射
\[
H^1(\mathbb{Q}, E_n)\to H^1(\mathbb{Q}, E)_n\to\prod_{p = 2, 3, \dots, \infty}H^1(\mathbb{Q}_p, E)_n,
\]
得到基本正合列
\begin{equation} \label{eq:29}
0\to E(\mathbb{Q})/nE(\mathbb{Q})\to S^{(n)}(E/\mathbb{Q})\to\Sha(E/\mathbb{Q})_n\to 0.
\tag{29}
\end{equation}
我们将通过证明 $S^{(n)}(E/\mathbb{Q})$ 有限来证明 $E(\mathbb{Q})/nE(\mathbb{Q})$ 有限.

注意 $\Sha(E/\mathbb{Q})$ 是挠群 (见 \ref{thm:4.5.5}). 稍后我们将给出 $\Sha(E/\mathbb{Q})$ 的几何解释, 它表明 $\Sha(E/\mathbb{Q})$ 提供了 $\mathbb{Q}$ 上亏格 $1$ 曲线 Hasse 原理失败的量度. 我们还将看到这些定义可推广到所有数域上的椭圆曲线.

\begin{remarkstar}
群 $H^1(k, E)$ 也称为 Weil--Ch\^atelet 群, 记为 $\operatorname{WC}(E/k)$ (\cite{lang1958}; \cite{tate1958}). Selmer 群由 \cite{cassels1962a} 命名, 他同时承担了用西里尔字母 $\Sha$ 表示 Tate--Shafarevich 群的责任 (\cite[p.~109]{cassels1991}),\footnote{他声称这是他对该主题``最持久的贡献''.}
\end{remarkstar}

\section{Selmer 群的有限性}\label{sec:4.3}

本节证明以下定理.

\begin{theorem} \label{thm:4.3.1} 
对 $\mathbb{Q}$ 上的每条椭圆曲线 $E$ 与整数 $n\ge 1$, Selmer 群 $S^{(n)}(E/\mathbb{Q})$ 是有限的.
\end{theorem}

事实上, 我们证明任意数域 $L$ 上的相应陈述.

\subsection*{预备}

\begin{lemma} \label{thm:4.3.2} 
设 $E$ 为 $\mathbb{Q}_p$ 上具有好约化的椭圆曲线, $n$ 为与 $p$ 互素的整数. 点 $P\in E(\mathbb{Q}_p)$ 对某个 $Q\in E(\mathbb{Q}_p)$ 具有 $nQ$ 的形式, 当且仅当它的像 $\tilde{P}\in\tilde{E}(\mathbb{F}_p)$ 对某个 $\tilde{Q}\in\tilde{E}(\mathbb{F}_p)$ 具有 $n\tilde{Q}$ 的形式.
\end{lemma}

\begin{proof}
因 $P\mapsto\tilde{P}$ 是同态, 必要性显然; 充分性由图追查
\[\begin{tikzcd}
	0 & {E^1(\mathbb{Q}_p)} & {E(\mathbb{Q}_p)} & {\tilde{E}(\mathbb{F}_p)} & 0 \\
	0 & {E^1(\mathbb{Q}_p)} & {E(\mathbb{Q}_p)} & {\tilde{E}(\mathbb{F}_p)} & 0
	\arrow[from=1-1, to=1-2]
	\arrow[from=1-2, to=1-3]
	\arrow["n", from=1-2, to=2-2]
	\arrow["\simeq"', draw=none, from=1-2, to=2-2]
	\arrow[from=1-3, to=1-4]
	\arrow["n", from=1-3, to=2-3]
	\arrow[from=1-4, to=1-5]
	\arrow["n", from=1-4, to=2-4]
	\arrow[from=2-1, to=2-2]
	\arrow[from=2-2, to=2-3]
	\arrow[from=2-3, to=2-4]
	\arrow[from=2-4, to=2-5]
\end{tikzcd}\]
推出, 其中第一个垂直箭头是同构 (II\ref{thm:2.4.2}). 详细地, 设 $P\in E(\mathbb{Q}_p)$ 使 $\tilde{P} = n\tilde{Q}$ ($Q\in E(\mathbb{Q}_p)$); 则 $P - nQ$ 在 $\tilde{E}(\mathbb{F}_p)$ 中映射为零, 故位于 $E^1(\mathbb{Q}_p)$ 中. 因此 $P - nQ = nQ'$ 对某个 $Q'\in E^1(\mathbb{Q}_p)$, 从而 $P = n(Q + Q')$.
\end{proof}

我们将用到一些局部代数数论. 对 $\mathbb{Q}_p$ 的有限扩张 $K$, $\mathbb{Z}_p$ 在 $K$ 中的整闭包 $\mathcal{O}_K$ 是带唯一极大理想 $(\varpi)$ 的主理想整环. 故 $p = \text{单位}\cdot\varpi^e$, 其中 $e$ 称为 $K$ 在 $\mathbb{Q}_p$ 上的分歧指数\index{fenqizhishu@分歧指数(ramification index)} (ramification index). 当 $e = 1$ (即 $\mathcal{O}_K$ 中极大理想是 $(p)$) 时, 称 $K$ 在 $\mathbb{Q}_p$ 上非分歧\index{feifenqi@非分歧(unramified)} (unramified).

\begin{lemma} \label{thm:4.3.3} 
对 $\mathbb{F}_p$ 的每个有限扩张 $k$, 存在 $\mathbb{Q}_p$ 的次数为 $[k:\mathbb{F}_p]$ 的非分歧扩张 $K$, 使 $\mathcal{O}_K/p\mathcal{O}_K = k$.
\end{lemma}

\begin{proof}
设 $a$ 为 $k$ 在 $\mathbb{F}_p$ 上的本原元, $f_0(X)$ 为 $a$ 在 $\mathbb{F}_p$ 上的最小多项式, 故
\[
k = \mathbb{F}_p[a]\cong\mathbb{F}_p[X]/(f_0(X))
\]
(\cite[5.1]{milne-ft}). 设 $f(X)$ 为 $f_0(X)$ 到 $\mathbb{Z}_p[X]$ 中首一多项式的任一提升. 因 $a$ 是 $f_0$ 的单根, 它提升为 $f(X)$ 的根 $\alpha$ (见I~\ref{thm:1.2.14}), 且 $\alpha\in\mathcal{O}_K$. 则 $K = \mathbb{Q}_p[\alpha]$ 具有所需性质.
\end{proof}

\begin{remark} \label{thm:4.3.4} 
与上面引理一样, 设 $K\supset\mathcal{O}_K\to k$. 设 $q$ 为 $k$ 的阶, 故 $k$ 中元素是 $X^q - X$ 的根. 则 Hensel 引理 (I\ref{thm:1.2.13}) 对 $\mathcal{O}_K$ 成立, 故 $X^q - X$ 在 $k$ 中的所有根都提升到 $\mathcal{O}_K$. 因此 $K$ 含 $X^q - X$ 的分裂域, 且等于它.
\end{remark}

设 $K$ 如引理 \ref{thm:4.3.3} 中. 因 $\mathcal{O}_K$ 是以 $p$ 为唯一素元 (在相差单位的意义下) 的主理想整环, $K^{\times}$ 中每个元素 $\alpha$ 可唯一写成 $up^m$, 其中 $u\in\mathcal{O}_K^{\times}$, $m\in\mathbb{Z}$. 定义 $\ord\nolimits_p(\alpha) = m$. 则 $\ord\nolimits_p$ 是延拓 $\ord\nolimits_p\colon\mathbb{Q}_p^{\times}\to\mathbb{Z}$ 的同态 $K^{\times}\to\mathbb{Z}$.

\hyperref[ch:2]{第二章}, \S\ref{sec:2.4} 对把 $\mathbb{Q}_p$ 换成非分歧扩张 $K$ 逐字成立, 只是此时
\[
E^0(K)/E^1(K)\cong\tilde{E}_{\mathrm{ns}}(k), \qquad
E^n(K)/E^{n+1}(K)\cong k.
\]
因此, 引理 \ref{thm:4.3.2} 在把 $\mathbb{Q}_p$ 换成 $K$、$\mathbb{F}_p$ 换成 $k$ 后仍有效.

\begin{lemma} \label{thm:4.3.5} 
设 $E$ 为 $\mathbb{Q}_p$ 上具有好约化的椭圆曲线, $n$ 为与 $p$ 互素的整数. 对每个 $P\in E(\mathbb{Q}_p)$, 存在 $\mathbb{Q}_p$ 的有限非分歧扩张 $K$ 使 $P\in nE(K)$.
\end{lemma}

\begin{proof}
设 $P\in E(\mathbb{Q}_p)$. 则 $\tilde{P}\in n\tilde{E}(k)$ 对 $\mathbb{F}_p$ 的某个有限扩张 $k$ (见~\ref{thm:4.2.1}). 设 $K$ 为剩余域 $k$ 的非分歧扩张 (见~\ref{thm:4.3.3}). 则用 $K$ 代替 $\mathbb{Q}_p$, 由引理 \ref{thm:4.3.2} 得 $P\in nE(K)$.
\end{proof}

\begin{proposition} \label{thm:4.3.6} 
设 $E$ 为 $\mathbb{Q}$ 上的椭圆曲线. 对每个 $\xi\in S^{(n)}(E/\mathbb{Q})$ 与不整除 $n$ 的好约化素数 $p$, 存在 $\mathbb{Q}_p$ 的有限非分歧扩张 $K$ 使 $\xi$ 在 $H^1(K, E_n)$ 中映射为零.
\end{proposition}

\begin{proof}
由 Selmer 群的定义, 存在映射到 $\xi$ 的像 $\xi_p\in H^1(\mathbb{Q}_p, E_n)$ 的 $P\in E(\mathbb{Q}_p)$. 因 $E$ 在 $p$ 处有好约化, 存在 $\mathbb{Q}_p$ 的非分歧扩张 $K$ 使 $P\in nE(K)$. 由此 $\xi_p$ 在 $H^1(K, E_n)$ 中映射为零,
\[
\begin{tikzcd}[
    row sep=large, 
    column sep=large,
    % 设置元素映射箭头的样式
    mapsto/.style={|->}
]
    % 第一行
    & & H^1(\mathbb{Q}, E_n) \dar & & & \gamma \dar[mapsto] \\
    % 第二行
    E(\mathbb{Q}_p) \rar["n"] \dar & E(\mathbb{Q}_p) \rar \dar & H^1(\mathbb{Q}_p, E_n) \dar & & P \rar[mapsto] \dar[mapsto] & \gamma_p \\
    % 第三行
    E(K) \rar["n"] & E(K) \rar & H^1(K, E_n) & Q \rar[mapsto] & P & 
\end{tikzcd}
\]
\end{proof}

\subsection*{特殊情形的有限性证明}

我们证明在 $E$ 的 2 阶点坐标属于 $\mathbb{Q}$ 的情形, $S^{(2)}(E/\mathbb{Q})$ 有限. 此条件意味着 $E$ 的方程形如
\[
Y^2Z = (X-\alpha Z)(X-\beta Z)(X-\gamma Z), \qquad \alpha, \beta, \gamma\in\mathbb{Q}
\]
(I\ref{thm:1.3.4}). 它蕴含
\[
E_2(\mathbb{Q}^{\mathrm{al}})\cong E_2(\mathbb{Q})\cong(\mathbb{Z}/2\mathbb{Z})^2\cong(\mu_2)^2,
\]
全都带 $\Gal(\mathbb{Q}^{\mathrm{al}}/\mathbb{Q})$ 的平凡作用, 故
\[
H^1(\mathbb{Q}, E_2)\cong H^1(\mathbb{Q}, \mu_2)^2\cong(\mathbb{Q}^{\times}/\mathbb{Q}^{\times2})^2
\]
(\ref{thm:4.1.8}b). 这些同构只依赖同构 $E_2(\mathbb{Q})\cong(\mathbb{Z}/2\mathbb{Z})^2$ 的选取, 我们固定之.

设 $\xi\in S^{(2)}(E/\mathbb{Q})\subset H^1(\mathbb{Q}, E_2)$. 对每个不整除 $2\Delta$ 的素数 $p_0$, 由~\ref{thm:4.3.6} 存在 $\mathbb{Q}_{p_0}$ 的有限非分歧扩张 $K$ 使 $\xi_{p_0}$ 在 $H^1(K, E_2)$ 中映射为零,
\[
\begin{tikzcd}[row sep=large, column sep=1.5cm]
    % 第一行：元素
    \xi & \xi_{p_0} & 0 \\
    % 第二行：第一层群
    H^1(\mathbb{Q}, E_2) \arrow[r] \arrow[d, "{\ref{thm:4.1.8}}"', "\simeq"] & 
    H^1(\mathbb{Q}_{p_0}, E_2) \arrow[r] \arrow[d, "{\ref{thm:4.1.8}}"', "\simeq"] & 
    H^1(K, E_2) \arrow[d, "{\ref{thm:4.1.8}}"', "\simeq"] \\
    % 第三行：第二层群
    (\mathbb{Q}^\times / \mathbb{Q}^{\times 2})^2 \arrow[r] & 
    (\mathbb{Q}_{p_0}^\times / \mathbb{Q}_{p_0}^{\times 2})^2 \arrow[r] & 
    (K^\times / K^{\times 2})^2.
\end{tikzcd}
\]
设 $\xi$ 在第一个同构下对应于
\[
\Bigl(((-1)^{\varepsilon(1)}\prod_p p^{\varepsilon(p)}), \ ((-1)^{\varepsilon'(1)}\prod_p p^{\varepsilon'(p)})\Bigr), \qquad 0\le\varepsilon(p), \varepsilon'(p)\le 1.
\]
则 $\xi_{p_0}$ 在第二个同构下对应于 $(p_0^{\varepsilon(p_0)}, p_0^{\varepsilon'(p_0)})$, 而它在 $(K^{\times}/K^{\times2})^2$ 中映射为零. 因 $p_0$ 是 $K$ 的整数环的素元, 这蕴含 $\varepsilon(p_0) = \varepsilon'(p_0) = 0$. 由于这对所有不整除 $2\Delta$ 的 $p_0$ 成立, $\xi$ 只有有限多种可能.

\begin{remark} \label{thm:4.3.7} 
在此情形可不提上同调群证明 $E(\mathbb{Q})/2E(\mathbb{Q})$ 有限. 考虑椭圆曲线
\[
Y^2Z = (X-\alpha Z)(X-\beta Z)(X-\gamma Z), \qquad \alpha, \beta, \gamma\in\mathbb{Z}.
\]
定义 $\varphi_\alpha\colon E(\mathbb{Q})/2E(\mathbb{Q})\to\mathbb{Q}^{\times}/\mathbb{Q}^{\times2}$ 为
\[
\varphi_\alpha((x:y:z)) = \begin{cases}
(x/z - \alpha)\mathbb{Q}^{\times2}, & z\neq 0, x\neq\alpha z,\\
(\alpha-\beta)(\alpha-\gamma)\mathbb{Q}^{\times2}, & z\neq 0, x = \alpha z,\\
\mathbb{Q}^{\times2}, & (x:y:z) = (0:1:0).
\end{cases}
\]
\footnote{校对注：原书在无穷远点的分支写为 $\mathbb{Q}^{\times}$，但值域是 $\mathbb{Q}^{\times}/\mathbb{Q}^{\times2}$，且群同态须把单位元送到单位平方类；此处及后文的同一公式均改为 $\mathbb{Q}^{\times2}$.} 类似定义 $\varphi_\beta$. 可直接证明 $\varphi_\alpha$ 与 $\varphi_\beta$ 是同态, 映射 $(\varphi_\alpha, \varphi_\beta)\colon E(\mathbb{Q})\to(\mathbb{Q}^{\times}/\mathbb{Q}^{\times2})^2$ 的核是 $2E(\mathbb{Q})$, 且 $\varphi_\alpha(P)$ 与 $\varphi_\beta(P)$ 由整除 $2\Delta$ 的素数的 $\pm$ 乘积表示 (见 \cite[Chap. 15]{cassels1991}).
\end{remark}

\subsection*{一般情形的有限性证明: 概述}

我们在\hyperref[ch:1]{第一章}看到 $\mathbb{Q}$ 对每个非零素理想 $(p)$ 有一个赋值, 从而有一个完备化 $\mathbb{Q}_p$, 还有一个方便记为 $\mathbb{Q}_\infty$ 的完备化 $\mathbb{R}$. 类似地, 数域 $L$ 对 $\mathcal{O}_L$ 的每个素理想有一个赋值 (从而完备化), 对 $L$ 到 $\mathbb{R}$ 的每个嵌入或 $L$ 到 $\mathbb{C}$ 的复共轭嵌入对有一个赋值. 记 $P(p)$ 为 $L$ 上延拓 $|\cdot|_p$ 的赋值集. 则 (\cite[8.2]{milne-ant}),
\[
L\otimes_{\mathbb{Q}}\mathbb{Q}_p\cong\prod_{v\in P(p)}L_v,
\]
其中 $L_v$ 是 $L$ 对 $v$ 的完备化. 令 $P = \bigcup_{p = 2, 3, \dots, \infty}P(p)$.

对数域 $L$ 上的椭圆曲线 $E$, 定义
\[
S^{(n)}(E/L) = \operatorname{Ker}\Bigl(H^1(L, E_n)\to\prod_{v\in P}H^1(L_v, E)\Bigr).
\]
下面的引理表明, 在证明 $S^{(n)}(E/L)$ 有限时, 可用更大的域替换 $L$.

\begin{lemma} \label{thm:4.3.8} 
设 $E$ 为数域 $L$ 上的椭圆曲线, $L'$ 为 $L$ 的有限 Galois 扩张. 对所有 $n\ge 1$, 映射 $S^{(n)}(E/L)\to S^{(n)}(E/L')$ 的核有限.
\end{lemma}

\begin{proof}
\footnote{校对注：原书两处写为 $\Gal(L'/\mathbb{Q})$. 本引理的 Galois 扩张为 $L'/L$，依注记 \ref{thm:4.1.7} 的膨胀--限制正合列，核应为 $H^1(\Gal(L'/L),E_n(L'))$.}
{\emergencystretch=3em 因 $S^{(n)}(E/L)$ 与 $S^{(n)}(E/L')$ 分别是 $H^1(L, E_n)$ 与 $H^1(L', E_n)$ 的子群, 只需证 $H^1(L, E_n)\to H^1(L', E_n)$ 的核有限. 但 (cf.~\ref{thm:4.1.7}) 此核是 $H^1(\Gal(L'/L), E_n(L'))$, 因 $\Gal(L'/L)$ 与 $E_n(L')$ 都有限, 故有限.\par}
\end{proof}

在特殊情形的有限性证明中, 我们用了以下事实:
\begin{enumerate}[label=\textup{(\Alph*)}]
\item $\mathbb{Q}$ 含 $1$ 的本原平方根;
\item $E(\mathbb{Q})_2 = E(\mathbb{Q}^{\mathrm{al}})_2$ (对 $E$ 的假设);
\item 对任何有限素数集 $T$, 映射
\[
r\mapsto(\ord\nolimits_p(r)\bmod 2)\colon\mathbb{Q}^{\times}/\mathbb{Q}^{\times2}\to\bigoplus_{p\notin T}\mathbb{Z}/2\mathbb{Z}
\]
的核有限.
\end{enumerate}
设 $E$ 为数域 $L$ 上的椭圆曲线. 在证明 $S^{(n)}(E/L)$ 有限时, 引理 \ref{thm:4.3.8} 允许我们假设 $L$ 含 $1$ 的 $n$ 次本原根, 且 $E(L)$ 含 $E(\mathbb{Q}^{\mathrm{al}})$ 中所有 $n$ 阶点. 在命题 \ref{thm:4.3.13} 中, 我们证明 (C) 对 $L$ 与 $n$ 的类似物由代数数论课程中证明的三个基本定理推出. 现在, 特殊情形的有限性证明可搬到 $E/L$ 的情形.

\subsection*{代数数论复习}

本节中, $L$ 是 $\mathbb{Q}$ 的有限扩张, $\mathcal{O}_L$ 是 $L$ 中全部代数整数的环 ($\mathbb{Z}$ 在 $L$ 中的整闭包); 见 III~\ref{thm:3.3.19}.

$\mathcal{O}_L$ 中每个元素是不可约 (即``不可分解'') 元素的乘积, 但此分解可能不唯一. 例如在 $\mathbb{Z}[\sqrt{-5}]$ 中,
\[
6 = 2\cdot 3 = (1+\sqrt{-5})(1-\sqrt{-5}),
\]
且 $2$, $3$, $1+\sqrt{-5}$, $1-\sqrt{-5}$ 不可约且没有两个相差一个单位. Kummer 与 Dedekind 为补救此问题而提出扩充``理想数'' (现在称为理想) 之集, 以恢复唯一分解. 对理想 $\mathfrak{a}$ 与 $\mathfrak{b}$, 有限和 $\sum a_ib_i$ ($a_i\in\mathfrak{a}$, $b_i\in\mathfrak{b}$) 的集合是理想, 记为 $\mathfrak{ab}$. 环 $\mathcal{O}_L$ 是唯一分解整环当且仅当它是主理想整环.

\begin{theorem}[Dedekind] \label{thm:4.3.9} 
$\mathcal{O}_L$ 中每个理想都可唯一写成素理想的乘积.
\end{theorem}

\begin{proof}
见 \cite[3.7, 3.29]{milne-ant} 或 \cite[\S 3.4]{samuel2008}.
\end{proof}

例如, 在 $\mathbb{Z}[\sqrt{-5}]$ 中,
\[
(6) = (2, 1+\sqrt{-5})^2(3, 1+\sqrt{-5})(3, 1-\sqrt{-5}).
\]
对元素 $a\in\mathcal{O}_L$ 与 $\mathcal{O}_L$ 中素理想 $\mathfrak{p}$, 记 $\ord\nolimits_\mathfrak{p}(a)$ 为 $\mathfrak{p}$ 在理想 $(a)$ 的唯一分解中的指数, 故
\[
(a) = \prod_{\mathfrak{p}}\mathfrak{p}^{\ord\nolimits_\mathfrak{p}(a)}.
\]
对 $x = \frac{a}{b}\in L$, 定义 $\ord\nolimits_\mathfrak{p}(x) = \ord\nolimits_\mathfrak{p}(a) - \ord\nolimits_\mathfrak{p}(b)$. $\mathcal{O}_L$ (或 $L$) 的理想类群 (ideal class group)\index{lixiangleiqun@理想类群(ideal class group)} $C$ 定义为零化的同态
\[
L^{\times}\xrightarrow{x\mapsto(\ord\nolimits_\mathfrak{p}(x))}\bigoplus_{\mathfrak{p}\subset\mathcal{O}_L, \ \mathfrak{p}\text{ 素}}\mathbb{Z}\to C\to 0
\]
的余核. 定理 \ref{thm:4.3.9} 允许我们把 $\mathcal{O}_L$ 中理想幺半群等同于 $\bigoplus_{\mathfrak{p}}\mathbb{N}\subset\bigoplus_{\mathfrak{p}}\mathbb{Z}$. $C$ 的每个元素由 $\mathcal{O}_L$ 中的非零理想表示, 且两个非零理想 $\mathfrak{a}, \mathfrak{b}\in\mathcal{O}_L$ 表示 $C$ 的同一元素当且仅当 $(a)\mathfrak{a} = (b)\mathfrak{b}$ 对某些 $a, b\in\mathcal{O}_L$. 当且仅当 $\mathcal{O}_L$ 是主理想整环时 $C = 0$, 故 $C$ 度量了 $\mathcal{O}_L$ 中分解不唯一性的程度.

\begin{theorem}[类数有限性] \label{thm:4.3.10} 
理想类群 $C$ 有限.
\end{theorem}

\begin{proof}
理想 $\mathfrak{a}$ 的数值范数 (numerical norm) $\Nm\mathfrak{a}$ 定义为指标 $(\mathcal{O}_L:\mathfrak{a})$. 若 $\mathfrak{a} = \prod\mathfrak{p}_i^{r_i}$, 则 $\Nm\mathfrak{a} = \prod(\Nm\mathfrak{p})^{r_i}$. 由此, 以及 $\mathbb{Z}$ 的每个素理想 $(p)$ 之上只有有限多个 $\mathcal{O}_L$ 的素理想 $\mathfrak{p}$, 我们看到对任何实数 $B$, 只有有限多个满足 $\Nm\mathfrak{a}\le B$ 的理想 $\mathfrak{a}$. 为证 $C$ 有限, (更精确地) 证明 $C$ 的每个元素由满足
\[
\Nm\mathfrak{a}\le\frac{n!}{n^n}\Bigl(\frac{4}{\pi}\Bigr)^s|\Delta_L|^{1/2}
\]
的理想表示, 其中 $n = [L:\mathbb{Q}]$, $2s$ 是 $L$ 的非实复嵌入数, $\Delta_L$ 是 $\mathcal{O}_L$ 在 $\mathbb{Z}$ 上的判别式. 见 \cite[4.3]{milne-ant} 或 \cite[\S 4.3]{samuel2008}.
\end{proof}

我们还需理解 $\mathcal{O}_L$ 的单位群 $\mathcal{O}_L^{\times}$. 群 $\mathbb{Z}^{\times} = \{\pm 1\}$, 但已经对 $\mathbb{Q}[\sqrt{2}]$ 中整数环 $\mathbb{Z}[\sqrt{2}]$, 该群无限, 因为 $\sqrt{2}+1$ 是 $\mathbb{Z}[\sqrt{2}]$ 中的单位:
\[
(\sqrt{2}+1)(\sqrt{2}-1) = 1.
\]
事实上
\[
\mathbb{Z}[\sqrt{2}]^{\times} = \{\pm(1+\sqrt{2})^n \mid n\in\mathbb{Z}\}\cong\mathbb{Z}/2\mathbb{Z}\oplus\mathbb{Z}.
\]

\begin{theorem}[Dedekind 单位定理] \label{thm:4.3.11} 
$\mathcal{O}_L$ 的单位群 $\mathcal{O}_L^{\times}$ 有限生成.
\end{theorem}

\begin{proof}
$L$ 中单位根群 $\mu(L)$ 有限 (因为对所有 $n\ge 1$, $n$ 次分圆多项式在 $\mathbb{Q}$ 上不可约; \cite[5.9]{milne-ft}) 且循环 (因为对所有 $n\ge 1$, 它含至多 $n$ 个阶整除 $n$ 的元素).

选取同构
\[
L\otimes_{\mathbb{Q}}\mathbb{R}\to\mathbb{R}^r\times\mathbb{C}^s, \qquad
\alpha\otimes 1\mapsto(\sigma_1\alpha, \dots, \sigma_r\alpha, \sigma_{r+1}\alpha, \dots, \sigma_{r+s}\alpha).
\]
这给出同态
\[
\mathcal{O}_L^{\times}\to\mathbb{R}^{r+s}, \qquad
\alpha\mapsto(\log|\sigma_1\alpha|, \dots, \log|\sigma_r\alpha|, \log|\sigma_{r+1}\alpha|, \dots, \log|\sigma_{r+s}\alpha|).
\]
定理的证明是证明该同态的核是 $L^{\times}$ 的有限子群 $\mu(L^{\times})$, 像是超平面
\[
H\colon x_1 + \cdots + x_r + 2x_{r+1} + \cdots + 2x_{r+s}
\]
中 $\mathbb{R}^{r+s}$ 的满格, 故
\[
\mathcal{O}_L^{\times}\cong\mu(L)\oplus\mathbb{Z}^{r+s-1}.
\]
见 \cite[Chapter 5]{milne-ant} 或 \cite[\S 4.4]{samuel2008}.
\end{proof}

$\mathcal{O}_L$ 的元素 $a$ 是单位当且仅当 $(a) = \mathcal{O}_L$. 这等价于对所有素理想 $\mathfrak{p}$ 有 $\ord\nolimits_\mathfrak{p}(a) = 0$, 故有正合序列
\[
0\to U\to L^{\times}\to\bigoplus_{\mathfrak{p}}\mathbb{Z}\to C\to 0
\]
其中 $U$ 有限生成, $C$ 有限. 更一般地, 有下面的陈述.

\begin{corollary} \label{thm:4.3.12} 
设 $T$ 为 $\mathcal{O}_L$ 的有限素理想集, 由正合序列
\begin{align*}
    0\longrightarrow U_T\longrightarrow L^{\times}& \longrightarrow\bigoplus_{\mathfrak{p}\notin T}\mathbb{Z}\longrightarrow C_T\longrightarrow 0\\
    x & \longmapsto (\ord\nolimits_\mathfrak{p}(x)).
\end{align*}
定义 $U_T$ 与 $C_T$. 则 $U_T$ 有限生成, $C_T$ 有限.
\end{corollary}

\begin{proof}
映射对 $L^{\times}\longrightarrow\bigoplus_{\text{所有 }\mathfrak{p}}\mathbb{Z}\longrightarrow\bigoplus_{\mathfrak{p}\notin T}\mathbb{Z}$ (投影) 的核--余核序列 (\ref{thm:4.2.2}) 是正合序列
\[
0\longrightarrow U\longrightarrow U_T\longrightarrow\bigoplus_{\mathfrak{p}\in T}\mathbb{Z}\longrightarrow C\longrightarrow C_T\longrightarrow 0.
\]
\end{proof}

\subsection*{一般情形的有限性证明: 完成}

先推广陈述 (C).

\begin{proposition} \label{thm:4.3.13} 
对代数数域 $L$ 与 $n\ge 1$, 设 $N$ 为映射 $a\mapsto(\ord\nolimits_\mathfrak{p}(a)\bmod n)\colon L^{\times}/L^{\times n}\to\bigoplus_{\mathfrak{p}\notin T}\mathbb{Z}/n\mathbb{Z}$ 的核. 则有正合序列
\[
0\to U_T/U_T^n\to N\to(C_T)_n\to 0.
\]
\end{proposition}

\begin{proof}
令 $P = \operatorname{Ker}\bigl(\bigoplus_{\mathfrak{p}\notin T}\mathbb{Z}\to C_T\bigr)$. 所需序列是映射对
\[
L^{\times}/L^{\times n}\xrightarrow{a}P/nP\xrightarrow{b}\bigoplus_{\mathfrak{p}\notin T}\mathbb{Z}/n\mathbb{Z}
\]
的核--余核序列. 为计算 $a$ 与 $b$ 的核与余核, 把蛇引理应用于图
\[\begin{tikzcd}
	& {U_T} & {L^\times} & P & 0 & P & {\bigoplus_{\mathfrak{p}\notin T}\Z} & {C_T} & 0 \\
	0 & {U_T} & {L^\times} & P & 0 & P & {\bigoplus_{\mathfrak{p}\notin T}\Z} & {C_T}.
	\arrow[from=1-2, to=1-3]
	\arrow["n", from=1-2, to=2-2]
	\arrow[from=1-3, to=1-4]
	\arrow["n", from=1-3, to=2-3]
	\arrow[from=1-4, to=1-5]
	\arrow["n", from=1-4, to=2-4]
	\arrow[from=1-6, to=1-7]
	\arrow["n", from=1-6, to=2-6]
	\arrow[from=1-7, to=1-8]
	\arrow["n", from=1-7, to=2-7]
	\arrow[from=1-8, to=1-9]
	\arrow["n", from=1-8, to=2-8]
	\arrow[from=2-1, to=2-2]
	\arrow[from=2-2, to=2-3]
	\arrow[from=2-3, to=2-4]
	\arrow[from=2-5, to=2-6]
	\arrow[from=2-6, to=2-7]
	\arrow[from=2-7, to=2-8]
\end{tikzcd}\]
\end{proof}

\begin{corollary} \label{thm:4.3.14} 
群 $N$ 有限.
\end{corollary}

\begin{proof}
这由~\ref{thm:4.3.12} 与 \ref{thm:4.3.13} 推出.
\end{proof}

现在与特殊情形相同的论证证明 $S^{(n)}(E/L)$ 有限. 由引理 \ref{thm:4.3.8}, 可设 $L$ 含 $1$ 的 $n$ 次本原根且 $E(L)_n = E(L^{\mathrm{al}})_n$, 故
\begin{equation} \label{eq:30}
H^1(L, E_n)\cong L^{\times}/L^{\times n}\times L^{\times}/L^{\times n}.
\tag{30}
\end{equation}
设 $T$ 为 $\mathcal{O}_L$ 中与 $n$ 不互素或 $E$ 有坏约化的 (有限) 素理想集. 设 $\xi\in S^{(n)}(E/L)$. 对每个 $\mathfrak{p}\notin T$, 有交换图
\[\begin{tikzcd}[row sep=large, column sep=large]
    H^1(L, E_n) \arrow[r] \arrow[d, "\simeq"'] & 
    H^1(L_{\mathfrak{p}}, E_n) \arrow[r] \arrow[d, "\simeq"'] & 
    H^1(L'_{\mathfrak{p}}, E_n) \arrow[d, "\simeq"'] \\
    (L^\times/L^{\times n})^2 \arrow[r] & 
    (L_{\mathfrak{p}}^\times/L_{\mathfrak{p}}^{\times n})^2 \arrow[r] \arrow[d, "(\text{ord}_{\mathfrak{p}})^2"] & 
    (L_{\mathfrak{p}}'^\times/L_{\mathfrak{p}}'^{\times n})^2 \arrow[d, "(\text{ord}_{\mathfrak{p}})^2"] \\
    & (\mathbb{Z}/n\mathbb{Z})^2 \arrow[r, equal] & 
    (\mathbb{Z}/n\mathbb{Z})^2
\end{tikzcd}\]
其中 $L_\mathfrak{p}'$ 是 $L_\mathfrak{p}$ 的非分歧扩张, 使 $\xi$ 在 $H^1(L_\mathfrak{p}', E_n)$ 中映射为零 (cf.~\ref{thm:4.3.6}). 由此, 在同构~\eqref{eq:30} 下, $\xi$ 映射入有限群 $N\times N$.

\begin{aside} \label{thm:4.3.15} 
上述 Selmer 群有限性的证明比标准证明简单, 标准证明不必要地``把 $E(L)/nE(L)$ 的推测有限性转译成关于 $L$ 的某些域扩张的陈述''. 然而, 其他论证可能更明显地算法化.
\end{aside}

\section{高度；有限基定理的证明}\label{sec:4.4}

设 $P = (a_0:\dots:a_n)\in\mathbb{P}^n(\mathbb{Q})$. 若 $a_0, \dots, a_n\in\mathbb{Z}$ 且 $\gcd(a_0, \dots, a_n) = 1$, 则称 $(a_0, \dots, a_n)$ 为 $P$ 的本原代表元 (primitive representative).\index{benyuandaibiaoyuan@本原代表元(primitive representative)} $P$ 的高度 (height) $H(P)$ 定义为
\[
H(P) = \max_i|a_i|,
\]
其中 $|\cdot|$ 是通常的绝对值. $P$ 的对数高度 (logarithmic height)\index{gaodu@高度(height)!duishu@对数(logarithmic)} $h(P)$ 定义为 $\log H(P)$.

\subsection*{$\mathbb{P}^1$ 上的高度}

设 $F(X, Y)$ 与 $G(X, Y)$ 为 $\mathbb{Q}[X, Y]$ 中在 $\mathbb{P}^1(\mathbb{Q}^{\mathrm{al}})$ 上无公共零点的 $m$ 次齐次多项式. 则 $F$ 与 $G$ 定义映射
\[
\varphi\colon\mathbb{P}^1(\mathbb{Q})\to\mathbb{P}^1(\mathbb{Q}), \qquad
(x:y)\mapsto(F(x,y):G(x,y)).
\]

\begin{lemma} \label{thm:4.4.1} 
存在常数 $B$ 使
\[
|h(\varphi(P)) - mh(P)|\le B, \qquad \forall P\in\mathbb{P}^1(\mathbb{Q}).
\]
\end{lemma}

\begin{proof}
因用非零常数乘 $F$ 与 $G$ 不改变 $\varphi$, 可设 $F$ 与 $G$ 有整数系数. 设 $P\in\mathbb{P}^1(\mathbb{Q})$, $(a:b)$ 为其本原代表元. 对任何单项式 $cX^iY^{m-i}$,
\[
|ca^ib^j|\le|c|\max(|a|^m, |b|^m),
\]
故
\[
|F(a,b)|, |G(a,b)|\le C\cdot\max(|a|, |b|)^m
\]
其中 $C = (m+1)\max(|F \text{ 或 } G \text{ 的系数}|)$. 因 $F(a,b), G(a,b)\in\mathbb{Z}$,
\begin{align}
    H(\varphi(P))& \le\max(|F(a,b)|, |G(a,b)|) \tag{31}\label{eq:31}\\
    & \le C\cdot\max(|a|, |b|)^m = C\cdot H(P)^m. \tag{32}\label{eq:32}
\end{align}
取对数, 得不等式 $h(\varphi(P))\le mh(P) + \log C$.

证明反向不等式的困难在于 $F(a,b)$ 与 $G(a,b)$ 可能有大的公共因子, 从而 \eqref{eq:31} 可能严格. 我们用 $F$ 与 $G$ 在 $\mathbb{Q}^{\mathrm{al}}$ 上无公共零点的假设来限制此问题.

该假设说 $F$ 与 $G$ (作为齐次多项式) 的结式 $R$ 非零 (I\ref{thm:1.1.26}). 考虑 $Y^{-m}F(X, Y) = F\bigl(\frac{X}{Y}, 1\bigr)$ 与 $Y^{-m}G(X, Y) = G\bigl(\frac{X}{Y}, 1\bigr)$. 视为单变量 $\frac{X}{Y}$ 的多项式, $F\bigl(\frac{X}{Y}, 1\bigr)$ 与 $G\bigl(\frac{X}{Y}, 1\bigr)$ 有与 $F(X, Y)$ 与 $G(X, Y)$ 相同的结式, 故 (见 I, \S\ref{sec:1.1}) 存在次数 $m-1$ 的多项式 $U\bigl(\frac{X}{Y}\bigr)$, $V\bigl(\frac{X}{Y}\bigr)\in\mathbb{Z}\bigl[\frac{X}{Y}\bigr]$ 使
\[
U\Bigl(\frac{X}{Y}\Bigr)F\Bigl(\frac{X}{Y}, 1\Bigr) + V\Bigl(\frac{X}{Y}\Bigr)G\Bigl(\frac{X}{Y}, 1\Bigr) = R.
\]
两边乘以 $Y^{2m-1}$ 并把 $Y^{m-1}U\bigl(\frac{X}{Y}\bigr)$ 改名为 $U(X, Y)$, 把 $Y^{m-1}V\bigl(\frac{X}{Y}\bigr)$ 改名为 $V(X, Y)$, 得方程
\[
U(X, Y)F(X, Y) + V(X, Y)G(X, Y) = RY^{2m-1}.
\]
类似地, 存在次数 $m-1$ 的齐次多项式 $U'(X, Y)$ 与 $V'(X, Y)$ 使
\[
U'(X, Y)F(X, Y) + V'(X, Y)G(X, Y) = RX^{2m-1}.
\]
把 $(a, b)$ 代入 $(X, Y)$, 得方程
\begin{equation} \label{eq:33}
U(a, b)F(a, b) + V(a, b)G(a, b) = Rb^{2m-1}, \qquad
U'(a, b)F(a, b) + V'(a, b)G(a, b) = Ra^{2m-1}.
\tag{33}
\end{equation}
由这些方程, 推出 $\gcd(F(a,b), G(a,b))$ 整除 $\gcd(Ra^{2m-1}, Rb^{2m-1}) = R$. 如证明第一部分, 存在 $C > 0$ 使
\[
U(a,b), U'(a,b), V(a,b), V'(a,b)\le C\cdot\max(|a|, |b|)^{m-1}.
\]
因此, 方程~\eqref{eq:33} 表明
\[
2C\cdot\max(|a|, |b|)^{m-1}\cdot\max(|F(a,b)|, |G(a,b)|)\ge |R||a|^{2m-1}, \quad |R||b|^{2m-1},
\]
即
\[
\max(|F(a,b)|, |G(a,b)|)\ge\frac{|R|\cdot H(P)^{2m-1}}{2C\cdot H(P)^{m-1}} = \frac{|R|}{2C}H(P)^m.
\]
\footnote{校对注：原书上式分母误用 $h(P)$，末项又写为 $H(P)^{m-1}$. 依 $H(P)=\max(|a|,|b|)$ 及前一式，分母应为 $H(P)^{m-1}$，而指数 $(2m-1)-(m-1)=m$；此处统一更正.} 因 $\gcd(F(a,b), G(a,b))\mid R$,
\[
H(\varphi(P))\ge\frac{1}{|R|}\max(|F(a,b)|, |G(a,b)|)\ge\frac{1}{2C}H(P)^m.
\]
取对数, 得不等式 $h(\varphi(P))\ge mh(P) - \log 2C$.
\end{proof}

\begin{lemma} \label{thm:4.4.2} 
设 $R$ 为 $(P, Q)$ 在映射
\[
(a:b), (c:d)\mapsto(ac:ad+bc:bd)\colon\mathbb{P}^1(\mathbb{Q})\times\mathbb{P}^1(\mathbb{Q})\to\mathbb{P}^2(\mathbb{Q})
\]
下的像. 则
\[
\frac{1}{2}\le\frac{H(R)}{H(P)H(Q)}\le 2.
\]
\end{lemma}

\begin{proof}
选 $P$ 与 $Q$ 的本原代表元 $(a:b)$ 与 $(c:d)$. 则
\[
H(R)\le\max(|ac|, |ad+bc|, |bd|)\le 2\max(|a|, |b|)\max(|c|, |d|) = 2H(P)H(Q).
\]
若素数 $p$ 同时整除 $ac$ 与 $bd$, 则要么它整除 $a$ 与 $d$ 但不整除 $b$ 与 $c$, 要么反之. 任一种情形它都不整除 $ad + bc$, 故 $(ac, ad+bc, bd)$ 是 $R$ 的本原代表元. 剩下证
\[
\max(|ac|, |ad+bc|, |bd|)\ge\frac{1}{2}(\max(|a|, |b|))(\max(|c|, |d|)),
\]
但这是初等练习 (例如把 $a, b, c, d$ 视为实数, 缩放两对使 $a = 1 = c$).
\end{proof}

\subsection*{$E$ 上的高度}

设 $E$ 为椭圆曲线
\[
E\colon Y^2Z = X^3 + aXZ^2 + bZ^3, \qquad a, b\in\mathbb{Q}, \ \Delta = 4a^3 + 27b^2\neq 0.
\]
对 $P\in E(\mathbb{Q})$, 定义
\[
H(P) = \begin{cases} H((x(P):z(P))), & \text{若 } z(P)\neq 0,\\ 1, & \text{若 } P = (0:1:0), \end{cases}
\qquad
h(P) = \log H(P).
\]
还有 $h$ 的其他定义, 但它们相差有界量, 因此给出相同的典范高度 (见下).

\begin{lemma} \label{thm:4.4.3} 
对任何常数 $B$, 集合 $\{P\in E(\mathbb{Q}) \mid h(P) < B\}$ 有限.
\end{lemma}

\begin{proof}
当然, 对任何 $B$, 集合 $\{P\in\mathbb{P}^1(\mathbb{Q}) \mid H(P)\le B\}$ 有限, 但对每个点 $(x_0:z_0)\in\mathbb{P}^1(\mathbb{Q})$, 至多有两个点 $(x_0:y:z_0)\in E(\mathbb{Q})$, 故 $\{P\in E(\mathbb{Q}) \mid H(P)\le B\}$ 有限.
\end{proof}

\begin{lemma} \label{thm:4.4.4} 
存在常数 $A$ 使
\[
|h(2P) - 4h(P)|\le A.
\]
\end{lemma}

\begin{proof}
设 $P = (x:y:z)$, $2P = (x_2:y_2:z_2)$. 由\hyperref[eq:dup]{倍点公式}，
\[
(x_2:z_2) = (F(x):G(x)),
\]
其中 $F(X, Z)$ 与 $G(X, Z)$ 是满足
\[
\frac{F(X, 1)}{G(X, 1)} = \frac{(3X^2+a)^2 - 8X(X^3+aX+b)}{4(X^3+aX+b)}
\]
的 $4$ 次齐次多项式. 因 $X^3 + aX + b$ 与其导数 $3X^2 + a$ 无公共根, $F(X, 1)$ 与 $G(X, 1)$ 也无公共根, 故引理 \ref{thm:4.4.1} 表明对某个常数 $A$,
\[
|h(2P) - 4h(P)|\le A.
\]
\end{proof}

\begin{lemma} \label{thm:4.4.5} 
至多存在一个满足下列条件的函数 $\hat{h}\colon E(\mathbb{Q})\to\mathbb{R}$:
\begin{enumerate}[label=\textup{(\alph*)}]
\item $\hat{h}(P) - h(P)$ 在 $E(\mathbb{Q})$ 上有界;
\item $\hat{h}(2P) = 4\hat{h}(P)$.
\end{enumerate}
\end{lemma}

\begin{proof}
若 $\hat{h}$ 满足 (a) 且界为 $B$, 则
\[
\Bigl|\hat{h}(2^nP) - h(2^nP)\Bigr|\le B.
\]
若它还满足 (b), 则
\[
\Bigl|\hat{h}(P) - \frac{h(2^nP)}{4^n}\Bigr|\le\frac{B}{4^n},
\]
故
\[
\hat{h}(P) = \lim_{n\to\infty}h(2^nP)/4^n.
\]
\end{proof}

满足引理中条件 (a) 与 (b) 的任何函数 $\hat{h}\colon E(\mathbb{Q})\to\mathbb{R}$ 都称为典范 (canonical) 或 N\'eron--Tate 高度函数.\index{gaodu@高度(height)!Neron-Tate@N\'eron--Tate(Néron-Tate)}\footnote{文献中典范高度的归一化不同, 其中一个是另一个的两倍 (见 \cite[3.4]{cremona1997}).} 由证明, 若 $\hat{h}$ 存在, 则 $\hat{h}(P)$ 必是序列 $h(2^nP)/4^n$ 的极限.

\begin{lemma} \label{thm:4.4.6} 
对所有 $P\in E(\mathbb{Q})$, 序列 $h(2^nP)/4^n$ 在 $\mathbb{R}$ 中是 Cauchy 列.
\end{lemma}

\begin{proof}
由引理 \ref{thm:4.4.4}, 存在常数 $A$ 使对所有 $P$,
\[
|h(2P) - 4h(P)|\le A.
\]
对 $N\ge M\ge 0$ 与 $P\in E(\mathbb{Q})$,
\begin{align*}
    \Bigl|\frac{h(2^NP)}{4^N} - \frac{h(2^MP)}{4^M}\Bigr|
    & = \Bigl|\sum_{n=M}^{N-1}\Bigl(\frac{h(2^{n+1}P)}{4^{n+1}} - \frac{h(2^nP)}{4^n}\Bigr)\Bigr|
    \\ & \le\sum_{n=M}^{N-1}\frac{1}{4^{n+1}}|h(2^{n+1}P) - 4h(2^nP)|
    \\ & \le\sum_{n=M}^{N-1}\frac{1}{4^{n+1}}A 
    \\ & \le\frac{A}{4^{M+1}}\Bigl(1 + \frac{1}{4} + \frac{1}{4^2} + \cdots\Bigr) 
    \\ & = \frac{A}{3\cdot 4^M}.
\end{align*}
故序列 $h(2^nP)/4^n$ 是 Cauchy 列.
\end{proof}

引理允许我们定义
\[
\hat{h}(P) = \lim_{n\to\infty}\frac{h(2^nP)}{4^n}, \qquad \forall P\in E(\mathbb{Q}).
\]

\begin{theorem} \label{thm:4.4.7} 
函数 $\hat{h}\colon E(\mathbb{Q})\to\mathbb{R}$ 是 N\'eron--Tate 高度函数. 且
\begin{enumerate}[label=\textup{(\alph*)}]
\item 对所有 $C\ge 0$, 集合 $\{P\in E(\mathbb{Q}) \mid \hat{h}(P)\le C\}$ 有限;
\item $\hat{h}(P)\ge 0$, 等号当且仅当 $P$ 是有限阶点.
\end{enumerate}
\end{theorem}

\begin{proof}
在引理 \ref{thm:4.4.6} 的证明中取 $M = 0$, 不等式变为
\[
\Bigl|\frac{h(2^NP)}{4^N} - h(P)\Bigr|\le\frac{A}{3}.
\]
令 $N\to\infty$, 得 $\hat{h}$ 满足引理 \ref{thm:4.4.5} 的条件 (a). 对条件 (b), 注意
\[
\hat{h}(2P) = \lim_{n\to\infty}\frac{h(2^{n+1}P)}{4^n} = 4\cdot\lim_{n\to\infty}\frac{h(2^{n+1}P)}{4^{n+1}} = 4\cdot\hat{h}(P).
\]
故 $\hat{h}$ 是 N\'eron--Tate 高度函数.

集合 $\{P \mid \hat{h}(P)\le C\}$ 有限, 因为 $h$ 有此性质且 $\hat{h}(P) - h(P)$ 有界.

因 $H(P)$ 是 $\ge 1$ 的整数, $h(P)$ 与 $\hat{h}(P)$ 都 $\ge 0$. 若 $P$ 是挠点, 则 $\{2^nP \mid n\ge 0\}$ 有限, 故 $\hat{h}$ 在其上有界, 设界为 $D$, 则对所有 $n$, $\hat{h}(P) = \hat{h}(2^nP)/4^n\le D/4^n$; 故 $\hat{h}(P) = 0$. 另一方面, 若 $P$ 无限阶, 则 $\{2^nP \mid n\ge 0\}$ 无限, 且 (由 (a)) $\hat{h}$ 在其上无界. 故对某个 $n$, $\hat{h}(2^nP) > 1$, 从而 $\hat{h}(P) > 4^{-n} > 0$.
\end{proof}

设 $f\colon M\to K$ 为从 Abel 群 $M$ 到特征 $\neq 2$ 的域 $K$ 的函数. 若 $f(2x) = 4f(x)$ 且
\[
B(x, y) \overset{\text{def}}{=} f(x+y) - f(x) - f(y)
\]
双加性, 则称 $f$ 为二次型 (quadratic form). 则 $B$ 对称, 且它是满足 $f(x) = \frac{1}{2}B(x, x)$ 的唯一对称双加性形式 $B\colon M\times M\to K$. 我们需要下面的判据.

\begin{lemma} \label{thm:4.4.8} 
从 Abel 群到特征 $\neq 2$ 的域 $K$ 的函数 $f\colon M\to K$ 若满足平行四边形律 (parallelogram law)\footnote{在初等线性代数中, 平行四边形定律说的是, 对于$\R^n$中的向量$u,v$, 成立\[\|u+v\|^2 + \|u-v\|^2 = 2\|u\|^2+2\|v\|^2.\]}
\[
f(x+y) + f(x-y) = 2f(x) + 2f(y), \qquad \forall x, y\in M,
\]
则是二次型.
\end{lemma}

\begin{proof}
设 $f$ 满足平行四边形律. 在其中取 $x = y = 0$, 得 $f(0) = 0$; 取 $x = y$, 得 $f(2x) = 4f(x)$; 取 $x = 0$, 得 $f(-y) = f(y)$. 因 $B(x, y)$ 对称, 剩下证 $B(x+y, z) = B(x, z) + B(y, z)$, 即
\[
f(x+y+z) - f(x+y) - f(x+z) - f(y+z) + f(x) + f(y) + f(z) = 0.
\]
现在四次应用平行四边形律表明
\begin{align*}
    f(x+y+z) + f(x+y-z) - 2f(x+y) - 2f(z) = 0,\\
    f(x+y-z) + f(x-y+z) - 2f(x) - 2f(y-z) = 0,\\
    f(x+y+z) + f(x-y+z) - 2f(x+z) - 2f(y) = 0,\\
    2f(y+z) + 2f(y-z) - 4f(y) - 4f(z) = 0.
\end{align*}
这些方程的交替和是 (两倍的) 所需方程.
\end{proof}

\begin{proposition} \label{thm:4.4.9} 
高度函数 $\hat{h}\colon E(\mathbb{Q})\to\mathbb{R}$ 是二次型.
\end{proposition}

我们必须证明平行四边形律.

\begin{lemma} \label{thm:4.4.10} 
存在常数 $C$ 使
\[
H(P_1+P_2)\cdot H(P_1-P_2)\le C\cdot H(P_1)^2\cdot H(P_2)^2
\]
对所有 $P_1, P_2\in E(\mathbb{Q})$ 成立.
\end{lemma}

\begin{proof}
设 $P_1 + P_2 = P_3$, $P_1 - P_2 = P_4$, $P_i = (x_i:y_i:z_i)$. 则
\[
(x_3x_4:x_3z_4 + x_4z_3:z_3z_4) = (w_0:w_1:w_2),
\]
其中
\begin{align*}
    w_0 & = x_1^2x_2^2 - 2ax_1x_2z_1z_2 - 4b(x_1z_1z_2^2 + x_2z_1^2z_2) + a^2z_1^2z_2^2,\\
    w_1 & = 2(x_1x_2 + az_1z_2)(x_1z_2 + x_2z_1) + 4bz_1^2z_2^2,\\
    w_2 & = (x_2z_1 - x_1z_2)^2.
\end{align*}
例如, 由 II~\S\ref{sec:2.2} 的加法公式,
\begin{align*}
    x_3(x_1 - x_2)^2 = x_1x_2^2 + x_1^2x_2 - 2y_1y_2 + a(x_1 + x_2) + 2b,\\
    x_4(x_1 - x_2)^2 = x_1x_2^2 + x_1^2x_2 + 2y_1y_2 + a(x_1 + x_2) + 2b.
\end{align*}
相加这些方程并齐次化, 得方程 $x_3z_4 + x_4z_3 = w_1/w_2$. 类似地, $x_3x_4 = w_0/w_2$.

选取 $P_1$ 与 $P_2$ 的本原代表元 $(x_1:y_1:z_1)$ 与 $(x_2:y_2:z_2)$ 后, 显然对某个常数 $C$ 有
\[
H(w_0:w_1:w_2)\le C\cdot H(P_1)^2\cdot H(P_2)^2.
\]
这里 $H(w_0:w_1:w_2)$ 是 $(w_0:w_1:w_2)$ 作为 $\mathbb{P}^2(\mathbb{Q})$ 中点的极高. 由引理 \ref{thm:4.4.2},
\[
H(w_0:w_1:w_2)\ge\frac{1}{2}H(P_3)H(P_4).
\]
合并这两个不等式, 得所需不等式 (常数 $2C$).
\end{proof}

\begin{lemma} \label{thm:4.4.11} 
典范高度函数 $\hat{h}\colon E(\mathbb{Q})\to\mathbb{R}$ 满足平行四边形律:
\[
\hat{h}(P+Q) + \hat{h}(P-Q) = 2\hat{h}(P) + 2\hat{h}(Q).
\]
\end{lemma}

\begin{proof}
对前一引理取对数, 得
\[
h(P+Q) + h(P-Q)\le 2h(P) + 2h(Q) + \log C.
\]
把 $P$ 与 $Q$ 换成 $2^nP$ 与 $2^nQ$, 除以 $4^n$, 令 $n\to\infty$, 得不等式
\[
\hat{h}(P+Q) + \hat{h}(P-Q)\le 2\hat{h}(P) + 2\hat{h}(Q).
\]
在其中代入 $P' = P+Q$, $Q' = P-Q$, 得反向不等式:
\[
\hat{h}(P') + \hat{h}(Q')\le 2\hat{h}\Bigl(\frac{P'+Q'}{2}\Bigr) + 2\hat{h}\Bigl(\frac{P'-Q'}{2}\Bigr) = \frac{1}{2}\hat{h}(P'+Q') + \frac{1}{2}\hat{h}(P'-Q').
\]
\end{proof}

\begin{remark} \label{thm:4.4.12} 
设 $K$ 为数域. 当 $\mathcal{O}_K$ 不是主理想整环时, $\mathbb{P}^n(K)$ 中的点 $P$ 可能没有本原代表元,\footnote{更准确地说, 可能不存在 $a_0, \dots, a_n\in\mathcal{O}_K$ 使 $P = (a_0:\dots:a_n)$ 且由 $a_i$ 生成的理想为 $\mathcal{O}_K$.} 故我们对 $\mathbb{P}^n(\mathbb{Q})$ 中极高所给的定义不能直接推广到数域. 相反, 我们需要略微不同的方法. 注意对 $c\in\mathbb{Q}^{\times}$,
\[
\prod_{p = 2, \dots, \infty}|c|_p = 1
\]
(积公式, product formula). 这里 $p = \infty$ 时 $|\cdot|_p$ 是通常绝对值, 否则是 I, \S\ref{sec:1.2} 定义的 $p$-进赋值. 故对 $P = (a_0:a_1:\dots:a_n)\in\mathbb{P}^n(\mathbb{Q})$,
\[
H(P) \overset{\text{def}}{=} \prod_{p = 2, \dots, \infty}\max_i(|a_i|_p)
\]
与 $P$ 代表元的选取无关. 且当选取本原代表元 $(a_0, \dots, a_n)$ 时, 对所有 $p\neq\infty$ 有 $\max_i|a_i|_p = 1$, 故 $H(P) = \max_i|a_i|_\infty$, 与先前定义一致. 对数域 $K$, 可归一化赋值使积公式成立 (见 \cite[8.8]{milne-ant}), 则
\[
H(P) \overset{\text{def}}{=} \prod_v \max_i(|a_i|_v), \qquad P = (a_0:a_1:\dots:a_n),
\]
给出 $\mathbb{P}^n(K)$ 上良好的极高概念. 用此定义, 本节所有结果都推广到数域上的椭圆曲线.
\end{remark}

\begin{remarkstar}
数域 $K$ 上射影簇的每个射影嵌入 $X\hookrightarrow\mathbb{P}^n$ 都在 $X(K)$ 上定义高度函数, 对应同一线丛的两个嵌入定义的对数高度函数相差有界量.\footnote{校对注：原书未限定两个嵌入所对应的线丛，但任意两个射影嵌入的高度未必相差有界量. 例如将 $\mathbb{P}^1$ 的恒等嵌入换为二次 Veronese 嵌入 $(a:b)\mapsto(a^2:ab:b^2)$，对数高度变为原来的两倍；此处补明比较高度时应固定的条件.} 在 1958 年国际数学家大会的短通讯中, N\'eron 猜想: 当 $X$ 是 阿贝尔簇时, 可能的高度函数 $X(K)\to\mathbb{R}$ 中有一个是二次型. Tate 于 1962 年用上述论证证明了它. N\'eron 后来找到二次高度的另一种构造, 更复杂, 但优点是把它表示为局部高度之和.
\end{remarkstar}

\subsection*{有限基定理的证明}

我们证明下面更精确的陈述.

\begin{proposition} \label{thm:4.4.13} 
设 $C > 0$ 使 $S \overset{\text{def}}{=} \{P\in E(\mathbb{Q}) \mid \hat{h}(P)\le C\}$ 含 $2E(\mathbb{Q})$ 在 $E(\mathbb{Q})$ 中的陪集代表元集, 则 $S$ 生成 $E(\mathbb{Q})$.
\end{proposition}

\begin{proof}
设存在 $Q\in E(\mathbb{Q})$ 不属于 $S$ 生成的子群. 因 $\hat{h}$ 取离散值, 可选 $Q$ 使 $\hat{h}(Q)$ 尽可能小. 由 $S$ 的定义, 存在 $P\in S$ 使 $Q = P + 2R$, $R\in E(\mathbb{Q})$. 显然 $R$ 不能属于 $S$ 生成的子群, 故 $\hat{h}(R)\ge\hat{h}(Q)$. 于是
\begin{align*}
    2\hat{h}(P) & = \hat{h}(P+Q) + \hat{h}(P-Q) - 2\hat{h}(Q)\qquad\text{(由~\ref{thm:4.4.11})}\\
    & \ge 0 + \hat{h}(2R) - 2\hat{h}(Q)  \qquad (\text{因 } h(\cdot)\ge 0)\\
    & = 4\hat{h}(R) - 2\hat{h}(Q)\qquad \text{(见~\ref{thm:4.4.5}b)}\\
    & \ge 2\hat{h}(Q),
\end{align*}
这与 $\hat{h}(P)\le C$ 且 $\hat{h}(Q) > C$ 矛盾.
\end{proof}

由引理 \ref{thm:4.4.12}, 该论证对任何数域 $K$ 无需修改即可成立.

\begin{remark} \label{thm:4.4.14} 
设 $C$ 如命题的陈述中, 令 $P_1, \dots, P_s$ 为 $E(\mathbb{Q})/2E(\mathbb{Q})$ 中元素的代表元集. 设 $Q\in E(\mathbb{Q})$ 有 $\hat{h}(Q) > C$, 写 $Q = P_i + 2R$, $R\in E(\mathbb{Q})$. 则 $\hat{h}(R) < \hat{h}(Q)$, 否则命题中的论证会给出矛盾. 故 $C$ 具有证明概述中使用的性质.
\end{remark}

\section{\texorpdfstring{计算 $E(\mathbb{Q})$ 的秩的问题}{计算 E(Q) 的秩的问题}}\label{sec:4.5}

据 Weil 的说法, 数学中两个最古老的未解决问题之一就是寻找确定群 $E(\mathbb{Q})$ 的算法.\footnote{以同余数问题的形式. 另一个是是否存在奇完全数的问题. 我曾向 Weil 建议后者不太有趣, 但他回答说在它被回答之前我们不会知道.} 我们知道 $E(\mathbb{Q})$ 有限生成, 故
\[
E(\mathbb{Q})\cong E(\mathbb{Q})_{\mathrm{tors}}\oplus\mathbb{Z}^r
\]
对某个称为秩 (rank) 的 $r\ge 0$ 成立. 由于我们知道如何计算 $E(\mathbb{Q})_{\mathrm{tors}}$ (见II, \ref{thm:2.5.1}), 问题归结为求 $r$ 的算法, 或更好, 求 $E(\mathbb{Q})/E(\mathbb{Q})_{\mathrm{tors}}$ 的基的算法. 我们把 Selmer 群 $S^{(n)}(E/\mathbb{Q})$ 视为 $r$ 的可计算上界, 以 $\Sha(E/\mathbb{Q})_n$ 为误差项 (见基本正合列~\eqref{eq:29}). 问题在于确定 $E(\mathbb{Q})$ 在 $S^{(n)}(E/\mathbb{Q})$ 中的像.

固定素数 $p$. 从
\[\begin{tikzcd}[column sep=large]
	0 & {E[p^n]} & E & E & 0 \\
	0 & {E[p^{n+1}]} & E & E & 0
	\arrow[from=1-1, to=1-2]
	\arrow[from=1-2, to=1-3]
	\arrow["{p^n}", from=1-3, to=1-4]
	\arrow[from=1-4, to=1-5]
	\arrow[from=2-1, to=2-2]
	\arrow["p", from=2-2, to=1-2]
	\arrow[from=2-2, to=2-3]
	\arrow["p", from=2-3, to=1-3]
	\arrow["{p^{n+1}}", from=2-3, to=2-4]
	\arrow["1", from=2-4, to=1-4]
	\arrow[from=2-4, to=2-5]
\end{tikzcd}\]
得交换图
\[\begin{tikzcd}
	0 & {E(\mathbb{Q})/p^n E(\mathbb{Q})} & {H^1(\mathbb{Q}, E_{p^n})} & {H^1(\mathbb{Q}, E)_{p^n}} & 0 \\
	0 & {E(\mathbb{Q})/p^{n+1} E(\mathbb{Q})} & {H^1(\mathbb{Q}, E_{p^{n+1}})} & {H^1(\mathbb{Q}, E)_{p^{n+1}}} & 0
	\arrow[from=1-1, to=1-2]
	\arrow[from=1-2, to=1-3]
	\arrow[from=1-3, to=1-4]
	\arrow[from=1-4, to=1-5]
	\arrow[from=2-1, to=2-2]
	\arrow[from=2-2, to=1-2]
	\arrow[from=2-2, to=2-3]
	\arrow[from=2-3, to=1-3]
	\arrow[from=2-3, to=2-4]
	\arrow["p"', from=2-4, to=1-4]
	\arrow[from=2-4, to=2-5]
\end{tikzcd}\]
从而得交换图
\[\begin{tikzcd}
	0 & {E(\mathbb{Q})/pE(\mathbb{Q})} & {S^{(p)}(E/\mathbb{Q})} & {\Sha(E/\mathbb{Q})_p} & 0 \\
	0 & {E(\mathbb{Q})/p^2E(\mathbb{Q})} & {S^{(p^2)}(E/\mathbb{Q})} & {\Sha(E/\mathbb{Q})_{p^2}} & 0 \\
	& \vdots & \vdots & \vdots \\
	0 & {E(\mathbb{Q})/p^nE(\mathbb{Q})} & {S^{(p^n)}(E/\mathbb{Q})} & {\Sha(E/\mathbb{Q})_{p^n}} & 0
	\arrow[from=1-1, to=1-2]
	\arrow[from=1-2, to=1-3]
	\arrow[from=1-3, to=1-4]
	\arrow[from=1-4, to=1-5]
	\arrow[from=2-1, to=2-2]
	\arrow[from=2-2, to=1-2]
	\arrow[from=2-2, to=2-3]
	\arrow[from=2-3, to=1-3]
	\arrow[from=2-3, to=2-4]
	\arrow["p", from=2-4, to=1-4]
	\arrow[from=2-4, to=2-5]
	\arrow[from=3-2, to=2-2]
	\arrow[from=3-3, to=2-3]
	\arrow["p", from=3-4, to=2-4]
	\arrow[from=4-1, to=4-2]
	\arrow[from=4-2, to=3-2]
	\arrow[from=4-2, to=4-3]
	\arrow[from=4-3, to=3-3]
	\arrow[from=4-3, to=4-4]
	\arrow["p", from=4-4, to=3-4]
	\arrow[from=4-4, to=4-5]
\end{tikzcd}\]
左边的垂直映射是自然商映射. 记 $S^{(p,n)}(E/\mathbb{Q})$ 为 $S^{(p^n)}(E/\mathbb{Q})$ 在 $S^{(p)}(E/\mathbb{Q})$ 中的像.

\begin{proposition} \label{thm:4.5.1} 
有包含
\[
E(\mathbb{Q})/pE(\mathbb{Q})\subset\bigcap_n S^{(p,n)}(E/\mathbb{Q}),
\]
当 $\Sha(E/\mathbb{Q})$ 中没有被 $p$ 的所有幂整除的非零元素时等式成立. 在后一情形,
\[
E(\mathbb{Q})/pE(\mathbb{Q})\cong S^{(p,n_0)}(E/\mathbb{Q})
\]
对一切充分大的 $n_0$ 成立.
\end{proposition}

\begin{proof}
因上面图中的左垂直箭头是满射, $E(\mathbb{Q})/pE(\mathbb{Q})$ 在 $S^{(p)}$ 中的像等于 $E(\mathbb{Q})/p^nE(\mathbb{Q})$ 的像, 由图的交换性它包含在 $S^{(p,n)}$ 中. 反之, 设 $\xi$ 位于 $\bigcap_n S^{(p,n)}$, 则对每个 $n$, 存在映射到 $\xi$ 的 $\xi_n\in S^{(p^n)}$. 设 $\delta_n$ 为 $\xi_n$ 在 $\Sha_{p^n}$ 中的像. 则对所有 $n$, $p^{n-1}\delta_n = \delta_1$, 故 $\delta_1$ 被 $p$ 的所有幂整除. 若 $\Sha$ 中唯一这样的元素为零, 则 $\xi$ 在 $E(\mathbb{Q})/pE(\mathbb{Q})$ 的像中.

群 $\Sha$ 是挠群, 且对所有 $n$, $\Sha_{p^n}$ 有限. 若 $\Sha$ 中没有被 $p$ 的所有幂整除的非零元素, 则 $\Sha$ 的 $p$-主分量 $\Sha(p)$ 有限.\footnote{设 $A$ 为挠 Abel 群, 对所有 $n$, $A_{p^n}$ 有限且非零. 对每个 $N$, 设 $A^{(N)}$ 为满足 $a_n\in A_{p^n}$ (对所有 $n$) 且对 $n\le N$ 有 $pa_n = a_{n-1}$ 的无限序列 $a_1, \dots, a_n, \dots$ 的集合. 每个 $A^{(N)}$ 是紧空间 $\prod_n A_{p^n}$ 的非空闭子集, 且 $A^{(1)}\supset A^{(2)}\supset\cdots$. 拓扑学的标准结果表明 $\bigcap_N A^{(N)}$ 非空, 故 $A$ 含被 $p$ 的所有幂整除的非零元素.} 因此存在 $n_0$ 使 $p^{n_0-1}\Sha(p) = 0$. 则 $S^{(p,n_0)}$ 在 $\Sha_p$ 中映射为零, 故
\[
E(\mathbb{Q})/pE(\mathbb{Q})\cong S^{(p,n_0)}(E/\mathbb{Q}).
\]
\end{proof}

\begin{remark} \label{thm:4.5.2} 
这给出计算 $r$ 的策略. 计算 $S^{(p)}$, 让计算机整夜运行计算由高度 $h(P)\le 10$ 的点生成的 $E(\mathbb{Q})$ 的子群 $T^{(1)}$. 若 $T^{(1)}$ 映满 $S^{(p)}$, 就找到了 $r$, 甚至找到 $E(\mathbb{Q})$ 的一组生成元. 若没有, 计算 $S^{(p^2)}$, 再让计算机整夜运行计算由高度 $h(P)\le 10^2$ 的点生成的 $E(\mathbb{Q})$ 的子群 $T^{(2)}$. 若 $T^{(2)}$ 在 $S^{(p)}$ 中的像是 $S^{(p,2)}$, 就找到了 $r$. 若没有, 继续……

噩梦般的可能性: Tate--Shafarevich 群含有被 $p$ 的所有幂整除的非零元素, 此时``我们注定要永远计算下去'' (\cite[p.~193]{tate1974}). 例如 $\Sha(E/\mathbb{Q})$ 含 $\mathbb{Q}_p/\mathbb{Z}_p$ 的拷贝时就会如此. 人们广泛猜想这不会发生.
\end{remark}

\begin{conjecture} \label{thm:4.5.3} 
Tate--Shafarevich 群总是有限的.
\end{conjecture}

等价地, $\Sha(E/\mathbb{Q})$ 无非零可除元素, 且对除有限多个 $p$ 外的所有 $p$, $\Sha(E/\mathbb{Q})_p = 0$.

直到 Kolyvagin\index{Kolyvagin} 与 Rubin\index{Rubin} 1987 年左右的工作之前 (\cite{kolyvagin1988a, kolyvagin1988b}; \cite{rubin1987}), 对 $\mathbb{Q}$ 上单条椭圆曲线都不知 $\Sha$ 有限. 在函数域上, 有限 Tate--Shafarevich 群的椭圆曲线例子在 20 年前就已知道 (\cite{milne1968})——见 \S\ref{sec:4.10}.

\subsection*{Selmer 猜想与 Cassels 配对}

基本正合列~\eqref{eq:29} 表明 $\Sha_m$ 度量了用 $S^{(m)}$ 估计 $E(\mathbb{Q})/mE(\mathbb{Q})$ 时的误差. 借助电子计算机, \cite{selmer1954} 在大量情形中经验性地发现, 由 $S^{(m)}$ 与 $S^{(m^2)}$ 得到的 $E(\mathbb{Q})$ 生成元数估计之差是偶数. 他猜想这总是对的.

Selmer 观察的可能解释如下. 设 $\Sha(E/\mathbb{Q})$ 有限, 且允许取值于 $\mathbb{Q}/\mathbb{Z}$ 的、以 $l(\xi, \xi) = 0$ ($\forall\xi\in\Sha$) 为交错的非退化双线性型 $l$. 则 $\Sha$ 的阶是平方.\footnote{若 $\Sha$ 被 $p$ 杀掉, 从而是 $\mathbb{F}_p$ 上的向量空间, 这是线性代数的标准结果. 有限 Abel 群的一般陈述可由此经归纳论证推出.} 当限制在 $\Sha_m$ 上时, 双线性型 $l$ 的核是子群 $\Sha_m\cap m\Sha$. 这蕴含 $\Sha_m\cap m\Sha$ 在 $\Sha_m$ 中的指标是平方, 从而蕴含 Selmer 猜想.

此论证激励 Cassels 证明下面的定理.

\begin{theorem}[Cassels 1962b] \label{thm:4.5.4} 
设 $E$ 为 $\mathbb{Q}$ 上的椭圆曲线. 在 $\Sha(E/\mathbb{Q})$ 上存在取值于 $\mathbb{Q}/\mathbb{Z}$ 的典范交错双线性型 $l$, 其核恰是 $\Sha(E/\mathbb{Q})$ 的可除元素群.
\end{theorem}

这可能是 Cassels 对该主题最持久的贡献. 由定理, 若 $\Sha(E/\mathbb{Q})$ 没有被 $p$ 的所有幂整除的元素, 则 Selmer 猜想对 $p$ 成立. 换言之, 若 Selmer 猜想对单个 $p$ 失败, 则猜想 \ref{thm:4.5.3} 为假.

\begin{remarkstar}
设 $C$ 为使序列
\[
0\to\Sha(E/\mathbb{Q})\to H^1(\mathbb{Q}, E)\to\bigoplus_{p = 2, \dots, \infty}H^1(\mathbb{Q}_p, E)\to C\to 0
\]
正合的群, 并给序列中每个群配离散拓扑. 当 $\Sha(E/\mathbb{Q})$ 有限时, \cite{cassels1964} 证明该序列的 Pontryagin 对偶是正合序列
\begin{equation} \label{eq:34}
0\to\widehat{\Sha}(E/\mathbb{Q})\to\prod_{p = 2, \dots, \infty}H^1(\mathbb{Q}_p, E)\to\widehat{E(\mathbb{Q})}\to 0,
\tag{34}
\end{equation}
其中 $\widehat{E(\mathbb{Q})}$ 是 $E(\mathbb{Q})$ 的 pro-有限完备化, 即 $\widehat{E(\mathbb{Q})} = \varprojlim_n E(\mathbb{Q})/nE(\mathbb{Q})$.

Tate 把 Cassels 的工作推广为 阿贝尔簇与其对偶的 Tate--Shafarevich 群之间的对偶 (Cassels--Tate 对偶), 并为 阿贝尔簇构造了对偶正合序列~\eqref{eq:34} (Cassels--Tate 对偶正合序列). 见 \cite{tate1963}.
\end{remarkstar}

\subsection*{秩的显式计算}

显式计算 $E(\mathbb{Q})$ 的秩 $r$ 可能困难\footnote{当然, 人们总是可以在 Pari 中输入 ellgenerators, 希望别人已经做了艰苦工作.} (也许不可能), 但偶尔它直截了当. 为避免处理 $\mathbb{Q}$ 以外数域 $L$ 的问题, 我们假设椭圆曲线在 $\mathbb{Q}$ 上的所有 2 阶点都有理, 故它由方程
\[
Y^2Z = (X-\alpha Z)(X-\beta Z)(X-\gamma Z), \qquad \alpha, \beta, \gamma \text{ 两两不同的整数},
\]
给出. $(X-\alpha)(X-\beta)(X-\gamma)$ 的判别式是
\[
\Delta = (\alpha-\beta)^2(\beta-\gamma)^2(\gamma-\alpha)^2.
\]

\begin{proposition} \label{thm:4.5.5} 
$E(\mathbb{Q})$ 的秩 $r$ 满足不等式
\[
r\le 2\cdot\#\{p \mid p \text{ 整除 } 2\Delta\}.
\]
\end{proposition}

\begin{proof}
设 $T = E(\mathbb{Q})_{\mathrm{tors}}$. 因 $E(\mathbb{Q})\cong T\oplus\mathbb{Z}^r$, 有 $E(\mathbb{Q})/2E(\mathbb{Q})\cong T/2T\oplus(\mathbb{Z}/2\mathbb{Z})^r$. 因 $T$ 有限, 映射 $T\xrightarrow{2}T$ 的核与余核同阶, 故 $T/2T\cong(\mathbb{Z}/2\mathbb{Z})^2$. 有单射
\[
E(\mathbb{Q})/2E(\mathbb{Q})\hookrightarrow(\mathbb{Q}^{\times}/\mathbb{Q}^{\times2})^2,
\]
其像包含在由 $-1$ 与 $E$ 有坏约化 (即整除 $2\Delta$) 的素数生成的 $\mathbb{Q}^{\times}/\mathbb{Q}^{\times2}$ 的子群的乘积中 (见~\ref{thm:4.3.7}).
\end{proof}

可以改进此估计. 设 $T_1$ 为使约化为结点型的整除 $\Delta$ 的素数集, $T_2$ 为使约化为尖点型的素数集. 故 $T_1$ 由使 $(X-\alpha)(X-\beta)(X-\gamma)$ 的根中两个模该素数重合的那些素数组成, $T_2$ 由三个根都重合的那些素数组成. 分别记 $T_1$ 与 $T_2$ 的元素数为 $t_1$ 与 $t_2$.

\begin{proposition} \label{thm:4.5.6} 
$E(\mathbb{Q})$ 的秩 $r$ 满足 $r\le t_1 + 2t_2 - 1$.
\end{proposition}

\begin{proof}
如 (\ref{thm:4.3.7}) 定义 $\varphi_\alpha\colon E(\mathbb{Q})/2E(\mathbb{Q})\to\mathbb{Q}^{\times}/\mathbb{Q}^{\times2}$:
\[
\varphi_\alpha((x:y:z)) = \begin{cases}
(x/z - \alpha)\mathbb{Q}^{\times2}, & z\neq 0, x\neq\alpha z,\\
(\alpha-\beta)(\alpha-\gamma)\mathbb{Q}^{\times2}, & z\neq 0, x = \alpha z,\\
\mathbb{Q}^{\times2}, & (x:y:z) = (0:1:0).
\end{cases}
\]
类似定义 $\varphi_\beta$——映射
\[
P\mapsto(\varphi_\alpha(P), \varphi_\beta(P))\colon E(\mathbb{Q})/2E(\mathbb{Q})\to(\mathbb{Q}^{\times}/\mathbb{Q}^{\times2})^2
\]
是单射. 对每个素数 $p$, 令 $\varphi_p(P)$ 为 $(\mathbb{Z}/2\mathbb{Z})^2$ 中分量是 $\ord\nolimits_p(\varphi_\alpha(P))\bmod 2$ 与 $\ord\nolimits_p(\varphi_\beta(P))\bmod 2$ 的元素, 令 $\varphi_\infty(P)$ 为 $\{\pm\}^2$ 中分量是 $\operatorname{sign}(\varphi_\alpha(P))$ 与 $\operatorname{sign}(\varphi_\beta(P))$ 的元素.

为证命题, 只需证以下陈述:
\begin{enumerate}[label=\textup{(\alph*)}]
\item 若 $p\nmid\Delta$, 则对所有 $P$, $\varphi_p(P) = 0$;
\item 若 $p\in T_1$, 则对所有 $P$, $\varphi_p(P)$ 包含在 $\mathbb{F}_2^2$ 的对角线中;
\item 当 $\alpha < \beta < \gamma$ 排序时, 对所有 $P$, $\varphi_\infty(P)$ 等于 $(+,+)$ 或 $(+,-)$.
\end{enumerate}
除 $p = 2$ 外, (a) 在 (\ref{thm:4.3.7}) 前一段已证.

我们证 $P = (x:y:1)$, $x\neq\alpha,\beta,\gamma$ 且 $\alpha\equiv\beta\pmod{p}$ 情形下的 (b). 令
\[
a = \ord\nolimits_p(x-\alpha), \qquad b = \ord\nolimits_p(x-\beta), \qquad c = \ord\nolimits_p(x-\gamma).
\]
因 $(x-\alpha)(x-\beta)(x-\gamma)$ 是平方, $a + b + c\equiv 0\pmod{2}$.

若 $a < 0$, 则 (因 $\alpha\in\mathbb{Z}$) $p^{-a}$ 作为 $x$ (最简形式) 分母的因子出现, 由此 $b = a = c$. 因 $a+b+c\equiv 0\pmod{2}$, 这蕴含 $a\equiv b\equiv c\equiv 0\pmod{2}$, 故 $\varphi_p(P) = 0$. 若 $b < 0$ 或 $c < 0$, 同样的论证适用.

若 $a > 0$, 则 $p$ 整除 $x-\alpha$ 的分子. 因 $p$ 不整除 $(\alpha-\gamma)$, 它不整除 $(\alpha-\gamma) + (x-\alpha) = (x-\gamma)$, 故 $c = 0$. 现在 $a + b\equiv 0\pmod{2}$ 蕴含 $\varphi_p(P)$ 位于 $\mathbb{F}_2^2$ 的对角线中. 若 $b > 0$ 或 $c > 0$, 类似论证适用.

(b) 的其余情形类似证明.

证 (c). 设 $P = (x:y:1)$, $x\neq\alpha,\beta,\gamma$. 可设 $\alpha < \beta < \gamma$, 故 $x-\alpha > x-\beta > x-\gamma$. 则 $\varphi_\infty(P) = (+,+)$, $(+,-)$ 或 $(-,-)$. 但因 $(x-\alpha)(x-\beta)(x-\gamma)$ 是 $\mathbb{Q}$ 中的平方, 序对 $(-,-)$ 不可能. $x = \alpha$ 等情形同样容易.
\end{proof}

\begin{example} \label{thm:4.5.7} 
曲线
\[
E\colon Y^2Z = X^3 - XZ^2
\]
属于上述形式, $(\alpha, \beta, \gamma) = (-1, 0, 1)$. 唯一坏素数是 $2$, 而此处约化为结点型. 故 $r = 0$, $E$ 无无限阶点,
\[
E(\mathbb{Q})\cong(\mathbb{Z}/2\mathbb{Z})^2.
\]
\end{example}

\begin{aside} \label{thm:4.5.8} 
$E(\mathbb{Q})$ 的秩是否有界或可任意大, 是一个古老但仍未解决的问题.\footnote{\cite[p.~257]{cassels1966} 写道: ``人们曾广泛猜想秩有只依赖基域的上界. 这在我看来不可信, 因为理论清楚地表明 [椭圆曲线] 只有在由系数非常大的方程定义时才可能有高秩.'' \cite[p.~194]{tate1974} 写道``我猜测秩没有界.'' 此后, 观点又转回认为秩有界. \cite{park2019} 提出启发式, 表明 $\mathbb{Q}$ 上秩大于 $21$ 的椭圆曲线只有有限多个. 也见 \cite{poonen2018}.} 目前 (2020) 已知的最高秩是 $28$——这自本书第 1 版 (2006) 以来没有改变. 对有限域 $k$ 上 $k(T)$ 上的椭圆曲线, 已知秩可任意大 (\cite{tateshafarevich1967}).
\end{aside}

\begin{exercise} \label{thm:4.5.9} 
做下面两个问题之一 (知道二次互反律的应做 \hyperref[item:4.5.9.2]{\ref*{item:4.5.9.2}}).
\begin{enumerate}[label=\textup{(\arabic*)}]
\item 证明若 $E$ 有方程 $Y^2Z = X^3 - 4XZ^2$, 则 $E(\mathbb{Q})$ 有限. 提示: 设 $P$ 为 $E(\mathbb{Q})$ 中无限阶点, 证明在可能把 $P$ 换成 $P + Q$ ($2Q = 0$) 后, $\varphi_2(P)$ 为零. 然后证明 $\varphi_\infty(P) = (+,+)$——矛盾.
\item\label{item:4.5.9.2} 设 $E$ 为椭圆曲线 $Y^2Z = X^3 - p^2XZ^2$, $p$ 为奇素数. 证明 $E(\mathbb{Q})$ 的秩 $r$ 满足:
\[
r\le 2 \quad \text{若 } p\equiv 1\pmod{8}, \qquad r = 0 \quad \text{若 } p\equiv 3\pmod{8}, \qquad r\le 1 \quad \text{否则}.
\]
提示: 设 $P$ 为 $E(\mathbb{Q})$ 中无限阶点, 证明在可能把 $P$ 换成 $P + Q$ ($2Q = 0$) 后, $\varphi_p(P)$ 为零.
\end{enumerate}
这些是标准例子.
\end{exercise}

\section{N\'eron--Tate 配对}\label{sec:4.6}

我们在 \S\ref{sec:4.4} 中看到有典范的 $\mathbb{Z}$-双线性配对
\[
B\colon E(\mathbb{Q})\times E(\mathbb{Q})\to\mathbb{R}, \qquad
B(x, y) = \hat{h}(x+y) - \hat{h}(x) - \hat{h}(y).
\]
此配对唯一地延拓为 $\mathbb{R}$-双线性配对
\[
B\colon E(\mathbb{Q})\otimes\mathbb{R}\times E(\mathbb{Q})\otimes\mathbb{R}\to\mathbb{R}.
\]
若 $\{e_1, \dots, e_r\}$ 是 $E(\mathbb{Q})/E(\mathbb{Q})_{\mathrm{tors}}$ 的 $\mathbb{Z}$-基, 则 $\{e_1\otimes 1, \dots, e_r\otimes 1\}$ 是 $E(\mathbb{Q})\otimes_{\mathbb{Z}}\mathbb{R}$ 的 $\mathbb{R}$-基, 相对于它 $B$ 有矩阵 $(B(e_i, e_j))$.

\begin{theorem} \label{thm:4.6.1} 
双线性配对
\[
B\colon E(\mathbb{Q})\otimes\mathbb{R}\times E(\mathbb{Q})\otimes\mathbb{R}\to\mathbb{R}
\]
是正定\index{zhengding@正定(positive definite)}的; 特别地, 它非退化.
\end{theorem}

这由定理 \ref{thm:4.4.7} 与下一个引理推出. $\mathbb{R}$-向量空间 $V$ 中的格 (lattice) 指由 $V$ 的基生成的 $\mathbb{Z}$-子模 (有时称为满格, full lattice).

\begin{lemma} \label{thm:4.6.2} 
设 $q\colon V\to\mathbb{R}$ 为有限维 $\mathbb{R}$-向量空间 $V$ 上的二次型. 若存在 $V$ 中的格 $\Lambda$ 满足
\begin{enumerate}[label=\textup{(\alph*)}]
\item 对每个常数 $C$, 集合 $\{P\in\Lambda \mid q(P)\le C\}$ 有限,
\item $\Lambda$ 中满足 $q(P) = 0$ 的唯一 $P$ 是 $P = 0$,
\end{enumerate}
则 $q$ 在 $V$ 上正定.
\end{lemma}

\begin{proof}
存在 $V$ 的基使 $q$ 取如下形式:
\[
q(X) = X_1^2 + \cdots + X_s^2 - X_{s+1}^2 - \cdots - X_t^2, \qquad t\le\dim V.
\]
假设 $s\neq\dim V$ 并推出矛盾. 用该基把 $V$ 等同于 $\mathbb{R}^n$. 设 $\delta$ 为 $\Lambda$ 中最短向量的长度, 即
\[
\delta = \inf\{q(P) \mid P\in\Lambda, P\neq 0\}.
\]
条件 (a, b) 蕴含 $\delta > 0$. 考虑集合
\[
B(\varepsilon) = \Bigl\{(x_i)\in\mathbb{R}^n \;\Big|\; x_1^2 + \cdots + x_s^2\le\frac{\delta}{2}, \ x_{s+1}^2 + \cdots + x_t^2\le\varepsilon\Bigr\}.
\]
$B(\varepsilon)$ 中每个向量 (用 $q$) 的长度 $\le\delta/2$, 故 $B(\varepsilon)\cap\Lambda = \{0\}$, 但取 $\varepsilon$ 大可使 $B(\varepsilon)$ 的体积任意大, 从而违反下面著名的 Minkowski 定理\index{dingli@定理(theorem)!Minkowski@Minkowski(Minkowski)}.
\end{proof}

\begin{theorem}[Minkowski] \label{thm:4.6.3} 
设 $\Lambda$ 为 $\mathbb{R}^n$ 中带基本平行六面体 $D_0$ 的格, $B$ 为 $\mathbb{R}^n$ 中紧、凸、关于原点对称的子集. 若
\[
\operatorname{Vol}(B)\ge 2^n\operatorname{Vol}(D_0),
\]
则 $B$ 含 $\Lambda$ 中异于原点的点.
\end{theorem}

\begin{proof}
先证 $\mathbb{R}^n$ 中满足 $\operatorname{Vol}(S) > \operatorname{Vol}(D_0)$ 的可测集 $S$ 含不同的点 $\alpha, \beta$ 使 $\alpha-\beta\in\Lambda$. 显然
\[
\operatorname{Vol}(S) = \sum\operatorname{Vol}(S\cap D),
\]
其中和对 $D_0$ 被 $\Lambda$ 中元素平移所得的所有 $D$ 进行. 基本平行六面体 $D_0$ 将含每个集合 $S\cap D$ 的 (被 $\Lambda$ 中元素平移所得) 唯一平移. 因 $\operatorname{Vol}(S) > \operatorname{Vol}(D_0)$, 这些集合中至少有两个重叠, 故存在元素 $\alpha, \beta\in S$ 使
\[
\alpha - \lambda = \beta - \lambda', \qquad \text{某些不同的 } \lambda, \lambda'\in\Lambda.
\]
则 $\alpha-\beta = \lambda-\lambda'\in\Lambda\setminus\{0\}$.

把此应用于 $S = \frac{1}{2}B \overset{\text{def}}{=} \{\frac{x}{2} \mid x\in B\}$. 它的体积为 $\frac{1}{2^n}\operatorname{Vol}(B) > \operatorname{Vol}(D_0)$, 故存在不同的 $\alpha, \beta\in B$ 使 $\alpha/2 - \beta/2\in\Lambda$. 因 $B$ 关于原点对称, $-\beta\in B$, 且因它凸, $(\alpha + (-\beta))/2\in B$.
\end{proof}

\begin{remark} \label{thm:4.6.4} 
由实向量空间 $V$、$V$ 中的格 $\Lambda$ 与 $V$ 上的正定二次型 $q$ 组成的系统在数学中极受关注——它们通常就简称为格 (lattices). 存在 $V$ 的基把 $(V, q)$ 等同于 $(\mathbb{R}^n, X_1^2 + \cdots + X_n^2)$. 寻找 $\mathbb{R}^n$ 中以格的点为中心的球堆积 (sphere packing)\index{qiuduiji@球堆积(sphere packing)} 的稠密堆积, 归结为寻找使
\[
\frac{\|\text{最短向量}\|^n}{\operatorname{Vol}(\text{基本平行六面体})}
\]
大的格 $\Lambda$. 许多格, 例如 Leech 格, 有非常有趣的自同构群. 见 \cite{conway1999}.

$\mathbb{Q}$ 上椭圆曲线 $E$ 给出这样的系统, 即 $V = E(\mathbb{Q})\otimes\mathbb{R}$, $\Lambda = E(\mathbb{Q})/E(\mathbb{Q})_{\mathrm{tors}}$, $q = \hat{h}$. 据我所知, 它们并不有趣——目前没有已知的椭圆曲线有 $\operatorname{rank}(E(\mathbb{Q})) > 28$. 但把数域换成一元函数域时, 得到非常有趣的格的无限族. 我们在下面的 \S\ref{sec:4.11} 中讨论.
\end{remark}

\section{上同调群的几何解释；Jacobi 簇}\label{sec:4.7}

本节全文中, $k$ 为完美域, $\Gamma = \Gal(k^{\mathrm{al}}/k)$.

对 $k$ 上的椭圆曲线 $E$ 与整数 $n > 1$, 我们有上同调群的正合列 (见 \S\ref{sec:4.2}):
\[
0\to E(k)/nE(k)\to H^1(k, E_n)\to H^1(k, E)_n\to 0.
\]
这里 $H^1(k, E_n)$ 与 $H^1(k, E)$ 定义为从 $\Gamma$ 到 $E(k^{\mathrm{al}})_n$ 与 $E(k^{\mathrm{al}})$ 的连续交叉同态模主交叉同态的群. 本节给出这些群——从而也给 Selmer 与 Tate--Shafarevich 群——的几何解释. 我们将给 $k$ 上亏格 $1$ 的曲线 $W$ (可能没有有理点) 配上椭圆曲线 $E$, 称为它的 Jacobi 簇 (jacobian)\index{Jacobi@Jacobi 簇(jacobian variety)}, 并且将看到椭圆曲线 $E$ 的 Tate--Shafarevich 群对以 $E$ 为 Jacobi 簇、使 Hasse 原理失败的 $k$ 上亏格 $1$ 曲线分类, 即那些在每个 $\mathbb{Q}_p$ 与 $\mathbb{R}$ 中都有点但没有 $\mathbb{Q}$ 中点的曲线.

一般地, $H^1(k, G)$ 对在 $k^{\mathrm{al}}$ 上同构于自同构群为 $G$ 的固定对象的 $k$ 上对象分类. 我们将看到几个例子.

\subsection*{集合的主齐性空间}

设 $A$ 为交换群. $A$ 的右 $A$-集 (right $A$-set)
\[
(w, a)\mapsto w + a\colon W\times A\to W
\]
若满足 $W\neq\varnothing$ 且映射
\[
(w, a)\mapsto(w, w+a)\colon W\times A\to W\times W
\]
双射, 即对每对 $w_1, w_2\in W$, 存在唯一的 $a\in A$ 使 $w_1 + a = w_2$, 则称为 $A$ 的主齐性空间 (principal homogeneous space).\index{zhuqixingkongjian@主齐性空间(principal homogeneous space)}

\begin{example} \label{thm:4.7.1} 
(a) 加法 $A\times A\to A$ 使 $A$ 成为 $A$ 的主齐性空间.
(b) 向量空间的主齐性空间 (例如按 Newton 的宇宙) 称为仿射空间\index{fangshekongjian@仿射空间(affine space)} (affine space). 本质上仿射空间是没有优选原点的向量空间.
\end{example}

主齐性空间的态射 $\varphi\colon W\to W'$ 只是 $A$-集的映射. 我们把下面的陈述留给读者验证.

\begin{aside} \label{thm:4.7.2} 
设 $W$ 与 $W'$ 为 $A$ 的主齐性空间.
\begin{enumerate}[label=\textup{(\alph*)}]
\item 对任何点 $w_0\in W$, $w_0'\in W'$, 存在唯一的态射 $\varphi\colon W\to W'$ 把 $w_0$ 送到 $w_0'$.
\item 每个态射 $W\to W'$ 都是同构 (即也有逆态射).
\end{enumerate}
\end{aside}

\begin{aside} \label{thm:4.7.3} 
设 $W$ 为 $A$ 的主齐性空间.
\begin{enumerate}[label=\textup{(\alph*)}]
\item 对任何点 $w_0\in W$, 存在唯一态射 $A\to W$ (主齐性空间的) 把 $0$ 送到 $w_0$.
\item $a\in A$ 定义 $W$ 的自同构 $w\mapsto w + a$, 而 $W$ 的每个自同构对唯一的 $a\in A$ 都有此形式.
\end{enumerate}
最后这个陈述表明 $\Aut(W)\cong A$, 故我们定义了一类以 $A$ 为自同构群的对象.
\end{aside}

\subsection*{曲线的主齐性空间}

设 $E$ 为域 $k$ 上的椭圆曲线. $E$ 的主齐性空间\footnote{也使用``torsor''一词. 主齐性空间是算术几何中拓扑学主丛 (principal bundles) 的类比.}是 $k$ 上的曲线 $W$ 连同由正则映射给出的 $E$ 的右作用
\[
(w, P)\mapsto w + P\colon W\times E\to W,
\]
使
\[
(w, P)\mapsto(w, w + P)\colon W\times E\to W\times W
\]
是代数簇的同构. 这些条件蕴含: 对每个域 $K\supset k$, 集合 $W(K)$ 要么为空, 要么是群 $E(K)$ 的主齐性空间.

\begin{example} \label{thm:4.7.4} 
加法 $E\times E\to E$ 使 $E$ 成为 $E$ 的主齐性空间. 同构于这个主齐性空间的主齐性空间称为平凡的 (trivial).
\end{example}

$E$ 的主齐性空间的态射是满足
\[
\begin{array}{ccc}
W\times E & \to & W\\
\downarrow{\varphi\times\mathrm{id}_E} & & \downarrow{\varphi}\\
W'\times E & \to & W'
\end{array}
\]
交换的正则映射 $\varphi\colon W\to W'$. 上一小节中的陈述经必要修改后推广到椭圆曲线的主齐性空间.

\begin{aside} \label{thm:4.7.5} 
设 $W$ 与 $W'$ 为 $E$ 的主齐性空间.
\begin{enumerate}[label=\textup{(\alph*)}]
\item 对任何域 $K\supset k$ 与点 $w_0\in W(K)$, $w_0'\in W'(K)$, 存在唯一态射 $\varphi\colon W\to W'$ (在 $K$ 上) 把 $w_0$ 送到 $w_0'$.
\item 每个态射 $W\to W'$ 都是同构.
\end{enumerate}
\end{aside}

\begin{aside} \label{thm:4.7.6} 
设 $W$ 为 $E$ 的主齐性空间. 对任何点 $w_0\in W(k)$, 存在唯一同态 $E\to W$ (主齐性空间的) 把 $0$ 送到 $w_0$. 故 $W$ 平凡当且仅当 $W(k)\neq\varnothing$. 因 $W$ 在 $k$ 的某个有限扩张中有坐标在该扩张中的点 (由零点定理), 它在该扩张上变平凡.
\end{aside}

\begin{aside} \label{thm:4.7.7} 
设 $W$ 为 $E$ 的主齐性空间. 点 $P\in E(K)$ 定义 $W$ 的自同构 $w\mapsto w + P$, 而 $W$ 在 $K$ 上的每个自同构对唯一的 $P\in E(K)$ 都有此形式. 特别地, $\Aut(W_{k^{\mathrm{al}}})\cong E(k^{\mathrm{al}})$.
\end{aside}

\subsection*{主齐性空间的分类}

设 $W$ 为 $E$ 在 $k$ 上的主齐性空间. 则 $\Gamma$ 连续作用在 $W(k^{\mathrm{al}})$ 上. 选点 $w_0\in W(k^{\mathrm{al}})$. 对 $\sigma\in\Gamma$, 对唯一的 $f(\sigma)\in E(k^{\mathrm{al}})$ 有 $\sigma w_0 = w_0 + f(\sigma)$. 比较等式
\begin{align*}
(\sigma\tau)w_0 &= w_0 + f(\sigma\tau) && \qquad (f \text{ 的定义})\\
(\sigma\tau)w_0 &= \sigma(\tau w_0) = \sigma(w_0 + f(\tau)) = \sigma w_0 + \sigma(f(\tau)) = w_0 + f(\sigma) + \sigma f(\tau),
\end{align*}
我们看到 $f$ 是交叉同态 $\Gamma\to E(k^{\mathrm{al}})$. 因 $w_0$ 的坐标在 $k$ 的有限扩张中, $f$ 连续. 第二个点 $w_1\in W(k^{\mathrm{al}})$ 由 $\sigma w_1 = w_1 + f_1(\sigma)$ 定义另一个交叉同态 $f_1$; 但 $w_1 = w_0 + P$ 对某个 $P\in E(\mathbb{Q}^{\mathrm{al}})$, 故
\[
\sigma w_1 = \sigma(w_0 + P) = \sigma w_0 + \sigma P = w_0 + f(\sigma) + \sigma P = w_1 + f(\sigma) + \sigma P - P,
\]
这表明 $f$ 与 $f_1$ 相差一个主交叉同态. 故 $f$ 的上同调类只依赖 $W$. 若此类为零, 则 $f(\sigma) = \sigma P - P$ 对某个 $P\in E(k^{\mathrm{al}})$, 且
\[
\sigma(w_0 - P) = \sigma w_0 - \sigma P = w_0 + \sigma P - P - \sigma P = w_0 - P.
\]
这蕴含 $w_0 - P\in W(k)$, 故 $W$ 是平凡主齐性空间\index{zhuqixingkongjian@主齐性空间(principal homogeneous space)!pingfan@平凡(trivial)} (见~\ref{thm:4.7.6}).

\begin{theorem} \label{thm:4.7.8} 
映射 $W\mapsto[f]$ 定义双射
\begin{equation} \label{eq:35}
\{E \text{ 的主齐性空间}\}\xrightarrow{\ \sim\ }H^1(k, E),
\tag{35}
\end{equation}
把平凡主齐性空间送到零元素.
\end{theorem}

\begin{proof}
{\emergencystretch=3em 设 $\varphi\colon W\to W'$ 为 $E$ 的主齐性空间 (在 $k$ 上) 的同构, $w_0\in W(k^{\mathrm{al}})$. 则 $(W, w_0)$ 与 $(W', \varphi(w_0))$ 定义相同的交叉同态, 故映射\par}
\[
\{E \text{ 的主齐性空间}\}\to H^1(k, E)
\]
在同构类上为常数.

若 $W$ 与 $W'$ 定义相同的上同调类, 则可选 $w_0\in W(k^{\mathrm{al}})$ 与 $w_0'\in W'(k^{\mathrm{al}})$ 使 $(W, w_0)$ 与 $(W', w_0')$ 定义相同的交叉同态. 由~\ref{thm:4.7.5}, 存在 $k^{\mathrm{al}}$ 上主齐性空间 $E$ 的唯一同构 $\varphi\colon W\to W'$ 把 $w_0$ 送到 $w_0'$. 设 $w\in W(k^{\mathrm{al}})$, 且 $w = w_0 + P$, $P\in E(k^{\mathrm{al}})$. 则
\begin{align*}
\varphi(\sigma w) &= \varphi(\sigma(w_0 + P)) = \varphi(\sigma w_0 + \sigma P) = \varphi(w_0 + f(\sigma) + \sigma P)\\
&= w_0' + f(\sigma) + \sigma P = \sigma w_0' + \sigma P = \sigma\varphi(w),
\end{align*}
这蕴含映射 $\varphi$ 定义在 $k$ 上 (见I, \ref{thm:1.5.5}). 故映射~\eqref{eq:35} 一对一.

最后设 $f\colon\Gamma\to E(k^{\mathrm{al}})$ 为连续交叉同态. 令 $W$ 表示代数曲线 $E_{k^{\mathrm{al}}}$, 但 $\Gamma$ 按如下规则作用:
\[
{}^{\sigma}P = \sigma P + f(\sigma), \qquad P\in W(k^{\mathrm{al}}) = E(k^{\mathrm{al}}).
\]
则
\[
\begin{aligned}
{}^{\sigma\tau}P
&= \sigma\tau P + f(\sigma\tau)
= \sigma\tau P + f(\sigma) + \sigma f(\tau),\\
{}^{\sigma}({}^{\tau}P)
&= \sigma\bigl(\tau P + f(\tau)\bigr) + f(\sigma)\\
&= \sigma\tau P + \sigma f(\tau) + f(\sigma),
\end{aligned}
\]
故这是 $\Gamma$ 在 $W(k^{\mathrm{al}})$ 上的作用. 该作用连续 (因 $f$ 连续) 且正则, 故来自 $W$ 在 $k$ 上的模型 $W_0$ (I\ref{thm:1.5.5}). 曲线 $W_0$ 有 $E$ 在 $k$ 上主齐性空间的自然结构.
\end{proof}

\begin{aside} \label{thm:4.7.9} 
记 $\operatorname{WC}(E/k)$ 为 $E$ 在 $k$ 上主齐性空间的同构类集. 定理中的双射在 $\operatorname{WC}(E/k)$ 上定义交换群结构, 可描述如下: 对主齐性空间 $W$ 与 $W'$, 定义 $W\wedge W'$ 为 $W\times W'$ 在 $E$ 的反对角作用下的商, 故
\[
(W\wedge W')(k^{\mathrm{al}}) = (W(k^{\mathrm{al}})\times W'(k^{\mathrm{al}}))/{\sim}, \qquad
(w, w')\sim(w + P, w' - P), \ P\in E(k^{\mathrm{al}}).
\]
\footnote{校对注：原书写为对角作用 $(w+P,w'+P)$. 若要使商代表上同调类之和，应取反对角作用 $(w+P,w'-P)$；选取几何基点后，商中的和坐标具有两余循环之和. 商上的 $E$-作用可由仅在第一因子上平移诱导.} 则 $W\wedge W'$ 有主齐性空间的自然结构, 且表示 $W$ 与 $W'$ 在 $\operatorname{WC}(E/k)$ 中的和.
\end{aside}

\subsection*{$H^1(\mathbb{Q}, E_n)$ 的几何解释}

现在给出 $H^1(k, E_n)$ 的几何解释. $n$-覆盖 (n-covering\index{n-fugai@n-覆盖(n-covering)}) 是由 $E$ 的主齐性空间 $W$ 与 (定义在 $k$ 上的) 正则映射 $\alpha\colon W\to E$ 组成的序对 $(W, \alpha)$, 具有如下性质:
存在 $w_1\in W(k^{\mathrm{al}})$ 使对所有 $P\in E(k^{\mathrm{al}})$ 有 $\alpha(w_1 + P) = nP$.
$n$-覆盖的态射 $(W, \alpha)\to(W', \alpha')$ 是主齐性空间 $E$ 的态射 $\varphi\colon W\to W'$ 使 $\alpha = \alpha'\circ\varphi$.

对 $\sigma\in\Gamma$, 写
\begin{equation} \label{eq:36}
\sigma w_1 = w_1 + f(\sigma), \qquad f(\sigma)\in E(k^{\mathrm{al}}).
\tag{36}
\end{equation}
如前, $f$ 是取值于 $E(k^{\mathrm{al}})$ 的交叉同态. 对 \eqref{eq:36} 两边应用 $\alpha$ 并利用
\[
\alpha(\sigma w_1) = \sigma(\alpha w_1) = \sigma(\alpha(w_1 + O)) = O,
\]
我们发现 $nf = 0$, 故 $f$ 取值于 $E_n(k^{\mathrm{al}})$. 元素 $w_1$ 在换为 $w_1 + Q$ ($Q\in E_n(k^{\mathrm{al}})$) 的意义下唯一, 由此 $f$ 在 $H^1(k, E_n)$ 中的类与 $w_1$ 的选取无关.

\begin{theorem} \label{thm:4.7.10} 
映射 $(W, \alpha)\mapsto[f]$ 定义双射
\[
\dfrac{\{n\text{-覆盖}\}}{\approx}\xrightarrow{\ \sim\ }H^1(k, E_n).
\]
\end{theorem}

\begin{proof}
记 $\operatorname{WC}(E_n/k)$ 为 $n$-覆盖模同构的集合. ``忘记 $\alpha$'' 的映射 $(W, \alpha)\mapsto W$ 定义满射
\[
\operatorname{WC}(E_n/k)\to\operatorname{WC}(E/k)_n,
\]
其纤维以自然方式是群 $E(k)/nE(k)$ 的主齐性空间. 例如, 若 $W$ 平凡, 从而存在 $w_0\in W(k)$, 则 $\alpha(w_0)\in E(k)$; 若 $w_0'\in W(k)$ 也, 则 $w_0' = w_0 + P$ 对某个 $P\in E(k)$, 且 $\alpha(w_0') = \alpha(w_0) + nP$, 故 $\alpha(w_0)$ 作为 $E(k)/nE(k)$ 中元素良定义. 现在考虑
\[\begin{tikzcd}
	&& {\operatorname{WC}(E_n/k)} & {\operatorname{WC}(E/k)_n} & \\
	0 & {E(k)/nE(k)} & {H^1(k,E_n)} & {H^1(k, E)_n} & 0
	\arrow["b", from=1-3, to=1-4]
	\arrow["a", from=1-3, to=2-3]
	\arrow["\simeq", from=1-4, to=2-4]
	\arrow[from=2-1, to=2-2]
	\arrow[from=2-2, to=2-3]
	\arrow["c", from=2-3, to=2-4]
	\arrow[from=2-4, to=2-5]
\end{tikzcd}\]
图交换, 故 $a$ 把 $b$ 的纤维映到 $c$ 的纤维中. 因这些纤维是 $E(k)/nE(k)$ 的主齐性空间, 由 (\ref{thm:4.7.2}), $a$ 在每个纤维上双射, 故在整个集合上双射.
\end{proof}

证明中的图给出了正合列~\eqref{eq:28} 的几何解释.

传统上, 在实施评注 \ref{thm:4.5.2} 的策略时, 用 $p^n$-覆盖而非上同调类工作, 而 $p^n$-覆盖的构造在文献中通常称为第 $n$ 次下降 (nth descent).\index{xiajiang@下降(descent)!di-n-ci@第 n 次(nth descent)} 见 \cite[\S 26]{cassels1966}.

\begin{aside} \label{thm:4.7.11} 
设 $\varphi\colon E\to E'$ 为椭圆曲线的同态. $E$ 的主齐性空间 $W$ 给出 $E'$ 的良定义主齐性空间 $W' = \varphi_*W$ 与 $\varphi$-等变映射 $W\to W'$. 给出 $n$-覆盖等价于给出 $E$ 的主齐性空间 $W$ 与 $n_*E$ 的平凡化.
\end{aside}

\subsection*{椭圆曲线的扭}

本小节研究如下问题:
\begin{quote}
给定 $k$ 上的椭圆曲线 $E_0$, 求所有在 $k^{\mathrm{al}}$ 上同构于 $E_0$ 的 $k$ 上椭圆曲线 $E$.
\end{quote}
这样的曲线 $E$ 称为 $E_0$ 的扭 (twist).\index{niuquxian@扭曲线(twist)} 回忆 $k$ 上椭圆曲线 $E$ 有指定点 $O\in E(k)$. 为简单, 假设 $k$ 的特征 $\neq 2, 3$, 故可把椭圆曲线写为
\[
E_{(a,b)}\colon Y^2Z = X^3 + aXZ^2 + bZ^3, \qquad a, b\in k, \ \Delta\neq 0.
\]
回忆 (I\ref{thm:1.2.1}) $E_{(a,b)}$ 与 $E_{(a',b')}$ 同构当且仅当存在 $c\in k^{\times}$ 使 $a' = c^4a$, $b' = c^6b$. 每个这样的 $c$ 定义同构
\[
(x:y:z)\mapsto(c^2x:c^3y:z)\colon E_{(a,b)}\to E_{(a',b')},
\]
而所有同构都有此形式.

\begin{example} \label{thm:4.7.12} 
对所有 $d\in k^{\times}$,
\[
E_d\colon dY^2Z = X^3 + aXZ^2 + bZ^3
\]
是 $E_{(a,b)}$ 的扭. 事实上, 作变量替换 $dZ\leftrightarrow Z$ 后, 方程变为
\[
Y^2Z = X^3 + \frac{a}{d^2}XZ^2 + \frac{b}{d^3}Z^3,
\]
故 $E_d$ 在 $d$ 为平方的任何域上与 $E_{(a,b)}$ 同构.
\end{example}

为能应用上同调, 需要计算固定零元素的 $E$ 的自同构群 $\Aut(E, O)$. 这样的自同构与群结构交换 (I\ref{thm:1.1.5}), 故可把 II\ref{thm:2.2.1}b 应用于 $(a', b') = (a, b)$.

\paragraph{} 
\begin{description}
    \item[情形 $ab\neq 0$:] 此时求满足 $c^4 = 1 = c^6$ 的 $c\in k^{\times}$. 这些方程蕴含 $c = \pm 1$, 故 $(E, O)$ 的唯一自同构是恒等映射与
    \[
    (x:y:z)\mapsto(x:-y:z).
    \]
    \item[情形 $a = 0$:] 此时 $c$ 可以是 $k$ 中 $1$ 的任意 $6$ 次根 $\zeta$, 而 $(E, O)$ 的自同构是映射\[
    (x:y:z)\mapsto(\zeta^{2i}x:\zeta^{3i}y:z), \qquad i = 0, 1, 2, 3, 4, 5.
    \]
    \item[情形 $b = 0$:] 此时 $c$ 可以是 $k$ 中 $1$ 的任意 $4$ 次根 $\zeta$, 而 $(E, O)$ 的自同构是映射
    \[
    (x:y:z)\mapsto(\zeta^{2i}x:\zeta^{3i}y:z), \qquad i = 0, 1, 2, 3.
    \]
\end{description}

\begin{proposition} \label{thm:4.7.13} 
设 $(E, O)$ 为特征 $\neq 2, 3$ 的域 $k$ 上的椭圆曲线. $(E, O)$ 的自同构群为 $\{\pm 1\}$, 除非 $j(E) = 0$ 或 $1728$, 此时它分别同构于 $\mu_6(k)$ 或 $\mu_4(k)$.
\end{proposition}

\begin{proof}
因 $j(E) = \dfrac{1728\cdot 4a^3}{4a^3 + 27b^2}$, 这由上面的讨论推出.
\end{proof}

对 $k$ 上的椭圆曲线 $E_0$, 记 $\Aut_{k^{\mathrm{al}}}(E) = \Aut(E_{k^{\mathrm{al}}})$.

\begin{remark} \label{thm:4.7.14} 
\begin{enumerate}[label=\textup{(\alph*)}]
\item 注意命题与 III\ref{thm:3.3.17} 一致, 后者说 (在 $\mathbb{C}$ 上) $\operatorname{End}(E)$ 同构于 $\mathbb{Z}$ 或域 $\mathbb{Q}[\sqrt{-d}]$ ($d > 0$) 的整数环的子环. 这样的环中唯一单位是单位根, 且只有 $\mathbb{Q}[\sqrt{-1}]$ 与 $\mathbb{Q}[\sqrt{-3}]$ 含除 $\pm 1$ 外的单位根.
\item 若允许 $k$ 有特征 $2$ 或 $3$, 则当 $j(E)\neq 0, 1728$ 时仍有 $\Aut(E, O) = \{\pm 1\}$.
然而, 若 $k$ 特征 $2$ 代数闭且 $j = 0 = 1728$, 则 $E$ 同构于曲线 $Y^2 - Y = X^3$. 此时 $\Aut(E, O)$ 是 $24$ 阶非交换群. 它在 $E_3(k)\cong\mathbb{F}_3^2$ 上的作用定义同构 $\Aut(E, O)\cong\SL_2(\mathbb{F}_3)$.
若 $k$ 特征 $3$ 代数闭且 $j = 0 = 1728$, 则 $E$ 同构于曲线 $Y^2 = X^3 - X$. 此时 $\Aut(E, O)$ 是 $12$ 阶非交换群. 它在 $E_2(k)\cong\mathbb{F}_2^2$ 上的作用定义同构 $\Aut(E, O)/\{\pm 1\}\cong\SL_2(\mathbb{F}_2)\cong S_3$.
\end{enumerate}
\end{remark}

固定 $k$ 上的椭圆曲线 $E_0$, 设 $E$ 为在 $k^{\mathrm{al}}$ 上同构于 $E_0$ 的 $k$ 上椭圆曲线. 选 $k^{\mathrm{al}}$ 上的同构 $\varphi\colon E_0\to E$. 对每个 $\sigma\in\Gamma$, 得到第二个同构 $\sigma\varphi \overset{\text{def}}{=} \sigma\circ\varphi\circ\sigma^{-1}\colon E_0\to E$ (在 $k^{\mathrm{al}}$ 上). 例如, 若 $\varphi$ 是 $(x:y:z)\mapsto(c^2x:c^3y:z)$, 则 $\sigma\varphi$ 是 $(x:y:z)\mapsto((\sigma c)^2x:(\sigma c)^3y:z)$. 两个同构 $\varphi, \sigma\varphi\colon E_0\to E$ (在 $k^{\mathrm{al}}$ 上) 相差 $E_0$ 在 $k^{\mathrm{al}}$ 上的自同构:
\[
\sigma\varphi = \varphi\circ\alpha(\sigma), \qquad \alpha(\sigma)\in\Aut_{k^{\mathrm{al}}}(E_0, O).
\]
比较等式
\begin{align*}
(\sigma\tau)\varphi &= \varphi\circ\alpha(\sigma\tau) && \qquad \text{(由定义)}\\
(\sigma\tau)\varphi &= \sigma(\tau\varphi) = \sigma(\varphi\circ\alpha(\tau)) = \varphi\circ\alpha(\sigma)\circ\sigma\alpha(\tau),
\end{align*}
我们看到 $\alpha$ 是交叉同态 $\Gamma\to\Aut_{k^{\mathrm{al}}}(E_0, O)$. 选不同的同构 $\varphi$ 把 $\alpha(\sigma)$ 换成它与主交叉同态的复合.

\begin{theorem} \label{thm:4.7.15} 
映射 $E\mapsto[\alpha]$ 定义一一对应
\[
\{E_0 \text{ 的扭}\}\xrightarrow{\ \sim\ }H^1(\Gamma, \Aut_{k^{\mathrm{al}}}(E_0)).
\]
\end{theorem}

\begin{proof}
证明与定理 \ref{thm:4.7.8} 类似.
\end{proof}

\begin{corollary} \label{thm:4.7.16} 
若 $j(E_0)\neq 0, 1728$, 则 $E_0$ 的扭恰是例 \ref{thm:4.7.12} 中的曲线 $E_d$.
\end{corollary}

\begin{proof}
此时 $\Aut_{k^{\mathrm{al}}}(E, O) = \mu_2$, 故 $H^1(\Gamma, \mu_2)\cong k^{\times}/k^{\times2}$ (见~\ref{thm:4.1.8}). 在定理的对应下, $E_d\leftrightarrow d\bmod k^{\times2}$.
\end{proof}

\begin{remark} \label{thm:4.7.17} 
同样的论证可用于得到 第~\pageref{subsec:twisted-multiplicative}~页中扭乘性群的描述. 代数群 $\mathbb{G}_m = \mathbb{A}^1\setminus\{0\}$ 的自同态是映射 $t\mapsto t^m$, $m\in\mathbb{Z}$. 故 $\operatorname{End}(\mathbb{G}_m)\cong\mathbb{Z}$, $\Aut(\mathbb{G}_m) = (\operatorname{End}(\mathbb{G}_m))^{\times} = \{\pm 1\}$. $\mathbb{G}_m$ 的扭形式由
\[
H^1(k, \{\pm 1\})\cong H^1(k, \mu_2)\cong k^{\times}/k^{\times2}
\]
分类. 对应于 $a\in k^{\times}/k^{\times2}$ 的扭乘性群是 $\mathbb{G}_m[a]$.
\end{remark}

\begin{remark} \label{thm:4.7.18} 
设 $\Aut(E)$ 为 $E$ 的所有自同构 (不必保持 $O$) 的群. 映射 $Q\mapsto t_Q$ (其中 $t_Q$ 是平移 $P\mapsto P + Q$) 把 $E(k)$ 等同于 $\Aut(E)$ 的子群. 我断言 $\Aut(E)$ 是半直积:
\[
\Aut(E) = E(k)\rtimes\Aut(E, O),
\]
即:
\begin{enumerate}[label=\textup{(\alph*)}]
\item $E(k)$ 是 $\Aut(E)$ 的正规子群;
\item $E(k)\cap\Aut(E, O) = \{0\}$;
\item $\Aut(E) = E(k)\cdot\Aut(E, O)$.
\end{enumerate}
设 $Q\in E(k)$, $\alpha\in\Aut(E, O)$. 则 $\alpha$ 是同态, 故对每个 $P\in E$,
\[
(\alpha\circ t_Q\circ\alpha^{-1})(P) = \alpha(\alpha^{-1}(P) + Q) = P + \alpha(Q) = t_{\alpha(Q)}(P),
\]
这证明了 (a). 断言 (b) 显然. 对 (c), 设 $\psi\in\Aut(E)$, 且 $\psi(0) = Q$; 则 $\psi = t_Q\circ(t_{-Q}\circ\psi)$, 而 $t_{-Q}\circ\psi\in\Aut(E, O)$.
\end{remark}

\subsection*{亏格 $1$ 的曲线}

设 $W$ 为 $k$ 上椭圆曲线 $E$ 的主齐性空间. 则 $W$ 在 $k^{\mathrm{al}}$ 上同构于 $E$, 故 $W$ 射影、非奇异、亏格 $1$ (至少在 $k^{\mathrm{al}}$ 上, 这蕴含它在 $k$ 上也). 下面的定理表明, 反之, $k$ 上每个射影非奇异亏格 $1$ 曲线 $W$ 都作为某个 $k$ 上椭圆曲线的主齐性空间出现.

\begin{theorem} \label{thm:4.7.19} 
设 $W$ 为 $k$ 上亏格 $1$ 的非奇异射影曲线. 存在 $k$ 上椭圆曲线 $E_0$ 使 $W$ 是 $E_0$ 的主齐性空间. 且 $E_0$ 在 (在 $k$ 上) 同构的意义下唯一.
\end{theorem}

\begin{proof}
选 $W(k^{\mathrm{al}})$ 中的点, 用它定义同构
\[
W_{k^{\mathrm{al}}}\xrightarrow{\varphi}E\subset\mathbb{P}^2
\]
从 $W_{k^{\mathrm{al}}}$ 到椭圆曲线
\[
Y^2Z + a_1XYZ + a_3YZ^2 = X^3 + a_2X^2Z + a_4XZ^2 + a_6Z^3, \qquad a_i\in k^{\mathrm{al}}
\]
(见II, \ref{thm:2.1.1} 中 (d)$\implies$(c) 的证明). 设 $\sigma\in\Gamma$. 则 $\sigma\varphi$ 是同构 $\sigma(W_{k^{\mathrm{al}}})\to\sigma E$. 这里 $\sigma(W_{k^{\mathrm{al}}})$ 与 $\sigma E$ 是分别把 $\sigma$ 应用于定义 $W_{k^{\mathrm{al}}}$ 与 $E$ 的多项式系数得到的. 例如 $\sigma E$ 由与 $E$ 相同但把 $a_i$ 换成 $\sigma a_i$ 的方程定义. 因 $W$ 由系数在 $k$ 中的多项式定义, $\sigma(W_{k^{\mathrm{al}}}) = W_{k^{\mathrm{al}}}$. 故 $E\cong W_{k^{\mathrm{al}}}\cong\sigma E$, 从而 $j(E) = j(\sigma E) = \sigma j(E)$. 因这对所有 $\sigma\in\Gamma$ 成立, $j(E)\in k$. 由 II, \ref{thm:2.2.3}, 存在 $j(E_0) = j(E)$ 的 $k$ 上曲线 $E_0$. 事实上, $k$ 上会有许多这样的曲线, 故必须确保选对. 我们选一个 $E_0$, 再把它扭成正确的那个.

选 $k^{\mathrm{al}}$ 上的同构 $\varphi\colon E_0\to W$, 对 $\sigma\in\Gamma$, 令 $\sigma\varphi = \varphi\circ\alpha(\sigma)$, 其中 $\alpha(\sigma)\in\Aut_{k^{\mathrm{al}}}(E_0)$. 则 $\sigma\mapsto\alpha(\sigma)$ 是到 $\Aut_{k^{\mathrm{al}}}(E_0)$ 中的交叉同态, 从而定义 $H^1(k, \Aut_{k^{\mathrm{al}}}(E_0))$ 中的类 $[\alpha]$. 由~\ref{thm:4.7.18}, 有正合序列
\[
1\to E_0(k^{\mathrm{al}})\to\Aut_{k^{\mathrm{al}}}(E_0)\to\Aut_{k^{\mathrm{al}}}(E_0, O)\to 1.
\]
若 $[\alpha]$ 位于 $H^1(k, \Aut_{k^{\mathrm{al}}}(E_0))$ 的子群 $H^1(k, E_0)$ 中, 则 $W$ 是 $E_0$ 的主齐性空间. 若否, 用 $[\alpha]$ 在 $H^1(k, \Aut_{k^{\mathrm{al}}}(E_0, O))$ 中的像扭 $E_0$, 得到 $k$ 上第二个具有相同 $j$-不变量的曲线 $E_1$. 现在可验证交叉同态 $[\alpha]$ 的类位于 $H^1(k, E_1)$ 中, 故 $W$ 是 $E_1$ 的主齐性空间.
\end{proof}

定理给出的曲线 $E_0$ 称为 $W$ 的 Jacobi 簇 (Jacobian vartiety). 它由如下性质刻画: 存在 $k^{\mathrm{al}}$ 上的同构 $\varphi\colon E_0\to W$ 使对每个 $\sigma\in\Gamma$, 存在点 $Q_\sigma\in E_0(k^{\mathrm{al}})$ 使
\[
(\sigma\varphi)(P) = \varphi(P + Q_\sigma), \qquad \forall P\in E(k^{\mathrm{al}}).
\]
为构造 $\mathbb{Q}$ 上曲线 $W$ 的 Jacobi 簇, 用以下步骤:
\begin{enumerate}[label=\textup{(\alph*)}]
\item 经 $\mathbb{Q}^{\mathrm{al}}$ 上的变量替换, 得同构 $W\cong E$, 其中 $E$ 是 $\mathbb{Q}^{\mathrm{al}}$ 上标准形式的椭圆曲线;
\item 写下 $\mathbb{Q}$ 上标准形式、在 $\mathbb{Q}^{\mathrm{al}}$ 上同构于 $E$ 的椭圆曲线 $E_0$;
\item 必要时修改 $E_0$, 使它具有刻画 Jacobi 簇的性质.
\end{enumerate}

\begin{remark} \label{thm:4.7.20} 
上面证明中用了取值于不必交换的群 $\Aut_{k^{\mathrm{al}}}(E_0)$ 的交叉同态. 但 $M$ 不交换时仍可定义 $H^1(G, M)$. 用乘法写 $M$. 如交换情形, 交叉同态是满足 $f(\sigma\tau) = f(\sigma)\cdot\sigma f(\tau)$ 的映射 $f\colon G\to M$. 若存在 $m\in M$ 使对所有 $\sigma\in G$ 有 $g(\sigma) = m^{-1}\cdot f(\sigma)\cdot\sigma m$, 则称两个交叉同态 $f$ 与 $g$ 等价, 而 $H^1(G, M)$ 定义为交叉同态的等价类集. 它是带指定元素 (常数映射 $\sigma\mapsto 1$ 的类) 的集合.
\end{remark}

\begin{remarkstar}
设 $C$ 为 $k$ 上亏格 $1$ 的曲线, $D$ 为 $C$ 上次数 $n+1\ge 2$ 的除子. 由 Riemann--Roch 定理, $L(D)$ 的维数为 $n+1$. 基 $f_0, f_1, \dots, f_n$ 定义态射 $P\mapsto(f_0(P):\dots:f_n(P))\colon C\to\mathbb{P}^n$. 这提出两个问题: 如何找到描述 $C$ 在 $\mathbb{P}^n$ 中像的方程; 以及由这些方程, 如何找到 $C$ 的 Jacobi 簇的 Weierstrass 方程. 关于这些问题的最新文章, 见 \cite{fisher2018}.
\end{remarkstar}

\subsection*{$\mathbb{Q}$ 上亏格 $1$ 曲线的分类}

总结上述结果.

\begin{aside} \label{thm:4.7.21} 
设 $(E, O)$ 为 $\mathbb{Q}$ 上的椭圆曲线. 给它附上不变量 $j(E)\in\mathbb{Q}$. $\mathbb{Q}$ 中每个元素都作为 $\mathbb{Q}$ 上椭圆曲线的 $j$-不变量出现, 且 $\mathbb{Q}$ 上两条椭圆曲线有相同的 $j$-不变量当且仅当它们在 $\mathbb{Q}^{\mathrm{al}}$ 上同构 (II\ref{thm:2.2.1}).
\end{aside}

\begin{aside} \label{thm:4.7.22} 
固定 $j\in\mathbb{Q}$, 考虑 $j(E) = j$ 的 $\mathbb{Q}$ 上椭圆曲线 $(E, O)$. 选一条, 例如 II\ref{thm:2.2.3} 中那条. 这样的曲线的同构类与 $H^1(\mathbb{Q}, \Aut(E, O))$ 的元素自然一一对应. 例如, 当 $j\neq 0, 1728$ 时, $\Aut(E, O) = \mu_2$, $H^1(\mathbb{Q}, \Aut(E, O))\cong\mathbb{Q}^{\times}/\mathbb{Q}^{\times2}$, 而对应于 $d\in\mathbb{Q}^{\times}$ 的曲线是例 \ref{thm:4.7.12} 中的曲线 $E_d$.
\end{aside}

\begin{aside} \label{thm:4.7.23} 
固定 $\mathbb{Q}$ 上的椭圆曲线 $(E, O)$, 考虑以 $E$ 为 Jacobi 簇的 $\mathbb{Q}$ 上亏格 $1$ 曲线. 这样的曲线有 $E$ 的主齐性空间的结构, 而每个 $E$ 的主齐性空间以 $E$ 为 Jacobi 簇. $E$ 的主齐性空间由群 $H^1(\mathbb{Q}, E)$ (非常大的群) 分类.
\end{aside}

\begin{aside} \label{thm:4.7.24} 
$\mathbb{Q}$ 上每条亏格 $1$ 曲线都有 Jacobi 簇, 它是 $\mathbb{Q}$ 上的椭圆曲线 $E$, 且该曲线是 $E$ 的主齐性空间.
\end{aside}

\begin{exercise} \label{thm:4.7.25} 
求曲线
\[
W\colon aX^3 + bY^3 + cZ^3 = 0, \qquad a, b, c\in\mathbb{Q}^{\times}
\]
的 Jacobi 簇. 提示: 曲线 $E\colon X^3 + Y^3 + dZ^3 = 0$ ($d\in\mathbb{Q}^{\times}$) 有基点 $O\colon(1:-1:0)$——序对 $(E, O)$ 是 $\mathbb{Q}$ 上的椭圆曲线. 它可通过变量替换 $X = X' + Y'$, $Y = X' - Y'$ 化为标准形式.
\end{exercise}

\section{Hasse (局部--整体) 原理的失败}\label{sec:4.8}

设 $E$ 为 $\mathbb{Q}$ 上的椭圆曲线. 回忆其 Tate--Shafarevich 群是 $H^1(\mathbb{Q}, E)\to\prod_p H^1(\mathbb{Q}_p, E)$ 的核. 因此, $E$ 的主齐性空间 $W$ 的上同调类位于 Tate--Shafarevich 群中当且仅当对所有 $p$ (包括 $p = \infty$) 有 $W(\mathbb{Q}_p)\neq\varnothing$. 故 $\Sha(E/\mathbb{Q})$ 的非零元素由以 $E$ 为 Jacobi 簇、使 Hasse 原理失败的 $\mathbb{Q}$ 上曲线表示.

本节描述一族 Tate--Shafarevich 群非零的曲线, 从而给出亏格 $1$ 曲线 Hasse 原理失败的例子.

\begin{proposition} \label{thm:4.8.1} 
若 $p\equiv 1\pmod{8}$, 则椭圆曲线
\[
E\colon Y^2Z = X^3 + pXZ^2
\]
的 $2$-Selmer 群 $S^{(2)}(E/\mathbb{Q})$ 同构于 $(\mathbb{Z}/2\mathbb{Z})^3$.
\end{proposition}

陈述中的曲线族与习题 \ref{thm:4.5.9}\hyperref[item:4.5.9.2]{\ref*{item:4.5.9.2}} 中的类似, 但因 $E$ 上只有 $2$ 阶点之一坐标在 $\mathbb{Q}$ 中, 我们没有 $H^1(\mathbb{Q}, E_2)$ 的简单描述. 当然, 可以过渡到 $\mathbb{Q}[\sqrt{p}]$, 但更简单是如下进行. 设 $E_0$ 为 $E$ 对由它的 $2$ 阶点 $P = (0:0:1)$ 生成的子群的商. 有同态
\[
E\xrightarrow{\pi}E_0\xrightarrow{\varphi}E
\]
其复合是乘以 $2$. 研究正合列
\begin{gather*}
0\to\langle P\rangle\to E(\mathbb{Q}^{\mathrm{al}})\xrightarrow{\pi}E_0(\mathbb{Q}^{\mathrm{al}})\to 0,\\
0\to\operatorname{Ker}\to E_0(\mathbb{Q}^{\mathrm{al}})\xrightarrow{\varphi}E(\mathbb{Q}^{\mathrm{al}})\to 0
\end{gather*}
{\emergencystretch=3em 的上同调序列, 可以得到关于 $E(\mathbb{Q})/2E(\mathbb{Q})$, $S^{(2)}(E/\mathbb{Q})$ 与 $\Sha(E/\mathbb{Q})_2$ 的信息. 例如, 映射对 $E(\mathbb{Q})\xrightarrow{\pi}E_0(\mathbb{Q})\xrightarrow{\varphi}E(\mathbb{Q})$ 的核--余核序列是正合序列\par}
\[
E(\mathbb{Q})/\pi(E(\mathbb{Q}))\to E(\mathbb{Q})/2E(\mathbb{Q})\to E(\mathbb{Q})/\varphi(E(\mathbb{Q}))\to 0.
\]
见 \cite[X, 6.2]{silverman2009}.

假设命题成立. 因 $E(\mathbb{Q})_2\cong\mathbb{Z}/2\mathbb{Z}$, 由基本正合列~\eqref{eq:29},
\[
\operatorname{rank}(E(\mathbb{Q})) + 1 + \dim_{\mathbb{F}_2}\Sha(E/\mathbb{Q})_2 = \dim_{\mathbb{F}_2}S^{(2)}(E/\mathbb{Q}).
\]
故 $E$ 的秩 $r$ 为 $0, 1$ 或 $2$. 若 $\Sha(E/\mathbb{Q})(2)$ 有限, 其阶已知是平方 (见~\ref{thm:4.5.4}), 故按猜想, 唯一可能是 $r = 0, 2$.

\begin{proposition} \label{thm:4.8.2} 
设 $E$ 如命题 \ref{thm:4.8.1} 中. 若 $2$ 模 $p$ 不是四次幂, 则 $\operatorname{rank}(E(\mathbb{Q})) = 0$ 且 $\Sha(E/\mathbb{Q})_2\cong(\mathbb{Z}/2\mathbb{Z})^2$.
\end{proposition}

我们在下面讨论证明.

\begin{remark} \label{thm:4.8.3} 
为应用命题, 必须找使 $2$ 模 $p$ 为四次幂的素数 $p$. Gauss 找到\index{Gauss}有效检验, 现在解释. 由基础代数, Gauss 整数环 $\mathbb{Z}[i]$ 是主理想整环. 奇素数 $p$ 要么在 $\mathbb{Z}[i]$ 中保持素数, 要么分解为 $p = (A+iB)(A-iB)$. 第一种情形恰在 $\mathbb{Z}[i]/p\mathbb{Z}[i]$ 是 $\mathbb{F}_p$ 上 $2$ 次域时发生. 故 $p$ 保持素数当且仅当 $\mathbb{F}_p$ 不含 $1$ 的 $4$ 次本原根. 因 $\mathbb{F}_p^{\times}$ 循环, 它含 $4$ 阶元素当且仅当 $4$ 整除其阶 $p-1$. 我们得出结论: 可表为 $p = A^2 + B^2$ ($A, B\in\mathbb{Z}$) 的奇素数 $p$ 恰是满足 $p\equiv 1\pmod{4}$ 的那些. 对满足 $p\equiv 1\pmod{8}$ 的素数 $p$, Gauss 表明 $2$ 是模 $p$ 的四次幂当且仅当 $8\mid AB$. 因此, 命题适用于素数
\[
17 = 1^2 + 4^2, \qquad 41 = 5^2 + 4^2, \qquad 97 = 9^2 + 4^2, \qquad 193 = 7^2 + 12^2, \dots
\]
Gauss 判据的证明相当初等, 见 \cite[X, 6.6]{silverman2009}. 数论学者会希望证明这样的素数 $p$ 有无限多个 (并求其密度).

为证明椭圆曲线的秩小于 Selmer 群给出的界, 必须展示 $\Sha(E/\mathbb{Q})_2$ 的非平凡元素. 对命题中曲线 $E\colon Y^2Z = X^3 + pZ^3$, 有以下曲线:
\[
Y^2 = 4pX^4 - 1, \qquad Y^2 = 2pX^4 - 2, \qquad -Y^2 = 2pX^4 - 2.
\]
可直接验证它们是 $E\colon Y^2Z = X^3 + pZ^3$ 的主齐性空间, 但从命题 \ref{thm:4.8.1} 的证明更容易推出 (\cite[X, 6.2(b)]{silverman2009}).
\end{remark}

\begin{remark} \label{thm:4.8.4} 
应解释上面列出的曲线的含义. 更一般地, 考虑曲线
\[
C\colon Y^2 = aX^4 + bX^3 + cX^2 + dX + e.
\]
假设右边的多项式无重根, 且特征 $\neq 2, 3$. 则这是非奇异仿射曲线, 但它的射影闭包
\[
\bar{C}\colon Y^2Z^2 = aX^4 + bX^3Z + cX^2Z^2 + dXZ^3 + eZ^4
\]
奇异, 因为令 $Y = 1$, 得方程
\[
Z^2 = aX^4 + bX^3Z + cX^2Z^2 + dXZ^3 + eZ^4,
\]
它在 $(0,0)$ 处显然奇异. $d$ 次平面射影曲线的亏格是
\[
g = \frac{(d-1)(d-2)}{2} - \sum_{P \text{ 奇异}}\delta_P.
\]
对 $P = (0,0)$, $\delta_P = 2$, 故 $\bar{C}$ 的亏格是 $3 - 2 = 1$. 可消解奇点得到非奇异曲线 $C'$ 与正则映射 $C'\to\bar{C}$, 它在奇异点之外是同构. 写 $C$ 时真正指的是 $C'$.
\end{remark}

我们将证明曲线
\[
C\colon Y^2 = 2 - 2pX^4
\]
在 $\mathbb{Q}$ 中没有点, 但在 $\mathbb{R}$ 与对所有 $p$ 的 $\mathbb{Q}_p$ 中有点. 为此需要二次互反律. 对不被素数 $p$ 整除的整数 $a$, Legendre 符号\index{Legendre-fuhao@Legendre 符号(Legendre symbol)}\index{Legendre} $\bigl(\frac{a}{p}\bigr)$ 在 $a$ 是模 $p$ 的平方 (即二次剩余) 时为 $+1$, 否则为 $-1$.

\begin{theorem}[二次互反律] \label{thm:4.8.5} 
对奇素数 $p, q$,
\[
\Bigl(\frac{p}{q}\Bigr)\Bigl(\frac{q}{p}\Bigr) = (-1)^{\frac{p-1}{2}\frac{q-1}{2}}.
\]
且
\[
\Bigl(\frac{2}{p}\Bigr) = (-1)^{\frac{p^2-1}{8}}.
\]
\end{theorem}

\begin{proof}
该定理的已发表证明可能比数学中任何其他定理都多.\footnote{与之竞争的有勾股定理与平行公设 (其全部证明都是错的).} 前六个完整证明由 Gauss 找到. 两个初等证明见 \cite[Chapter I]{serre1973}. 用分圆域的短证明见 ANT, Chap. 8, Examples, 或 \cite[\S 6.5]{samuel2008}.
\end{proof}

现在证明 $C\colon Y^2 = 2 - 2pX^4$ 没有坐标在 $\mathbb{Q}$ 中的点. 设 $(x, y)$ 为曲线上的点. 令 $x = r/t$, $r$ 与 $t$ 为无公因子的整数. 则
\[
y^2 = \frac{2t^4 - 2pr^4}{t^4}.
\]
右边的分子与分母是整数且无公因子, 故 $2t^4 - 2pr^4$ 是整数的平方, 它必为偶. 故存在整数 $s$ 使
\[
2s^2 = t^4 - pr^4.
\]
设 $q$ 为整除 $s$ 的奇素数. 则 $t^4\equiv pr^4\pmod{q}$, 故 $\bigl(\frac{p}{q}\bigr) = 1$. 由二次互反律, 这蕴含 $\bigl(\frac{q}{p}\bigr) = 1$. 由二次互反律, $\bigl(\frac{2}{p}\bigr) = 1$, 故 $s$ 的所有素因子都是模 $p$ 的平方. 于是 $s^2$ 是模 $p$ 的四次幂. 方程
\[
2s^2\equiv t^4 \pmod{p}
\]
现在表明 $2$ 是模 $p$ 的四次幂, 与假设矛盾.

还应确保无穷远处没有潜伏的点. $C$ 的射影闭包是
\[
\bar{C}\colon Y^2Z^2 = 2Z^4 - 2pX^4,
\]
而刚刚证明了 $\bar{C}$ 没有 $Z = 1$ 的有理点. 对 $Z = 0$, 有有理解 $(0:1:0)$, 但它是曲线 $Z^2 = 2Z^4 - 2pX^4$ 上的奇异点 $(0,0)$. 在 $\bar{C}$ 的消解 $C'\to\bar{C}$ 上, 没有 $\mathbb{Q}$ 点位于 $(0:1:0)$ 之上.

曲线 $C$ 显然在 $\mathbb{R}$ 中有点. Hensel 引理 (I\ref{thm:1.2.13}) 表明若 $C$ 模素数 $q$ 的约化在 $\mathbb{F}_q$ 中有非奇异点, 则 $C$ 在 $\mathbb{Q}_q$ 中有点. 对 $q\neq 2, p$, 曲线 $C$ 在 $q$ 处有好约化, 而下面的推论 \ref{thm:4.9.3} 表明它在 $\mathbb{F}_p$ 中有点. 故 $C$ 自动在 $\mathbb{Q}_q$ 中有点, 除非 $q$ 等于 $2$ 或 $p$. 这两个素数需要对 Hensel 引理更精细的应用, 留作读者练习.

\begin{remarkstar}
Hasse 原理失败的第一个亏格 $1$ 曲线例子是
\[
X^4 - 17 = 2Y^2
\]
(\cite{lind1940}). 这里, 跟随 \cite{cassels1966}, 是 Lind 的论证, 证明它没有 $\mathbb{Q}$-有理点. 设 $(u, v, w)$ 为
\[
u^4 - 17v^2 = 2w^2
\]
在互素整数中的解. 若 $p$ 是整除 $w$ 的奇素数, 则 $u^4 - 17v^2\equiv 0\pmod{p}$, 故 $17$ 是 $p$ 的二次剩余. 由二次互反律, 这蕴含 $p$ 是 $17$ 的二次剩余. 因 $2$ 是 $17$ 的二次剩余 (再由二次互反律), 我们看到 $w$ 的所有素因子都是 $17$ 的二次剩余, 故 $w$ 模 $17$ 是平方. 故 $w^2$ 模 $17$ 是四次幂, 而方程蕴含 $2$ 模 $17$ 是四次幂, 这不成立.
\end{remarkstar}

\section{有限域上的椭圆曲线}\label{sec:4.9}

固定 $\mathbb{F}_p$ 的代数闭包 $\bar{\mathbb{F}}$, 令 $\Gamma = \Gal(\bar{\mathbb{F}}/\mathbb{F}_p)$. 映射 $a\mapsto a^p$ 是 $\Gamma$ 的拓扑生成元, 称为 Frobenius 自同构\index{Frobenius@Frobenius(Frobenius)!zitonggou@自同构(automorphism)} (\cite[7.16]{milne-ft}).

记 $\mathbb{F}_{p^m}$ 为元素是 $X^{p^m} - X$ 的根的子域. 故 $\mathbb{F}_{p^m}$ 有 $p^m$ 个元素, 且 $\mathbb{F}_{p^m}\subset\mathbb{F}_{p^n}$ 当且仅当 $m\mid n$. $X^{p^m} - X$ 在 $\mathbb{F}_p[X]$ 中的不可约因子恰是 $\mathbb{F}_p[X]$ 中次数整除 $m$ 的首一不可约多项式 (\cite[4.22, \ref{thm:1.4.23}]{milne-ft}).

我们将利用微积分中许多幂级数恒等式在 $\mathbb{Q}$ 上有效的事实. 例如, 若在幂级数环 $\mathbb{Q}[[T]]$ 中定义
\begin{align*}
    \log(1+T) & = T - \frac{1}{2}T^2 + \frac{1}{3}T^3 - \frac{1}{4}T^4 + \frac{1}{5}T^5 - \cdots,\\
    \exp(T) & = 1 + T + \frac{1}{2!}T^2 + \cdots + \frac{1}{n!}T^n + \cdots,
\end{align*}
则
\begin{align*}
    \log\frac{1}{1-T} & = -\log(1-T) = T + \frac{1}{2}T^2 + \frac{1}{3}T^3 + \frac{1}{4}T^4 + \frac{1}{5}T^5 + \cdots,\\
    \exp(\log(1+T)) & = 1 + T.
\end{align*}

\subsection*{数 $a_p$}

设 $E_{(a,b)}\colon Y^2Z = X^3 + aXZ^2 + bZ^3$ 为 $\mathbb{F}_p$ ($p\neq 2$) 上的椭圆曲线, $N_p$ 为 $E$ 上坐标在 $\mathbb{F}_p$ 中的点数. 有一个``无穷远''点 $(0:1:0)$, 其余是 $Y^2 = X^3 + aX + b$ 在 $\mathbb{F}_p$ 中的解. 计数它们的一种方法是列出 $\mathbb{F}_p$ 中的平方, 再把 $x^3 + ax + b$ 的值 ($x\in\mathbb{F}_p$) 与列表对照. 对 $z\in\mathbb{F}_p^{\times}$, 若 $z$ 是 $\mathbb{F}_p^{\times}$ 中的平方则令 $\chi(z) = 1$, 否则 $-1$; 把 $\chi$ 延拓到 $\mathbb{F}_p$ 上令 $\chi(0) = 0$. 则
\[
N_p = 1 + \sum_{x\in\mathbb{F}_p}(\chi(x^3 + ax + b) + 1) = p + 1 - a_p,
\]
其中
\begin{equation} \label{eq:37}
-a_p = \sum_{x\in\mathbb{F}_p}\chi(x^3 + ax + b).
\tag{37}
\end{equation}
因 $\mathbb{F}_p^{\times}$ 偶阶循环, 恰有一半元素是平方. 由于没有理由期望 $x^3 + ax + b$ 比非平方更可能是平方, 可能期望 $-a_p$ 是 $p$ 个在 $+1$ 与 $-1$ 之间随机分布项的和. $|a_p|$ 的期望值将约为 $\sqrt{p}$ 的量级. 当固定 $a, b\in\mathbb{Z}$ 并对不同素数 $p$ 计算曲线 $E_{(a,b)}$ 模 $p$ 的 $a_p$ 时, 计算支持这一点. 但对好约化, 我们总是发现
\begin{equation} \label{eq:38}
|a_p| < 2\sqrt{p}.
\tag{38}
\end{equation}
而若 \eqref{eq:37} 中各项真正随机, 则 $0$ 与 $p-1$ 之间每个值都可能. 我们将稍后证明不等式~\eqref{eq:38}, 并在本书其余部分看到看似无害的数 $a_p$ 变成关键角色.

\subsection*{Frobenius 映射}

设 $C$ 为 $\mathbb{F}_p$ 上次数 $d$ 的射影平面曲线, 故 $C$ 由多项式
\[
F(X, Y, Z) = \sum_{i+j+k=d}a_{ijk}X^iY^jZ^k, \qquad a_{ijk}\in\mathbb{F}_p
\]
定义. 若 $P = (x:y:z)\in C(\bar{\mathbb{F}})$, 则 $\sum_{i+j+k=d}a_{ijk}x^iy^jz^k = 0$. 把此方程提升到 $p$ 次幂, 并利用我们处于特征 $p$ 中且对 $a\in\mathbb{F}_p$ 有 $a^p = a$, 得方程
\[
\sum_{i+j+k=d}a_{ijk}x^{ip}y^{jp}z^{kp} = 0,
\]
这说 $(x^p:y^p:z^p)$ 也在 $C$ 上. 故有映射
\[
(x:y:z)\mapsto(x^p:y^p:z^p)\colon C(\bar{\mathbb{F}})\to C(\bar{\mathbb{F}}),
\]
它由多项式定义, 故正则. 称为 Frobenius 映射.\index{Frobenius@Frobenius(Frobenius)!yingshe@映射(map)}

\begin{proposition} \label{thm:4.9.1} 
(a) Frobenius 映射的次数是 $p$.

(b) Frobenius 映射在 $E$ 的每一点的切线上为零.
\end{proposition}

\begin{proof}
(a) 由图
\[\begin{tikzcd}
	{k(x,y)} && {k(x,y)} \\
	{k(X)} && {k(X)}
	\arrow["{(x,y)\mapsto(x^p,y^p)}", from=1-1, to=1-3]
	\arrow["2"', no head, from=1-1, to=2-1]
	\arrow["2", no head, from=1-3, to=2-3]
	\arrow["{X\mapsto X^p}", from=2-1, to=2-3]
\end{tikzcd}\]
看到 $\deg(\operatorname{Frob}) \overset{\text{def}}{=} [k(x,y):k(x^p,y^p)] = [k(X):k(X^p)] = p$.
(b) 考虑由多项式 $f(X, Y)$ 与 $g(X, Y)$ 定义的正则映射 $\alpha\colon\mathbb{A}^2\to\mathbb{A}^2$, $(a, b)\mapsto(f(a, b), g(a, b))$. 映射 $(d\alpha)_P$ 在 $\mathbb{A}^2$ 的点 $P$ 处的切空间上由矩阵$\begin{pmatrix}
\bigl(\frac{\partial f}{\partial X}\bigr)_P & \bigl(\frac{\partial f}{\partial Y}\bigr)_P\\
\bigl(\frac{\partial g}{\partial X}\bigr)_P & \bigl(\frac{\partial g}{\partial Y}\bigr)_P
\end{pmatrix}$
给出. 若 $\alpha$ 把曲线 $C$ 映到曲线 $C'$, 则 $(d\alpha)_P$ 把 $C$ 的切空间映到 $C'$ 的. 对 Frobenius 映射, $(f(X, Y), g(X, Y)) = (X^p, Y^p)$, 故 $\frac{\partial f}{\partial X}$, $\frac{\partial f}{\partial Y}$, $\frac{\partial g}{\partial X}$, $\frac{\partial g}{\partial Y}$ 都恒为零. 故切空间上的映射都为零.
\end{proof}

\subsection*{$\mathbb{F}_p$ 上亏格 $1$ 的曲线}

\begin{proposition} \label{thm:4.9.2} 
设 $E$ 为 $\mathbb{F}_p$ 上的椭圆曲线. 则 $H^1(\mathbb{F}_p, E) = 0$, 故 $E$ 在 $\mathbb{F}_p$ 上的一切主齐性空间平凡.
\end{proposition}

\begin{proof}
为证命题, 必须证每个连续交叉同态 $f\colon\Gamma\to E(\bar{\mathbb{F}})$ 都是主的. 设 $\varphi\colon E\to E$ 为 Frobenius 映射. 它以 Frobenius 自同构 $\sigma\in\Gamma$ 作用在 $E(\bar{\mathbb{F}})$ 上. 正则映射 $\varphi - \mathrm{id}_E\colon E\to E$ 非常数, 故 $P\mapsto\varphi(P) - P\colon E(\bar{\mathbb{F}})\to E(\bar{\mathbb{F}})$ 是满射 (I\ref{thm:1.4.23}). 特别地, 存在 $P\in E(\bar{\mathbb{F}})$ 使 $\varphi(P) - P = f(\sigma)$, 即 $f(\sigma) = \sigma P - P$. 现在
\[
f(\sigma^2) = f(\sigma) + \sigma f(\sigma) = \sigma P - P + \sigma^2P - \sigma P = \sigma^2P - P,
\]
且归纳地, $f(\sigma^n) = \sigma^nP - P$. 故 $f$ 与主交叉同态 $\gamma\mapsto\gamma P - P$ 在一切 $\sigma^n$ 上一致. 因两者都连续, 这蕴含它们在整个 $\Gamma$ 上一致.

陈述的第二部分由第一部分经定理 \ref{thm:4.7.8} 推出.
\end{proof}

\begin{corollary} \label{thm:4.9.3} 
$\mathbb{F}_p$ 上亏格 $1$ 的每条非奇异射影曲线 $C$ 都有坐标在 $\mathbb{F}_p$ 中的点.
\end{corollary}

\begin{proof}
由定理 \ref{thm:4.7.19}, 存在 $\mathbb{F}_p$ 上椭圆曲线 $E$ 使 $C$ 有 $E$ 的主齐性空间的结构. 由命题, $C$ 平凡, 即 $C(\mathbb{F}_p)\neq\varnothing$.
\end{proof}

\subsection*{椭圆曲线 Riemann 假设的证明}

\begin{theorem}[Riemann 假设] \label{thm:4.9.4} 
对 $\mathbb{F}_p$ 上椭圆曲线 $E$, $|\#E(\mathbb{F}_p) - (p+1)|\le 2\sqrt{p}$.
\end{theorem}

\begin{proof}
映射 $\mathrm{id}_E - \varphi\colon E(\bar{\mathbb{F}})\to E(\bar{\mathbb{F}})$ 的核是 $E$ 上满足 $(x^p:y^p:z^p) = (x:y:z)$ 的点 $(x:y:z)$ 的集合. 这些正是有代表元 $x, y, z\in\mathbb{F}_p$ 的点 (见~\ref{thm:4.1.4}). 因 $\varphi$ 在 $O$ 处的切空间上为零, $\mathrm{id}_E - \varphi$ 作为恒等映射作用, 故 $\mathrm{id}_E - \varphi$ 可分 (I\ref{thm:1.4.24}b). 由此 (I\ref{thm:1.4.24}a),
\[
\#E(\mathbb{F}_p) = \deg(\mathrm{id}_E - \varphi).
\]
设 $f(T) = T^2 - sT + t$ 为 $\varphi$ 作为 $E$ 的自同态的特征多项式. 则 $s, t\in\mathbb{Z}$, $4t\ge s^2$, 且对所有 $n\in\mathbb{Z}$ 有 $f(n) = \deg(n - \varphi)$ (II\ref{thm:2.6.2}). 特别地,
\[
t = f(0) = \deg(\varphi) = p, \qquad 1 - s + p = f(1) = \deg(1 - \varphi) = \#E(\mathbb{F}_p),
\]
故 $|\#E(\mathbb{F}_p) - (p+1)| = |s|\le 2\sqrt{p}$.
\end{proof}

这个简短证明之后, 本节其余部分解释为什么此不等式称为 Riemann 假设.

前三节中所有内容在把 $p$ 换成 $q$ 后仍成立.

\begin{aside} \label{thm:4.9.5} 
设 $E$ 为 $\mathbb{Q}$ 上的椭圆曲线. 由定理 \ref{thm:4.9.4}, 对 $E$ 有好约化的素数 $p$, $E$ 模 $p$ 的点数是 $N_p = p + 1 - a_p$, 且 $|a_p|\le 2\sqrt{p}$. 可写 $a_p = 2\sqrt{p}\cos\theta_p$, $0\le\theta_p\le\pi$. 当 $E$ 在 $\mathbb{C}$ 上无复乘法 (即 $\operatorname{End}_{\mathbb{C}}(E) = \mathbb{Z}$) 时, 佐藤 (Mikio Sato\index{Sato})\footnote{据 Sarnak (2018 年演讲), 佐藤做的远多于计算, 可能第一个理解对称张量幂在此语境中的用途.} 于 1963 年 3 月从计算中发现 $\theta_p$ 似乎有密度分布 $\frac{2}{\pi}\sin^2\theta$, 即对 $0\le c < d\le\pi$,
\[
\lim_{N\to\infty}\frac{\#\{p \mid p \text{ 素数}, p\le N, c\le\theta_p\le d\}}{\#\{p \mid p \text{ 素数}, p\le N\}} = \frac{2}{\pi}\int_c^d\sin^2\theta\,d\theta.
\]
1963 年 12 月, Tate 宣布他著名的猜想, 把代数簇上的代数闭链与其 zeta 函数的极点联系起来. 在考察 $\mathbb{Q}$ 上椭圆曲线的幂时, Tate 发现 (启发式地) Sato 分布是唯一使 zeta 函数的极点与其猜想预测一致的对称密度分布. \cite{serre1968} (Chap.~I 附录) 给出 Tate 论证的严格处理, 并把 $\theta_p$ 按 $\frac{2}{\pi}\sin^2\theta$ 分布这一陈述称为 Sato--Tate 猜想. Sato 没有继续在数论中工作, 但 Sato 与 Tate 的名字在四十年后再次联系在一起, 因为他们都被授予 2002/3 年 Wolf 奖.

Sato--Tate 猜想现已被 Michael Harris, Richard Taylor 及其合作者证明. 见 Barnet-Lamb et al.\ 2011, 其中陈述: ``顺带我们注意到, Sato--Tate 猜想现在可对全实域上任意椭圆曲线证明.''
\end{aside}

\subsection*{数域的 zeta 函数}

首先回忆最初的 Riemann zeta 函数是
\begin{equation} \label{eq:39}
\zeta(s) = \prod_{p\text{ 素数}}\frac{1}{1-p^{-s}} = \sum_{n\ge 1}n^{-s}, \qquad s\in\mathbb{C}, \ \Re(s) > 1.
\tag{39}
\end{equation}
第二个表达式由把第一个展开 \[\prod_p\frac{1}{1-p^{-s}} = \prod_p(1 + p^{-s} + (p^{-s})^2 + \cdots)\] 得到各项之和 $(p_1^{r_1}\cdots p_t^{r_t})^{-s}$ 并应用唯一分解而得到.

\eqref{eq:39} 中的和与积都在 $\Re(s) > 1$ 收敛, 故 $\zeta(s)$ 在 $\Re(s) > 1$ 全纯且非零. 事实上, $\zeta(s)$ 延拓为整个复平面上的亚纯函数, 在 $s = 1$ 处有单极点. 且函数 $\Lambda(s) = \pi^{-s/2}\Gamma(\frac{s}{2})\zeta(s)$ 满足函数方程 $\Lambda(s) = \Lambda(1-s)$, 在 $s = 0, 1$ 处有单极点, 此外全纯. 这里 $\Gamma(s)$ 是 Gamma 函数. 因 $\Gamma(s)$ 在 $s = 0, -1, -2, -3, \dots$ 处有极点, 这迫使 $\zeta$ 在 $s = -2n$ ($n > 0$, $n\in\mathbb{Z}$) 处为零. 这些称为 zeta 函数的平凡零点 (trivial zeros).

\begin{conjecture}[Riemann 假设] \label{thm:4.9.6} 
$\zeta(s)$ 的非平凡零点位于直线 $\Re(s) = \frac{1}{2}$ 上.
\end{conjecture}

这可能是数学中最著名的未解决问题.

Dedekind 把 Riemann 的定义推广到每个数域 $K$, 为之附上 zeta 函数 $\zeta_K(s)$. 他定义
\[
\zeta_K(s) = \prod_{\mathfrak{p}}\frac{1}{1 - \Nm(\mathfrak{p})^{-s}} = \sum_{\mathfrak{a}}\Nm(\mathfrak{a})^{-s}, \qquad s\in\mathbb{C}, \ \Re(s) > 1,
\]
其中 $\mathfrak{p}$ 遍历 $\mathcal{O}_K$ 中非零素理想, $\mathfrak{a}$ 遍历 $\mathcal{O}_K$ 中非零理想. 回忆理想 $\mathfrak{a}$ 的数值范数 $\Nm\mathfrak{a}$ 是商环 $\mathcal{O}_K/\mathfrak{a}$ 的阶. 第二个表达式由第一个经唯一分解 $\mathfrak{a} = \prod\mathfrak{p}_i^{r_i}$ (见~\ref{thm:4.3.9}) 且 $\Nm(\mathfrak{a}) = \prod \Nm(\mathfrak{p}_i)^{r_i}$ 得到. 注意对 $K = \mathbb{Q}$, 此定义回到 $\zeta(s)$. 函数 $\zeta_K(s)$ 延拓为整个复平面上带 $s = 1$ 处单极点的亚纯函数, 且它的某个倍数 $\Lambda_K(s)$ 满足函数方程 $\Lambda_K(s) = \Lambda_K(1-s)$. 人们猜想 $\zeta_K(s)$ 的非平凡零点位于直线 $\Re(s) = \frac{1}{2}$ 上 (广义 Riemann 假设).

\subsection*{有限域上仿射平面曲线的 zeta 函数}

考虑 $\mathbb{F}_p$ 上的非奇异仿射平面曲线 $C\colon f(X, Y) = 0$. 仿照 Dedekind 的定义, 定义
\[
\zeta(C, s) = \prod_{\mathfrak{p}}\frac{1}{1 - \Nm\mathfrak{p}^{-s}}, \qquad \Re(s) > 1,
\]
其中 $\mathfrak{p}$ 遍历 $\mathbb{F}_p[x, y] \overset{\text{def}}{=} \mathbb{F}_p[X, Y]/(f(X, Y))$ 中非零素理想, $\Nm\mathfrak{p}$ 是商环 $\mathbb{F}_p[x, y]/\mathfrak{p}$ 的阶.

因 $\mathbb{F}_p[x, y]/\mathfrak{p}$ 有限且是整环, 它是域, 记其 $\mathbb{F}_p$ 上的次数为 $\deg(\mathfrak{p})$. 则 $\Nm\mathfrak{p} = p^{\deg\mathfrak{p}}$, 这使我们在 zeta 函数中可作变量替换. 当我们定义\footnote{$Z$ 是大写 zeta.}
\[
Z(C, T) = \prod_{\mathfrak{p}}\frac{1}{1 - T^{\deg\mathfrak{p}}},
\]
则 $\zeta(C, s) = Z(C, p^{-s})$. 乘积 $\prod\frac{1}{1 - T^{\deg\mathfrak{p}}}$ 对所有小 $T$ 收敛, 但通常只把它视为形式幂级数
\[
Z(C, T) = \prod_{\mathfrak{p}}(1 + T^{\deg\mathfrak{p}} + T^{2\deg\mathfrak{p}} + \cdots)\in\mathbb{Z}[[T]].
\]

\begin{example} \label{thm:4.9.7} 
可把 $\mathbb{A}^1$ 实现为仿射平面曲线 $Y = 0$, 故 $\mathbb{F}_p[x, y] = \mathbb{F}_p[X]$. $\mathbb{F}_p[X]$ 上非零理想与首一多项式 $f\in\mathbb{F}_p[X]$ 一一对应, 且理想 $(f)$ 素当且仅当 $f$ 不可约. 故
\[
Z(\mathbb{A}^1, T) = \prod_{f\text{ 首一不可约}}\frac{1}{1 - T^{\deg f}}.
\]
两边取对数并求导:
\begin{align*}
    \frac{Z'(\mathbb{A}^1, T)}{Z(\mathbb{A}^1, T)} & = \sum_f\frac{\deg f\cdot T^{\deg f - 1}}{1 - T^{\deg f}} \\
    & = \sum_f\sum_{n\ge 0}\deg f\cdot T^{(n+1)\deg f - 1}.
\end{align*}
在此幂级数中, $T^{m-1}$ 的系数是 $\sum\deg f$ ($f$ 遍历次数整除 $m$ 的首一不可约多项式), 即 $p^m$. 故
\[
\frac{Z'(\mathbb{A}^1, T)}{Z(\mathbb{A}^1, T)} = \sum p^mT^{m-1}.
\]
积分，得
\begin{equation}\label{eq:40}
\log Z(\mathbb{A}^1, T) = \sum_m\frac{p^mT^m}{m} = \log\frac{1}{1-pT}.
\tag{40}
\end{equation}
因此，$Z(\mathbb{A}^1, T) = \frac{1}{1-pT}$.
\end{example}

\subsection*{用 $C$ 的点表达 $Z(C, T)$}

设 $N_m = \#\mathbb{A}^1(\mathbb{F}_{p^m}) = p^m$. 则 \eqref{eq:40} 可改写为 $\log Z(\mathbb{A}^1, T) = \sum\frac{N_mT^m}{m}$. 我们证明对 $\mathbb{F}_p$ 上每个非奇异仿射曲线也有类似公式.

设 $C\colon f(X, Y) = 0$ 为 $\mathbb{F}_p$ 上非奇异仿射曲线, $\mathbb{F}_p[x, y] = \mathbb{F}_p[X, Y]/(f)$. 如 $\mathbb{A}^1$ 的情形,
\begin{align*}
    \frac{Z'(C, T)}{Z(C, T)} &= \sum_{\mathfrak{p}}\frac{\deg\mathfrak{p}\cdot T^{\deg\mathfrak{p}-1}}{1 - T^{\deg\mathfrak{p}}} \\ & = \sum_{\mathfrak{p}}\sum_{n\ge 0}\deg\mathfrak{p}\cdot T^{(n+1)\deg\mathfrak{p}}/T.
\end{align*}
在此幂级数中，$T^{m-1}$ 的系数是 $\sum\deg\mathfrak{p}$，其中 $\mathfrak{p}$ 遍历 $\deg\mathfrak{p}$ 整除 $m$ 的非零素理想. 但 $\deg\mathfrak{p} = [\mathbb{F}_p[x, y]/\mathfrak{p}:\mathbb{F}_p]$，故条件 $\deg\mathfrak{p}\mid m$ 意味着存在同态 $\mathbb{F}_p[C]/\mathfrak{p}\hookrightarrow\mathbb{F}_{p^m}$；由于 $\mathbb{F}_p[C]/\mathfrak{p}$ 在 $\mathbb{F}_p$ 上可分，这样的同态恰有 $\deg\mathfrak{p}$ 个. 反之，每个同态 $\mathbb{F}_p[x, y]\to\mathbb{F}_{p^m}$ 都通过某个满足 $\deg\mathfrak{p}\mid m$ 的素理想所对应的商 $\mathbb{F}_p[x, y]/\mathfrak{p}$ 因式分解. 因此，上述幂级数中 $T^{m-1}$ 的系数是 $\mathbb{F}_p$-代数同态 $\mathbb{F}_p[x, y]\to\mathbb{F}_{p^m}$ 的个数. 这样的同态由 $x, y$ 的像 $a, b$ 决定；反之，一对 $(a,b)\in\mathbb{F}_{p^m}^2$ 所定义的同态
\[
P(X,Y)\mapsto P(a,b)\colon\mathbb{F}_p[X,Y]\to\mathbb{F}_{p^m}
\]
当且仅当 $f(a,b)=0$ 时通过商 $\mathbb{F}_p[x,y]$ 因式分解. 因此有自然的一一对应 $\{\mathbb{F}_p\text{-代数同态 }\mathbb{F}_p[C]\to\mathbb{F}_{p^m}\}\leftrightarrow C(\mathbb{F}_{p^m})$，且
\[
\log Z(C, T) = \sum_{m\ge 1}N_m\frac{T^m}{m}, \qquad N_m = \#C(\mathbb{F}_{p^m}).
\]

\begin{proposition} \label{thm:4.9.8} 
设 $Z(C, T)$ 为 $\mathbb{F}_p$ 上非奇异仿射平面曲线 $C$ 的 zeta 函数, 则
\[
Z(C, T) = \exp\Bigl(\sum_{m\ge 1}N_m\frac{T^m}{m}\Bigr).
\]
\end{proposition}

\begin{proof}
对前面的等式取 exp.
\end{proof}

\subsection*{函数域的 zeta 函数}

设 $C$ 为 $\mathbb{F}_p$ 上非奇异仿射平面曲线. 若 $\mathfrak{p}$ 是 $k[C]$ 中的素理想, 则 $k[C]_\mathfrak{p}$ 是 $k(C)$ 中带极大理想 $\mathfrak{p}k[C]_\mathfrak{p}$ 与剩余域 $k[C]_\mathfrak{p}/\mathfrak{p}k[C]_\mathfrak{p}\cong k[C]/\mathfrak{p}$ 的离散赋值环. 故 $\zeta(C, s) = \prod_R\frac{1}{1-(R:\mathfrak{m}_R)^{-s}}$, 其中 $R$ 遍历对应 $k[C]$ 中素理想的离散赋值环, $\mathfrak{m}_R$ 是 $R$ 中的极大理想.

设 $K$ 为 $\mathbb{F}_p$ 上一元函数域. 上面的讨论提示定义 $\zeta(K, s) = \prod_R\frac{1}{1-(R:\mathfrak{m}_R)^{-s}}$, 其中 $R$ 遍历 $K$ 中以 $K$ 为分式域的离散赋值环. 类似定义 $Z(K, T)$. 则例如
\begin{align*}
\zeta(\mathbb{F}_p(T), s) &= \frac{1}{(1-p^{-s})}\zeta(\mathbb{A}^1, s) = \frac{1}{(1-p^{-s})(1-p^{1-s})},\\
Z(\mathbb{F}_p(T), s) &= \frac{1}{1-T}Z(\mathbb{A}^1, T) = \frac{1}{(1-T)(1-pT)}.
\end{align*}

\subsection*{射影曲线的 zeta 函数}

设 $C$ 为 $\mathbb{F}_p$ 上非奇异射影曲线. 把 $C$ 的 zeta 函数定义为它的函数域的 zeta 函数. 例如,
\[
Z(\mathbb{P}^1, T) = \frac{1}{(1-T)(1-pT)}, \qquad
\zeta(\mathbb{P}^1, s) = \frac{1}{(1-p^{-s})(1-p^{1-s})}.
\]
如仿射情形, 有表达式
\[
Z(C, T) = \exp\Bigl(\sum_{m\ge 1}N_m\frac{T^m}{m}\Bigr), \qquad N_m = \#C(\mathbb{F}_{p^m}); \qquad
\zeta(C, s) = Z(C, p^{-s}).
\]
例如, $N_m(\mathbb{P}^1) = N_m(\mathbb{A}^1) + 1$, 故 $\log Z(\mathbb{P}^1, T) = \log Z(\mathbb{A}^1, T) + \log\frac{1}{1-T}$, 从而
\[
Z(\mathbb{P}^1, T) = \frac{1}{(1-T)(1-pT)}.
\]
类似地, 对 $\mathbb{F}_p$ 上椭圆曲线 $E$, $Z(E, T) = \frac{1}{1-T}Z(E_{\mathrm{aff}}, T)$, 其中 $E_{\mathrm{aff}}$ 是仿射曲线 $E\cap\{Z\neq 0\}$.

\subsection*{计数固定次数的正除子}

\begin{proposition} \label{thm:4.9.9} 
设 $E$ 为 $\mathbb{F}_p$ 上的椭圆曲线. 对所有 $m\ge 1$, 次数为 $m$ 的正除子恰好有 $\#E(\mathbb{F}_p)(p^m-1)/(p-1)$ 个.
\end{proposition}

\begin{proof}
固定除子 $D_0$, 令 $P(D_0)$ 为 $E$ 上满足 $D\sim D_0$ (即 $D = D_0 + (f)$ 对某个 $f\in\mathbb{F}_p(E)^{\times}$) 的所有正除子 $D$ 的集合. 则 $f\mapsto D_0 + (f)\colon L(D_0)\setminus\{0\}\to P(D_0)$ 是满射, 且两个函数有相同像当且仅当一个是另一个的常数倍. 故有双射 $(L(D_0)\setminus\{0\})/\mathbb{F}_p^{\times}\to P(D_0)$. 当 $\deg D_0 = m\ge 1$ 时, Riemann--Roch 定理 (I\ref{thm:1.4.13}) 表明 $L(D_0)$ 的维数为 $m$, 故 $P(D_0)$ 有 $(p^m-1)/(p-1)$ 个元素.

设 $\operatorname{Pic}(E) = \operatorname{Div}(E)/\{(f) \mid f\in k(E)^{\times}\}$——由 \ref{thm:4.1.10}, 这与 I, \S\ref{sec:1.4} 的定义一致. 因主除子的次数为零, 次数映射穿过 $\operatorname{Pic}(E)$: $\deg\colon\operatorname{Pic}E\to\mathbb{Z}$. 注意 $\deg$ 是满射, 因为 $\mathfrak{p}_\infty\mapsto 1$. 定义 $\operatorname{Pic}^mE = \{d\in\operatorname{Pic}(E) \mid \deg d = m\}$. 则:
\begin{enumerate}[label=\textup{(\alph*)}]
\item 映射 $\operatorname{Pic}^0(E)\to\operatorname{Pic}^m(E)$, $d\mapsto d + m\mathfrak{p}_\infty$, 是双射 (显然);
\item 存在双射 $E(k)\to\operatorname{Pic}^0(E)$ (I~\eqref{eq:8}).
\end{enumerate}
现在能计数 $E$ 上次数 $m$ 的正除子类. 上面看到 $\operatorname{Pic}^m(E)$ 中每个类有 $\frac{p^m-1}{p-1}$ 个元素, 且类数 $\#\operatorname{Pic}^m(E) = \#E(\mathbb{F}_p)$. 因此, 总共 $\#E(\mathbb{F}_p)\frac{p^m-1}{p-1}$ 个次数为 $m$ 的正除子.
\end{proof}

\subsection*{椭圆曲线 zeta 函数的有理性}

\begin{theorem} \label{thm:4.9.10} 
设 $E$ 为 $\mathbb{F}_p$ 上的椭圆曲线. 则
\[
Z(E, T) = \frac{1 + (N_1 - p - 1)T + pT^2}{(1-T)(1-pT)}, \qquad N_1 = \#E(\mathbb{F}_p).
\]
\end{theorem}

\begin{proof}
这里 $Z(E, T) = \frac{1}{1-T}Z(E_{\mathrm{aff}}, T) = \frac{1}{1-T}\prod_{\mathfrak{p}}\frac{1}{1 - T^{\deg\mathfrak{p}}}$, 其中 $\mathfrak{p}$ 遍历 $\mathbb{F}_p[C]$ 的素理想. 展开乘积, 得 $Z(E, T) = \sum d_mT^m$, 其中 $d_0 = 1$, $d_m$ 是次数 $m$ 的正除子数. 由命题 \ref{thm:4.9.9}, $d_m = N_1\frac{p^m-1}{p-1}$. 故
\[
Z(E, T) = 1 + \frac{N_1}{p-1}\Bigl(\frac{1}{1-pT} - \frac{1}{1-T}\Bigr)
= \frac{1 + (N_1 - p - 1)T + pT^2}{(1-T)(1-pT)}.
\]
\end{proof}

\begin{remark} \label{thm:4.9.11} 
(a) 写 $1 + (N_1 - p - 1)T + pT^2 = (1-\alpha T)(1-\beta T)$, 故 $\alpha, \beta$ 是满足 $N_1 - p - 1 = -\alpha - \beta$, $\alpha\beta = p$ 的代数整数. 则
\[
\log Z(E, T) = \log\frac{(1-\alpha T)(1-\beta T)}{(1-T)(1-pT)} = \sum_m\frac{(1 + p^m - \alpha^m - \beta^m)T^m}{m},
\]
{\emergencystretch=3em 故 $N_m(E) = 1 + p^m - \alpha^m - \beta^m$. 因此, 若知道 $N_1$, 就能求 $\alpha$ 与 $\beta$, 从而知道整个序列 $N_1(E), N_2(E), N_3(E), \dots$.
(b) 用 (a) 的记号, $\zeta(E, s) = \frac{(1-\alpha p^{-s})(1-\beta p^{-s})}{(1-p^{-s})(1-p^{1-s})}$. 它在 $s = 0$ 与 $s = 1$ 有单极点, 在 $p^s = \alpha$ 与 $p^s = \beta$ 处有零点. 写 $s = \sigma + it$, 则 $|p^s| = p^\sigma$, 而 $\zeta(E, s)$ 的零点实部为 $\frac{1}{2}$ 当且仅当 $\alpha$ 与 $\beta$ 的绝对值为 $p^{1/2}$.\par}
\end{remark}

按定义, $\alpha$ 与 $\beta$ 是多项式 $1 + bT + pT^2$ ($b = N_1 - p - 1$) 的逆根. 若 $b^2 - 4p\le 0$, 则 $\alpha$ 与 $\beta$ 是复共轭, 因它们的乘积是 $p$, 这蕴含它们各有绝对值 $p^{1/2}$. 反之, 若 $|\alpha| = p^{1/2} = |\beta|$, 则 $|N_1 - p - 1| = |\alpha + \beta|\le 2\sqrt{p}$. 故由定理 \ref{thm:4.9.10}, $E$ 的 Riemann 假设等价于 $|N_1 - p - 1|\le 2\sqrt{p}$.

\begin{exercise} \label{thm:4.9.12} 
证明 $\mathbb{F}_p$ 上椭圆曲线 $E$ 的 zeta 函数满足函数方程 $\zeta(E, s) = \zeta(E, 1-s)$.
\end{exercise}

\begin{exercise} \label{thm:4.9.13} 
计算曲线 $E\colon Y^2Z + YZ^2 = X^3 - X^2Z$ 在域 $\mathbb{F}_2, \mathbb{F}_3, \mathbb{F}_5, \mathbb{F}_7$ 上的 zeta 函数, 并在每种情形验证 Riemann 假设. 该曲线在含 $625$ 个元素的域上有多少个点?
\end{exercise}

\begin{exercise} \label{thm:4.9.14} 
(a) 设 $E$ 为椭圆曲线 $E\colon Y^2Z = X^3 - 4X^2Z + 16Z^3$. 对所有素数 $3\le p\le 13$ 计算 $N_p \overset{\text{def}}{=} \#E(\mathbb{F}_p)$ (若用计算机可更多).
(b) 设 $F(q)$ 为无穷乘积 $F(q) = q\prod_{n=1}^{\infty}(1-q^n)^2(1-q^{11n})^2 = q - 2q^2 - q^3 + 2q^4 + \cdots$ 给出的 (形式) 幂级数. 对 $n\le 13$ 计算 $F(q)$ 中 $q^n$ 的系数 $M_n$ (若用计算机可更多).
(c) 对每个素数 $p$, 计算和 $M_p + N_p$. 对 $M_p + N_p$ 一般应为多少提出猜想.
(d) 证明你的猜想. [这可能非常困难, 甚至用课程中覆盖的信息或许不可能.] (你的猜想是, 或至少应该是 Eichler--Shimura 定理的特殊情形. Wiles (及他人) 的重大结果是 Eichler--Shimura 定理适用于 $\mathbb{Q}$ 上所有椭圆曲线.)
\end{exercise}

\subsection*{zeta 的简史}

自本书第 1 版出版以来, Roquette 对特征 $p$ 下 Riemann 假设早期历史的详尽论述已可获取 (\cite{roquette2018}). 我在修改这些简短注释时取材于它及 \cite{oort2016}, 并向读者推荐两者. 关于``从 Weil 到今日''的故事, 见 \cite{milne2016}. 故事像数论中多数故事一样, 从 Gauss 开始.

\paragraph{Gauss 1814}\label{para:gauss1814}
1796 到 1814 年, Gauss 记了一本数学笔记本, 其最后一条写道:
\begin{quote}\small
若 $a + bi$ 是素数, 且 $a - 1 + bi$ 被 $2 + 2i$ 整除, 则同余方程
\[
1\equiv xx + yy + xxyy \pmod{a + bi},
\]
(含 $x = 1$, $y = \pm i$, $x = \pm i$, $y = 1$) 的解数等于 $(a-1)^2 + bb$.
\end{quote}
Gauss 的曲线 $X^2 + Y^2 + X^2Y^2 = 1$ 仿射且奇异, 但其相伴的非奇异曲线是椭圆曲线 $E\colon Y^2Z = X^3 + 4XZ^2$, 而 Gauss 的陈述蕴含: 当 $p$ 是满足 $p\equiv 1\pmod{4}$ 的素数时, $|N - p - 1| < 2p^{1/2}$ ($N = \#E(\mathbb{F}_p)$). 对其余奇素数, 易见 $N = p + 1$, 故 Gauss 是在断言 Riemann 假设对 $E$ 成立. Herglotz 于 1921 年给出了 Gauss 陈述的证明.\footnote{曲线 $E$ 有 $i$ 的复乘法, 而 Gauss 更精确的陈述可视为椭圆曲线复乘法理论的开始.} 见 \cite[3.4]{roquette2018}.

\paragraph{Emil Artin 1921\index{Artin}}
在学位论文中, Artin 研究了形如 $\mathbb{F}_p(X)[\sqrt{f(X)}]$ ($f(X)\in\mathbb{F}_p[X]$) 的函数域的代数数论, 并定义了它们的 zeta 函数. 他证明 zeta 函数是 $p^{-s}$ 的有理函数且满足函数方程, 并对约 40 个多项式 $f$ 检验其零点在直线 $R(s) = 1/2$ 上. Artin 于 1921 年完成学位论文. 同年他得到进一步结果, 未发表, 但写在给导师 Herglotz 的信中. 例如, 他证明了下面的判据:
\begin{quote}\small
设 $C$ 为 $\mathbb{F}_p$ 上仿射曲线, $N_m = \#C(\mathbb{F}_{p^m})$; $C$ 的 Riemann 假设成立当且仅当
\[
|N_m - p^m| < Ap^{m/2}, \qquad m = 1, 2, 3, \dots
\]
对某个与 $m$ 无关的常数 $A$ 成立.
\end{quote}

\paragraph{F.~K. Schmidt\index{Schmidt} 1925--1931}
Schmidt 对任意域 $k$ 上非奇异射影曲线 $C$ 证明了 Riemann--Roch 定理. 当 $k = \mathbb{F}_q$ 时, 他用它 (如我们对椭圆曲线所做) 证明
\[
Z(C, T) = \frac{P(T)}{(1-T)(1-qT)}, \qquad P(T) = 1 + \cdots\in\mathbb{Z}[T], \ \deg P(T) = 2g.
\]
且 $Z(1/qT) = q^{1-g}T^{2-2g}Z(T)$, 故 $\zeta(1-s) = q^{1-g}(q^{2-2g})^{-s}\zeta(s)$. 写 $P(T) = \prod(1-\alpha_jT)$. 则公式表明 $N_m = 1 + q^m - \sum\alpha_j^m$. 故一旦知道 $\alpha_1, \dots, \alpha_{2g}$, 就知道所有 $N_m$. 但不同于椭圆曲线情形, $N_1$ 不决定 $\alpha_j$——需要知道若干 $N_m$.

\paragraph{Davenport\index{Davenport}, Mordell, Hasse 1932}
Davenport 与 Mordell 感兴趣的问题是: 设 $f(X, Y)\in\mathbb{Z}[X, Y]$, $p$ 为使 $f$ 模 $p$ 几何不可约的素数; 求 $f(X, Y) = 0$ 在 $\mathbb{F}_p$ 中解数 $N$ 的如下形式的估计:
\begin{equation} \label{eq:41}
|N - p|\le Ap^\theta
\tag{41}
\end{equation}
其中 $A$ 是正常数, 指数 $\theta < 1$ 尽可能小; $\theta$ 与 $A$ 都不应依赖 $p$. 他们用特设方法得到若干结果. 例如对多项式 $aX^m + bY^n + c$, Davenport 证明估计 \eqref{eq:41} 对最优可能的 $\theta = 1/2$ 成立. Hasse 改进 Davenport 的结果, 证明把 $p$ 换成 $p^r$ 时估计对相同常数 $A$ 仍成立. 当被告知此陈述时, Artin 向 Hasse 解释它等价于曲线 $aX^m + bY^n + c = 0$ 的 Riemann 假设, 且事实上 Davenport--Mordell 问题在延拓到任意有限基域 $\mathbb{F}_{p^r}$ 时等价于特征 $p$ 的 Riemann 假设. 正是这促使 Hasse 对函数域 Riemann 假设产生兴趣.\footnote{鉴于此, 令人惊讶的是公式 $Z(C, T) = \exp(\sum N_mT^m/m)$ 在 1940 年前似乎没有出现在文献中 (\cite[2.3]{oort2016}).}

\paragraph{Hasse, Deuring, Roquette 1932--1951}
Hasse\index{Hasse} 证明了有限域上椭圆曲线的 Riemann 假设. 他的证明如我们所给, 基于椭圆曲线的自同态环. Deuring 正确地看到把 Hasse 的证明推广到任意亏格曲线的关键, 是把椭圆曲线的自同态环换成曲线与自身的对应环 (ring of correspondences). 但德国学派所用的函数域方法很不适合研究对应, 这需要在``双重域''上工作, 直到 1951 年 Hasse 的学生 Roquette\index{Roquette} 才循此路线完成了有限域上曲线 Riemann 假设的证明.

\paragraph{Weil 1940--1948}
1940 年, Weil 宣布对一切曲线的 Riemann 假设的证明, 即对 $1\le i\le 2g$ 有 $|\alpha_i| = p^{1/2}$. 他的证明假设存在任意域上的代数几何理论, 包括 Jacobi 簇与 阿贝尔簇, 而当时只在 $\mathbb{C}$ 上知道. Weil 用 1940 年代大部分时间发展所需的代数几何, 并在 1948 年出版的书中给出了曲线 Riemann 假设的详细证明.

\paragraph{Weil 1949}
Weil 对曲线、阿贝尔簇及一些其他特殊代数簇的 zeta 函数的了解, 使他陈述出著名的``Weil 猜想''. 对 $\mathbb{F}_p$ 上非奇异射影簇 $V$, 可 (如曲线那样) 定义
\[
Z(V, T) = \exp\Bigl(\sum_{m\ge 1}N_m\frac{T^m}{m}\Bigr), \qquad N_m = \#V(\mathbb{F}_{p^m}); \qquad
\zeta(V, s) = Z(V, p^{-s}).
\]
Weil 猜想 $Z(V, T)$ 是 $T$ 的有理函数, $\zeta(V, s)$ 满足函数方程, 且 $\zeta(V, s)$ 的``Riemann 假设''成立. Weil 知道有理性与函数方程将随任意域上簇的上同调理论的存在而推出, 该理论给出``正确的'' Betti 数并允许 Lefschetz 不动点公式与 Poincar\'e 对偶定理.\footnote{不清楚 Weil 是否相信这样的上同调应存在. Hasse 的结果 (II~\ref{thm:2.6.9}) 表明特征 $p$ 中簇上不存在系数在 $\mathbb{Q}$ (甚至 $\mathbb{Q}_p$ 或 $\mathbb{R}$) 中的这样的上同调.}

\paragraph{Dwork\index{Dwork} 1960}
Weil 陈述其猜想后, 人们开始寻找具有证明有理性与函数方程所需性质的公理化上同调. Dwork 出人意料地找到 $Z(V, T)$ 是 $T$ 的有理函数的``初等''证明.

\paragraph{Grothendieck 等 1963/64}
Grothendieck 定义了平展上同调, 并在 M.~Artin 与 Verdier 的帮助下充分发展它, 以证明 $Z(V, T)$ 有理且满足函数方程. Grothendieck, Artin 与 Verdier 的结果表明, 对 $\mathbb{F}_p$ 上非奇异射影簇 $V$,
\[
Z(V, T) = \frac{P_1(T)P_3(T)\cdots P_{2d-1}(T)}{(1-T)P_2(T)P_4(T)\cdots P_{2d-2}(T)(1-p^dT)},
\]
其中 $d = \dim V$, $P_i(T)\in\mathbb{Z}[T]$. 且 $Z(V, 1/p^dT) = \pm T^\chi p^{d\chi/2}Z(V, T)$, 其中 $\chi$ 是 $V\times V$ 中对角线的自交数.

\paragraph{Deligne 1973}
Deligne 用平展上同调证明了剩余的 Weil 猜想, 即 Riemann 假设:
\[
P_i(T) = \prod_j(1-\alpha_{ij}T), \qquad |\alpha_{ij}| = p^{i/2}.
\]
最后这个陈述说 $\zeta(V, s)$ 的零点在直线 $\Re(s) = \frac{1}{2}, \frac{3}{2}, \dots, \frac{2d-1}{2}$ 上, 极点在线 $\Re(s) = 0, 1, 2, \dots, d$ 上.

\begin{exercise} \label{thm:4.9.15} 
设 $C$ 为 $\mathbb{F}_q$ 上亏格 $g$ 的非奇异射影曲线.
\begin{enumerate}[label=\textup{(\alph*)}]
	\item 证明
	\[
	Z(C, T) = \frac{P(T)}{(1-T)(1-qT)}, \qquad P(T) = 1 + \cdots\in\mathbb{Z}[T], \ \deg P(T) = 2g.
	\]
	\item 证明 $Z\bigl(\frac{1}{qT}\bigr) = q^{1-g}T^{2-2g}Z(T)$, 由此推出若 $\alpha$ 是 $P$ 的逆根, 则 $q/\alpha$ 也是.
	\item (Artin 判据) 证明 $C$ 的 Riemann 假设成立当且仅当
	\[
	|N_m - q^m - 1|\le 2g\sqrt{q^m}, \qquad \forall m\ge 1,
	\]
{\emergencystretch=3em 	其中 $N_m = \#C(\mathbb{F}_{q^m})$.\par}
	\begin{hint}
{\emergencystretch=3em 		对 (a), 见定理 \ref{thm:4.9.10} 的证明. 对 (b), 令 $Z(C, T) = \prod_{\mathfrak{p}}\frac{1}{1-T^{\deg(\mathfrak{p})}} = \sum_{n\ge 0}a_nT^n$, $a_n = \sum_d\frac{q^n-1}{q-1}$ (其中 $d$ 遍历次数 $n$ 的除子类), 并应用 Riemann--Roch 定理. 对 (c), 用不等式证明 $\log P(T^{-1}) = \log\prod_{1\le i\le 2g}(1-\alpha_iT^{-1}) = \sum_{m\ge 1}\frac{(N_m - q^m - 1)T^{-m}}{m}$ 在 $q^{1/2}|T|^{-1} < 1$ 时收敛, 并推出 $P(T^{-1})$ 在此区域无零点.\par}
	\end{hint}
\end{enumerate}
\end{exercise}

\begin{remarkstar}
Hasse 与 Deuring 在 1930 年代的工作为任意特征下椭圆曲线的现代理论奠定了基础. 除了计算椭圆曲线可能有的自同态环 (II~\S\ref{sec:2.6}) 与 Hasse 证明 Riemann 假设 (本节) 之外, Deuring 还得到若干其他重要结果.
\end{remarkstar}

\section{Birch\index{Birch} 与 Swinnerton-Dyer\index{Swinnerton-Dyer} 猜想}\label{sec:4.10}

\subsection*{引言}

回到计算 $E(\mathbb{Q})$ 秩的问题. 我们纯代数的方法只提供秩的上界 (用 Selmer 群表示), 而上界与实际秩之差由神秘的 Tate--Shafarevich 群度量. 很难决定 $S^{(2)}(E/\mathbb{Q})$ 的某个元素是来自无限阶元素, 还是存活下来给出 $\Sha(E/\mathbb{Q})$ 的非平凡元素——事实上, 没有已被证明的算法做到这一点. 显然, 有另一个方法会有帮助.

设 $E$ 为 $\mathbb{Q}$ 上的椭圆曲线. 对每个好约化素数 $p$, 令
\[
N_p = \#E_p(\mathbb{F}_p),
\]
其中 $E_p$ 表示 $E$ 在 $\mathbb{F}_p$ 上的约化.

一个想法是: 也许 $E(\mathbb{Q})$ 的秩应与群 $E_p(\mathbb{F}_p)$ 的阶有关. 对好素数 $p$, 有约化映射 $E(\mathbb{Q})\to E_p(\mathbb{F}_p)$, 但一般地它远非单射或满射. 例如, 若 $E(\mathbb{Q})$ 无限, 则其核也无限; 若 $E(\mathbb{Q})$ 有限 (故阶 $\le 16$), 则对所有大 $p$, 它将不可能是满射 (因 $\#E_p(\mathbb{F}_p)\ge p + 1 - 2\sqrt{p}$).

五十年代末, Birch 与 Swinnerton-Dyer 有了如下想法:\footnote{更详细地, Birch 在 1956--57 学年于 Princeton 期间得知了``线性代数群的自然 Tamagawa 测度对 Siegel 关于二次型工作的优美重述'', 他们想知道对 $\mathbb{Q}$ 上椭圆曲线 $E$ 是否有类似现象. 作为 $E$ 的 Tamagawa 数的替代, 他们取 $E(\mathbb{Q}_p)$ 的 $p$-进测度对小于某个界 $P$ 的素数的乘积. 对好约化的 $p$, 测度是 $N_p/p$, 故这本质上就是 $f(P)$.} 若 $E(\mathbb{Q})$ 大, 这应迫使 $N_p$ 比通常大. 由于他们能用当时仅存的少数计算机之一, 他们得以计算检验这一点. 对大的数 $P$ (取决于计算机速度), 令
\[
f(P) = \prod_{p\le P}\frac{N_p}{p}.
\]
回忆 $N_p$ 近似 $p$. 他们的计算引导他们提出下面的猜想.

\begin{conjecture} \label{thm:4.10.1} 
对 $\mathbb{Q}$ 上每条椭圆曲线 $E$, 存在常数 $c(E)$ 使
\[
\lim_{P\to\infty}f(P) = c(E)\log(P)^r,
\]
其中 $r = \operatorname{rank}(E(\mathbb{Q}))$.
\end{conjecture}

可把猜想更简洁地写成
\[
f(P)\sim c(E)\log(P)^r \qquad \text{当 } P\to\infty.
\]
注意此猜想的非凡性质: 它预测可以从数列 $N_p$ 确定 $E(\mathbb{Q})$ 的秩. 且连同对误差项的估计, 它将提供求 $r$ 的算法.

Birch 与 Swinnerton-Dyer 实际能相当一致地用此猜想预测 $r$, 但他们发现 $f(P)$ 随 $P$ 增加剧烈振荡, 且似乎没有希望以小于 10\% 的误差找到 $c(E)$. 相反, 在 Davenport 与 Cassels 的提示下, 他们用 $E$ 的 zeta 函数重述其猜想.

\subsection*{$\mathbb{Q}$ 上簇的 zeta 函数}

设 $V$ 为 $\mathbb{Q}$ 上非奇异射影簇. 这样的簇是一族齐次多项式
\[
F(X_0, \dots, X_n)\in\mathbb{Q}[X_0, \dots, X_n]
\]
{\emergencystretch=3em 的零点集. 缩放每个这样的多项式使其系数在 $\mathbb{Z}$ 中且无公因子, 令 $\bar{F}(X_0, \dots, X_n)\in\mathbb{F}_p[X_0, \dots, X_n]$ 为缩放后多项式模 $p$ 的约化. 若多项式 $\bar{F}$ 定义 $\mathbb{F}_p$ 上非奇异簇 $V_p$, 则称 $V$ 在 $p$ 处有好约化, 或 $p$ 对 $V$ 为好. 对给定簇, 除有限多个素数外都是好的.\par}

对每个好素数 $p$, 有 zeta 函数
\[
Z(V_p, T) = \exp\Bigl(\sum_{m\ge 1}N_m\frac{T^m}{m}\Bigr), \qquad N_m = \#V_p(\mathbb{F}_{p^m}); \qquad
\zeta(V_p, s) = Z(V_p, p^{-s}),
\]
并定义
\[
\zeta(V, s) = \prod_p\zeta(V_p, s).
\]
因 $V_p$ 的 Riemann 假设成立, 通过与 $\prod_p\frac{1}{1-p^{-s}}$ 比较, 乘积对 $\Re(s) > \dim V + 1$ 收敛.

设 $\ast$ 为 $\mathbb{Q}$上的点, 即 $\ast = \mathbb{A}^0 = \mathbb{P}^0$. 对此簇, 所有素数都好, 且
\[
\log Z(\ast_p, T) = \sum\frac{1\cdot T^m}{m} = \log\frac{1}{1-T}.
\]
故
\[
\zeta(\ast, s) = \prod_p\frac{1}{1-p^{-s}},
\]
这正是 Riemann zeta 函数, 已是有趣的函数.

设 $V = \mathbb{P}^n$. 所有素数都好, 且
\[
\#\mathbb{P}^n(\mathbb{F}_q) = \frac{q^{n+1}-1}{q-1} = 1 + q + \cdots + q^n, \qquad q = p^m,
\]
由此推出
\[
\zeta(\mathbb{P}^n, s) = \zeta(s)\zeta(s-1)\cdots\zeta(s-n),
\]
其中 $\zeta(s)$ 是 Riemann zeta 函数.

\begin{conjecture}[Hasse--Weil] \label{thm:4.10.2} 
对 $\mathbb{Q}$ 上任何非奇异射影簇 $V$, $\zeta(V, s)$ 可解析延拓为整个复平面上的亚纯函数, 且满足联系 $\zeta(V, s)$ 与 $\zeta(V, d+1-s)$ 的函数方程, 其中 $d = \dim V$.
\end{conjecture}

该猜想被广泛相信为真, 但只在少数情形已知. 注意上面计算表明, 对 $V = \mathbb{P}^n$, 它由 $\zeta(s)$ 的类似陈述推出.

当然, $\zeta(V, s)$ 预期满足 Riemann 假设, 但这对任何簇 (甚至点) 都未知.

\begin{remarkstar}
据 Weil 的回忆 (Collected Papers, II, p.~529), Hasse 对 $\mathbb{Q}$ 上椭圆曲线定义了 Hasse--Weil zeta 函数, 并把此情形的 Hasse--Weil 猜想设为学位论文问题! 起初 Weil 怀疑该猜想, 但他对 $Y^m = aX^n + b$ 形式的曲线在数域上证明了它, 方法是把它们的 zeta 函数用 Hecke $L$-函数表示. 特别地, Weil 表明椭圆曲线 $Y^2 = aX^3 + b$ 与 $Y^2 = aX^4 + b$ 的 zeta 函数可用 Hecke $L$-函数表示, 并提出对一切带复乘法的椭圆曲线应同样成立. 这在 Deuring 的``优美系列''论文中被证明. Deuring 的结果被 Shimura 与 Taniyama 及 Weil 推广到一切带复乘法的 阿贝尔簇. 在另一个方向上, Eichler 与 Shimura 通过把模椭圆曲线\index{motuoyuanquxian@模椭圆曲线(modular elliptic curve)}的 zeta 函数等同于模形式的 Mellin 变换\index{Mellin-bianhuan@Mellin 变换(Mellin transform)}, 证明了模椭圆曲线的 Hasse--Weil 猜想. 见\hyperref[ch:5]{第五章}.
\end{remarkstar}

\subsection*{$\mathbb{Q}$ 上椭圆曲线的 zeta 函数}

设 $E$ 为 $\mathbb{Q}$ 上椭圆曲线, $S$ 为 $E$ 有坏约化的素数集. 由上面定义
\[
\zeta(E, s) = \prod_{p\notin S}\frac{1 + (N_p - p - 1)p^{-s} + p^{1-2s}}{(1-p^{-s})(1-p^{1-s})}
= \frac{\zeta_S(s)\zeta_S(s-1)}{L_S(s)},
\]
其中 $\zeta_S(s)$ 是省去 $S$ 中素数对应因子的 Riemann zeta 函数,
\[
L_S(E, s) = \prod_{p\notin S}\frac{1}{1 + (N_p - p - 1)p^{-s} + p^{1-2s}}.
\]
写 $1 + (N_p - p - 1)T + pT^2 = (1-\alpha_pT)(1-\beta_pT)$, 故
\[
L_S(E, s) = \prod_p\frac{1}{1-\alpha_pp^{-s}}\frac{1}{1-\beta_pp^{-s}}.
\]
乘积 $\prod_p\frac{1}{1-p^{-s}}$ 对 $\Re(s) > 1$ 收敛, 且 $|\alpha_p| = p^{1/2} = |\beta_p|$, 故通过与 $\prod_p\frac{1}{1-p^{1/2}p^{-s}}$ 比较, $L_S(E, s)$ 对 $\Re(s) > \frac{3}{2}$ 收敛.

给 $L_S(E, s)$ 补充坏素数的因子. 定义
\[
L_p(T) = \begin{cases}
1 - a_pT + pT^2, & \text{若 } p \text{ 好},\ a_p = p + 1 - N_p,\\
1 - T, & \text{若 } E \text{ 有分裂乘性约化},\\
1 + T, & \text{若 } E \text{ 有非分裂乘性约化},\\
1, & \text{若 } E \text{ 有加性约化}.
\end{cases}
\]
在四种情形中 $L_p(p^{-1}) = \frac{N_p}{p}, \frac{p-1}{p}, \frac{p+1}{p}, \frac{p}{p}$, 每种情形都等于 $\#E_p^{\mathrm{ns}}(\mathbb{F}_p)/p$, 其中 $E_p^{\mathrm{ns}}$ 是曲线模 $p$ 的非奇异部分 (见 II~\S\ref{sec:2.3} 的表). 定义
\[
L(E, s) = \prod_p\frac{1}{L_p(p^{-s})},
\]
其中乘积现遍历一切素数. $E$ 的导子 (conductor) $N_{E/\mathbb{Q}}$ 定义为 $\prod_{p\text{ 坏}}p^{f_p}$, 其中
\[
f_p = \begin{cases} f_p = 1, & \text{若 } E \text{ 在 } p \text{ 处有乘性约化},\\ f_p\ge 2, & \text{若 } E \text{ 在 } p \text{ 处有加性约化; } p\neq 2, 3 \text{ 时等于 } 2. \end{cases}
\]
当 $E$ 有加性约化时 $f_p$ 的精确定义有点复杂. 但有一个 Ogg 公式可用于所有情形的计算:
\[
f_p = \ord\nolimits_p(\Delta) + 1 - m_p,
\]
其中 $m_p$ 是 N\'eron 模型 (不计数重数) 不可约分支数, $\Delta$ 是 $E$ 的极小 Weierstrass 方程 \eqref{eq:11}, p.~\pageref{eq:11} 的判别式.

定义
\[
\Lambda(E, s) = N_{E/\mathbb{Q}}^{s/2}(2\pi)^{-s}\Gamma(s)L(E, s).
\]
对 $E$ 在数域上的情形可类似定义完备 $L$-函数 $\Lambda(E, s)$. 下面是椭圆曲线情形 Hasse--Weil 猜想的精确版本.

\begin{conjecture} \label{thm:4.10.3} 
函数 $\Lambda(E, s)$ 可解析延拓为整个复平面上的亚纯函数, 且满足函数方程
\[
\Lambda(E, s) = w_E\Lambda(E, 2-s), \qquad w_E = \pm 1.
\]
\end{conjecture}

关于$w_E$的具体取值, 甚至已有明确的计算法则.

如前面所述, 对带复乘法的曲线, 该猜想在 1950 年代早期由 Deuring 证明. 关键点是总存在与 Gauss 对曲线 $Y^2Z = X^3 + 4XZ^2$ 证明的类似公式 (见 p.~\pageref{para:gauss1814}) 的 $N_p$``公式'', 它允许把 $E$ 的 $L$-函数等同于 Hecke 已定义并研究过的一类 $L$-函数 (``Hecke $L$-函数''), 对其已知解析延拓与函数方程.

本书其余大部分篇幅讨论的一个远为重要的结果是下面的. 令
\[
\Gamma_0(N) = \left\{\begin{pmatrix}a & b\\ c & d\end{pmatrix}\in\SL_2(\mathbb{Z}) \;\middle|\; c\equiv 0\pmod{N}\right\}.
\]
则 $\Gamma_0(N)$ 作用在复上半平面上, 商 $\Gamma_0(N)\setminus\mathbb{H}$ 有 Riemann 曲面结构. $\mathbb{Q}$ 上椭圆曲线 $E$ 称为模的 (modular), 若对某个 $N$ 存在 Riemann 曲面的非常数映射 $\Gamma_0(N)\setminus\mathbb{H}\to E(\mathbb{C})$. Eichler 与 Shimura (在五十、六十年代) 对模椭圆曲线证明了猜想 \ref{thm:4.10.3} 的略弱形式 (见\hyperref[ch:5]{第五章}, \S\ref{sec:5.7}).

回忆曲线在 $p$ 处有半稳定约化 (semistable reduction) 若它在该处有好或乘性约化; 等价地, 若 $p^2$ 不整除导子. Wiles (与 Richard Taylor) 证明无平方因子导子的椭圆曲线是模的. 此工作被多位作者推广到 $\mathbb{Q}$ 上所有椭圆曲线 (见\hyperref[ch:5]{第五章}). 故猜想 \ref{thm:4.10.3} 对 $\mathbb{Q}$ 上所有椭圆曲线已知.

\begin{remark} \label{thm:4.10.4} 
	\begin{enumerate}[label=\textup{(\alph*)}]
		\item Deuring 的结果对数域上带复乘法的椭圆曲线有效. Eichler, Shimura, Wiles 等人的结果只对 $\mathbb{Q}$ 上椭圆曲线有效. 即使今天, 对数域上椭圆曲线的 $L$-函数所知甚少, 除了一些全实域.
		\item 两个结果都证明远多于猜想 \ref{thm:4.10.3} 的简单陈述——它们成功识别了函数 $\Lambda(E, s)$.
	\end{enumerate}
\end{remark}

\subsection*{Birch 与 Swinnerton-Dyer 猜想的陈述}

设 $E$ 为 $\mathbb{Q}$ 上椭圆曲线. 令
\[
Y^2Z + a_1XYZ + a_3YZ^2 = X^3 + a_2X^2Z + a_4XZ^2 + a_6Z^3
\]
为 $E$ 在 $\mathbb{Q}$ 上的极小 Weierstrass 方程, 并令
\[
\omega = \frac{dx}{2y + a_1x + a_3}.
\]
回忆 (\S\ref{sec:4.6}) 有典范 N\'eron--Tate 高度配对
\[
\langle\cdot,\cdot\rangle\colon E(\mathbb{Q})\times E(\mathbb{Q})\to\mathbb{R}, \qquad
\langle P, Q\rangle = \hat{h}(P+Q) - \hat{h}(P) - \hat{h}(Q).
\]

\begin{conjecture}[Birch 与 Swinnerton-Dyer] \label{thm:4.10.5} 
设 $r$ 为 $E(\mathbb{Q})$ 的秩, $P_1, \dots, P_r$ 为 $E(\mathbb{Q})$ 中 $\mathbb{Z}$-线性无关的元素; 则
\[
L(E, s)\sim\Bigl(\Omega\prod_{p\text{ 坏}}c_p\Bigr)\frac{[\Sha(E/\mathbb{Q})]\det(\langle P_i, P_j\rangle)}{(E(\mathbb{Q}):\sum_i\mathbb{Z}P_i)^2}(s-1)^r \qquad \text{当 } s\to 1,
\]
其中
\begin{itemize}
	\item $[\ast] = \ast$的阶(别处记为 $\#\ast$);
	\item $\Omega = \int_{E(\mathbb{R})}|\omega|$;
	\item $c_p = (E(\mathbb{Q}_p):E^0(\mathbb{Q}_p)).$
\end{itemize}
特别地, 猜想预测 $E(\mathbb{Q})$ 的秩等于 $L(E, s)$ 在 $s = 1$ 处零点的阶. 在文献中, 此陈述称为弱 BSD (weak BSD) 猜想\index{caixiang@猜想(conjecture)!ruo-BSD@弱 BSD(weak BSD)}\index{caixiang@猜想(conjecture)!Birch-Swinnerton-Dyer@Birch--Swinnerton-Dyer(Birch and Swinnerton-Dyer)}, 而 \ref{thm:4.10.5} 称为全 BSD (full BSD) 猜想.\index{caixiang@猜想(conjecture)!quan-BSD@全 BSD(full BSD)}\footnote{``弱''猜想的证明已足以赢得百万美元的千禧年奖.}
\end{conjecture}

\begin{remark} \label{thm:4.10.6} 
	\begin{enumerate}[label=\textup{(\alph*)}]
		\item 对模椭圆曲线, 猜想中除 Tate--Shafarevich 群外的所有项都可计算, 事实上可由 Pari 计算 (见 II, \S \ref{sec:2.7}).
		\item 商 \[\frac{\det(\langle P_i, P_j\rangle)}{(E(\mathbb{Q}):\sum_i\mathbb{Z}P_i)^2}\] 与 $P_1, \dots, P_r$ 的选取无关. 当它们选为 $E(\mathbb{Q})$ 模挠的基时, 商等于 \[\frac{D}{[E(\mathbb{Q})_{\mathrm{tors}}]^2},\] 其中 $D$ 是 N\'eron--Tate 配对的判别式.
		\item 积分 $\int_{E(\mathbb{Q}_p)}|\omega|$ 有意义, 事实上它等于 $(E(\mathbb{Q}_p):E^1(\mathbb{Q}_p))/p$. 确实 (见II~\ref{thm:2.2.7}) 有双射 $E^1(\mathbb{Q}_p)\cong p\mathbb{Z}_p$, 在此下 $\omega$ 对应 $\mathbb{Z}_p$ 上使 $\mathbb{Z}_p$ 测度为 $1$ (从而 $p\mathbb{Z}_p$ 测度为 $1/(\mathbb{Z}_p:p\mathbb{Z}_p) = 1/p$) 的 Haar 测度. 故
		\[
		\int_{E(\mathbb{Q}_p)}|\omega| = (E(\mathbb{Q}_p):E^1(\mathbb{Q}_p))\int_{E^1(\mathbb{Q}_p)}|\omega| = (E(\mathbb{Q}_p):E^1(\mathbb{Q}_p))/p = c_p\frac{N_p}{p}.
		\]
		对包含所有 $E$ 有坏约化素数的任何有限素数集 $S$, 定义
		\[
		L_S^*(s) = \Bigl(\prod_{p\in S\cup\{\infty\}}\int_{E(\mathbb{Q}_p)}|\omega|\Bigr)^{-1}\prod_{p\notin S}\frac{1}{L_p(p^{-s})}, \qquad
		L_p(T) = 1 + (N_p - p - 1)T + pT^2.
		\]
		回忆 $p = \infty$ 时 $\mathbb{Q}_p = \mathbb{R}$. 当 $p$ 好时,
		\[
		L_p(p^{-1}) = \frac{N_p}{p} = \int_{E(\mathbb{Q}_p)}|\omega|,
		\]
		故 $L_S^*(s)$ 在 $s = 1$ 附近的行为与满足条件的 $S$ 无关,\footnote{更精确地, 对任何两个这样的集合 $S, S'$, $\lim_{s\to 1}L_S^*(s)/L_{S'}^*(s) = 1$.} 而 Birch--Swinnerton-Dyer 猜想可陈述为
		\[
		L_S^*(E, s)\sim\frac{[\Sha(E/\mathbb{Q})]\cdot D}{[E(\mathbb{Q})_{\mathrm{tors}}]^2}(s-1)^r \qquad \text{当 } s\to 1.
		\]
	\end{enumerate}
\end{remark}

\begin{aside} \label{thm:4.10.7} 
形式地 $L_p(1) = \prod_p\frac{p}{N_p}$, 故猜想 \ref{thm:4.10.5} 与猜想 \ref{thm:4.10.1} 有一致的气息. 假设 $E$ 模, \cite{goldfeld1982} 证明: 若 $\prod_{p\le x}N_p/p\sim C(\log x)^{-r}$ 对某 $C > 0$ 与 $r$, 则 $L(E, s)\sim C'\cdot(s-1)^r$ 当 $s\to 1$, 且 $L(E, s)$ 的广义 Riemann 假设成立, 即对 $\Re(s) > 1$ 有 $L(E, s)\neq 0$. 不用说, 猜想 \ref{thm:4.10.1} 对单条曲线都没有证明. 事实上, 它比广义 Riemann 假设强: 对固定的模椭圆曲线 $E/\mathbb{Q}$, 下面陈述等价:
\begin{align*}
\prod_{p\le x}N_p/p&\sim C(\log x)^r && \text{对某非零 } C \text{ 与 } r;\\
\Psi_E(x) &= o(x\log x), && \text{其中 } \Psi_E(x) = \sum_{p\le x}(\alpha_p^k + \bar{\alpha}_p^k)\log p.
\end{align*}
$L$-函数 $L(E, s)$ 的广义 Riemann 假设等价于 $\Psi_E(x) = O(x(\log x)^2)$, 故 $\Psi_E(x) = o(x\log x)$ 可视为它的更深 (但仍合理的) 形式.
\end{aside}

\begin{aside} \label{thm:4.10.8} 
椭圆曲线与高维 阿贝尔簇的一个区别是 阿贝尔簇成对出现: 对每个 阿贝尔簇 $A$ 有对偶 阿贝尔簇 $A^\vee$ ($A$ 的 Picard 簇). 对椭圆曲线, $A$ 与 $A^\vee$ 典范同构. N\'eron--Tate 高度给出配对 $A^\vee(\mathbb{Q})\times A(\mathbb{Q})\to\mathbb{R}$, 而 Cassels--Tate 配对给出 $\Sha(A^\vee/\mathbb{Q})$ 与 $\Sha(A/\mathbb{Q})$ 之间的对偶. 对 $\mathbb{Q}$ 上 阿贝尔簇 $A$, Birch--Swinnerton-Dyer 猜想推广为陈述
\[
L_S^*(A, s)\sim\frac{[\Sha(A/\mathbb{Q})]\cdot D}{[A^\vee(\mathbb{Q})_{\mathrm{tors}}][A(\mathbb{Q})_{\mathrm{tors}}]}(s-1)^r \qquad \text{当 } s\to 1,
\]
其中 $D$ 是高度配对 $A^\vee(\mathbb{Q})\times A(\mathbb{Q})\to\mathbb{R}$ 的判别式. 我们有 $L_S^*(A^\vee, s) = L_S^*(A, s)$, 因为 $A^\vee$ 与 $A$ 等种, 且 $[\Sha(A^\vee/\mathbb{Q})] = [\Sha(A/\mathbb{Q})]$, 因为 Cassels--Tate 对偶, 故该陈述在 $A$ 与 $A^\vee$ 之间对称. 见 \cite{tate1966b}.
\end{aside}

\subsection*{关于 Birch--Swinnerton-Dyer 猜想的已知结果}

始于 Birch 与 Swinnerton-Dyer \cite{birch1963, birch1965} 的工作, 大量计算证据不断积累支持这些猜想: 对数千条曲线计算了猜想中除 $\Sha$ 外的所有项; 计算预测的 $[\Sha]$ 值总是平方; 每当计算出 $[\Sha]$ 的某个 $p$-主分量的阶, 它都与猜想一致.

对 $\mathbb{Q}$ 上一对等种椭圆曲线, 全 BSD 猜想中的大多数项对两条曲线不同, 但 \cite{cassels1965} 证明若猜想对一条曲线成立, 则对两者都成立, 即椭圆曲线的全 BSD 猜想有效性只依赖曲线的等种类. 对函数域上某些椭圆曲线, 该猜想已知 (见下节).

到 1970 年代中期, 在 $\mathbb{Q}$ 上证明全 BSD 猜想几乎没有进展. 正如 \cite{tate1974} 所说: ``这个非凡的猜想把函数 $L$ 在目前未知其有定义的点处的行为, 与目前未知有限的群 $\Sha$ 的阶联系起来.''

Coates 与 Wiles \cite{coates1977} 证明: 若 $E$ 有类数 $1$ 的虚二次域的复乘法且 $E(\mathbb{Q})$ 无限, 则 $L(E, 1) = 0$.

从现在起, 假设 $E/\mathbb{Q}$ 模. 故 $L(E/\mathbb{Q}, s)$ 延拓到整个复平面, 且满足函数方程 (猜想 \ref{thm:4.10.3}).

\cite{heegner1952} 是第一个找到系统构造数域上椭圆曲线无限阶点的方法的人. 模曲线含某些称为``特殊''的点, 它们对应带复乘法的椭圆曲线, 而这些点可被推前给出模椭圆曲线上的有理点. Heegner\index{Heegner} 的论文风格缺乏说服力, 且假设了很少有人掌握的 \cite{weber1908} 的知识, 故在很大程度上被忽视, 直到 Birch 在 1960 年代末接起它.

对模椭圆曲线 $E/\mathbb{Q}$ 与 $\mathbb{Q}$ 的复二次扩张 $K$, Birch \cite{birch1969a, birch1969b, birch1970, birch1975} 定义``Heegner 点'' $P_K\in E(K)$, 并建议它常应为无限阶. 基于数值证据, 他观察到 Heegner 点的高度似乎与某些椭圆曲线的 $L$-函数在 $1$ 处的值及其一阶导数有关 (Harvard 研讨会, 1973). 证实 Birch 的观察, Gross 与 Zagier \cite{gross1986} 证明: 若 $E$ 为 $\mathbb{Q}$ 上模椭圆曲线, 则
\[
L'(E/K, 1) = C\cdot\hat{h}(P_K), \qquad C\neq 0.
\]
故 $P_K$ 无限阶当且仅当 $L'(E/K, 1)\neq 0$. 此工作的历史说明见 \cite{birch2004}.

\cite{kolyvagin1988a, kolyvagin1988b} 表明椭圆曲线上的 Heegner 点系统控制 Selmer 群的大小与结构. 到约 1990 年, 已知 (感谢 Gross, Zagier, Kolyvagin 等人):
\begin{quote}\small
($*$) 若 $L(E/\mathbb{Q}, s)$ 在 $s = 1$ 处零点的阶至多为 $1$, 则弱 BSD 猜想对 $E$ 成立且 $\Sha(E/\mathbb{Q})$ 有限.
\end{quote}
事实上 Kolyvagin 表明更多. 在 ($*$) 的假设下, 为完成 $E$ 全 BSD 猜想的证明, 只需对所有 $p$ 证明 $\Sha(E/\mathbb{Q})$ 的 $p$-主分量有预测的阶. 用 Iwasawa 理论, Kato, Rubin, Skinner, Urban 等对许多 $E$ 与 $p$ 做到了.

没有关于 $L(E/\mathbb{Q}, s)$ 零点阶的假设, 所知甚少. 例如, 不知道如何证明 $\operatorname{rank}(E(\mathbb{Q}))\le 1$ 蕴含弱 BSD 猜想. 预印本 \cite{bhargava2019} 表明, 按高度排序时, $\mathbb{Q}$ 上椭圆曲线 $E$ 的同构类中的多数满足 ($*$) 的假设. 关于 $\mathbb{Q}$ 上椭圆曲线 BSD 猜想工作的近期综述, 见 \cite{coates2013, coates2016}.

\subsection*{同余数}\label{subsec:congruent}

回忆自然数 $n$ 若是有理边的直角三角形的面积, 则称为同余的 (congruent).\footnote{此名来自 Fibonacci, \textit{Liber Quadratorum} (平方之书), 1225.} 有两个基本问题.
\begin{enumerate}[label=\textup{(\alph*)}]
\item 是否存在判定给定 $n$ 是否同余的算法?
\item 所有模 $8$ 同余于 $5, 6$ 或 $7$ 的整数 $n$ 是否都是同余数 (如数值证据所暗示)?
\end{enumerate}
先解释同余数与椭圆曲线的关系. 设 $n$ 同余, 故存在正有理数 $a, b, c$ 使 $a^2 + b^2 = c^2$ 且 $ab/2 = n$. 则 $x = n(a+c)/b$, $y = 2n^2(a+c)/b^2$ 是椭圆曲线
\[
E(n)\colon Y^2 = X^3 - n^2X
\]
上满足 $y\neq 0$ 的点. 反之, 设存在点 $(x, y)\in E(n)(\mathbb{Q})$ 且 $y\neq 0$. 则
\[
a = |(n^2 - x^2)/y|, \qquad b = |2nx/y|, \qquad c = |(n^2 + x^2)/y|
\]
是面积为 $n$ 的直角三角形的边.

$E(n)$ 的射影闭包上仅有的有限阶有理点是 $(0:1:0)$, $(0:0:1)$, $(n:0:1)$ 与 $(-n:0:1)$, 它们都有阶 $2$. 我们得出结论:
\[
n \text{ 同余} \iff E(n) \text{ 有无限阶有理点}.
\]
在他 1952 年的论文中, Heegner 实际上对同余数感兴趣. 他证明 (用 Heegner 点): 若 $p$ 是模 $8$ 同余于 $5$ 或 $7$ 的素数, 则 $p$ 是同余数; 若 $p$ 模 $8$ 同余于 $3$ 或 $7$, 则 $2p$ 是同余数.

设 $n$ 模 $8$ 同余于 $5, 6$ 或 $7$. 则 $L(E(n), s)$ 函数方程中的常数 $w_E$ 是 $-1$, 故 $L(E(n), s)$ 在 $s = 1$ 处有零点. 因此, 弱 BSD 猜想蕴含 $n$ 同余.\index{tongyushu@同余数(congruent number)} 或者, 可直接从 Selmer 猜想证明 $n$ 同余 (\cite{stephens1975}). 已知在模 $8$ 同余于 $5, 6$ 或 $7$ 的无平方因子整数中, 是同余数的那些密度为 $1$ (\cite{kriz2020}).

若 $E(n)$ 的 Tate--Shafarevich 群有限, 则存在求 $E(n)$ 上有理点的算法, 从而判定 $n$ 是否同余. Tunnell 给出更显式的算法. 设 $n$ 奇且无平方因子; 若 $n$ 同余, 则满足 $2x^2 + y^2 + 8z^2 = n$ 的三元组数是满足 $2x^2 + y^2 + 32z^2 = n$ 的三元组数的两倍, 且逆命题在弱 BSD 猜想对 $E(n)$ 成立时真. 对偶 $n$ 有类似陈述. 关于同余数工作的近期综述, 见 \cite{coates2017}.

\section{椭圆曲线与球堆积}\label{sec:4.11}

Birch--Swinnerton-Dyer 猜想预期不仅对 $\mathbb{Q}$ 上椭圆曲线成立, 而且对全局域 ($\mathbb{Q}$ 或 $\mathbb{F}_p(T)$ 的有限扩张) 上一切椭圆曲线成立. 对函数域, 该猜想已在一些重要情形被证明, 而 Elkies\index{Elkies}, Shioda\index{Shioda} 等人用它表明这样得到的格 $(E(K), \langle\cdot,\cdot\rangle)$ 给出非常稠密的球堆积.\footnote{``最令人兴奋的发展之一是 Elkies 与 Shioda 从函数域上椭圆曲线的 Mordell--Weil 群构造格堆积. 这样的格在约 80 到 4096 维中比以往已知的任何堆积都更稠密.'' \cite[p.~xviii]{conway1999}.}

\subsection*{函数域上的 Birch--Swinnerton-Dyer 猜想}

设 $K$ 为 $\mathbb{F}_p(T)$ 的有限扩张, $\mathbb{F}_q$ 为 $\mathbb{F}_p$ 在 $K$ 中的代数闭包. 存在 $\mathbb{F}_q$ 上非奇异射影曲线 $C$ 使 $K = \mathbb{F}_q(C)$ (I\ref{thm:1.6.3}). 如 \S\ref{sec:4.9} 所讨论, $C$ 在 $\mathbb{F}_q$ 上的 zeta 函数形如
\[
Z(C, T) = \frac{\prod_{i=1}^{2g}(1-\omega_iT)}{(1-T)(1-qT)}, \qquad |\omega_i| = q^{1/2},
\]
其中 $g$ 是 $C$ 的亏格. 设 $E$ 为 $K$ 上常数椭圆曲线 (constant elliptic curve), 即由系数 $a_i\in\mathbb{F}_q\subset K$ 的方程
\[
Y^2Z + a_1XYZ + a_3YZ^2 = X^3 + a_2X^2Z + a_4XZ^2 + a_6Z^3
\]
定义的曲线. 把 $E$ 在 $\mathbb{F}_q$ 上的 zeta 函数写为
\[
Z(E, T) = \frac{(1-\alpha_1T)(1-\alpha_2T)}{(1-T)(1-qT)}.
\]

\begin{theorem} \label{thm:4.11.1} 
弱 Birch--Swinnerton-Dyer 猜想对 $E/K$ 成立.
\end{theorem}

\begin{proof}
此时弱猜想说 $E(K)$ 的秩 $r$ 等于满足 $1\le i\le 2$, $1\le j\le 2g$ 且 $\alpha_i = \omega_j$ 的序对 $(i, j)$ 的个数. $E(K)$ 的秩等于 $\operatorname{Hom}(J, E)$ 的秩, 其中 $J$ 是 $C$ 的 Jacobi 簇. \cite{tate1966a} 证明 $\operatorname{Hom}(J, E)$ 的秩等于满足 $\alpha_i = \omega_j$ 的序对 $(i, j)$ 的个数.
\end{proof}

\begin{theorem} \label{thm:4.11.2} 
全 Birch--Swinnerton-Dyer 猜想对 $E/K$ 成立.
\end{theorem}

\begin{proof}
此时全猜想说
\[
q^g\prod_{\alpha_i\neq\omega_j}\Bigl(1 - \frac{\omega_j}{\alpha_i}\Bigr) = [\Sha(E/K)]\cdot D,
\]
其中 $D$ 是配对
\[
\operatorname{Hom}(J, E)\times\operatorname{Hom}(E, J)\xrightarrow{(a,b)\mapsto b\circ a}\operatorname{End}(J)\xrightarrow{\text{trace}}\mathbb{Z}
\]
的判别式. 这在 \cite{milne1968} 中证明.
\end{proof}

更一般地, 对函数域上常数 阿贝尔簇, 弱 BSD 猜想是 \cite{tate1966a} 的定理, 全 BSD 猜想是 \cite{milne1968} 的定理. 且对函数域上椭圆曲线, 已知弱猜想蕴含全猜想 (\cite{milne1975, milne1981}). 弱 BSD 猜想对由系数 $a_i = a_i(T)$ 满足 $\deg a_i(T)\le 2i$ 的多项式的 Weierstrass 方程在 $\mathbb{F}_q(T)$ 上定义的椭圆曲线已知 (\cite{artin1973}).

设 $E_1$ 与 $E_2$ 为 $\mathbb{F}_q$ 上椭圆曲线, 把 $E_2$ 视为域 $K = \mathbb{F}_q(E_1)$ 上的椭圆曲线. 若 $E_1$ 与 $E_2$ 不等种, 则 Tate--Shafarevich 群 $\Sha(E/K)$ 的阶为 $(n_1 - n_2)^2$, 其中 $n_1 = \#E_1(k)$, $n_2 = \#E_2(k)$. 若 $E_1 = E_2 = E$ 非超奇异\index{tuoyuanquxian@椭圆曲线(elliptic curve)!chaoqiyi@超奇异(supersingular)}, 则 Tate--Shafarevich 群 $\Sha(E/K)$ 的阶为 $(\operatorname{End}(E):\mathbb{Z}[F])^2$, 其中 $F$ 是 Frobenius 映射. (由上面公式的练习.)

\subsection*{球堆积}

如评注 \ref{thm:4.6.4} 所注意, 由有限秩自由 $\mathbb{Z}$-模 $L$ 与 $V \overset{\text{def}}{=} L\otimes\mathbb{R}$ 上的正定二次型 $q$ 组成的序对极受关注. 可选 $V$ 的基把 $(V, q)$ 等同于 $(\mathbb{R}^r, X_1^2 + \cdots + X_r^2)$. 与 $q$ 相伴的双线性型是 $\langle x, y\rangle = q(x+y) - q(x) - q(y)$. 给定这样的序对 $(L, q)$, 需计算的数是
\begin{enumerate}[label=\textup{(\alph*)}]
\item $L$ 的秩 $r$;
\item 最短向量长度的平方 $m(L) = \inf_{v\in L, v\neq 0}\langle v, v\rangle$;
\item $L$ 的判别式 (discriminant) $\disc L = \det(\langle e_i, e_j\rangle)$, 其中 $e_1, \dots, e_r$ 是 $L$ 的基.
\end{enumerate}
判别式与 $L$ 的基的选取无关. 令 $\delta(L) = m(L)/\disc (L)^{1/r}$. $L$ 的基本平行六面体体积为 $\sqrt{\disc L}$. 与 $L$ 相伴的球堆积由半径为 $\frac{1}{2}\sqrt{m(L)}$ 的球组成, 故其密度为
\[
d(L) = 2^{-r}b_r\delta(L)^{r/2},
\]
其中 $b_r = \pi^{r/2}/\Gamma(\frac{r+2}{2})$ 是 $r$ 维单位球的体积. 为使 $d(L)$ 最大, 需使 $\delta(L)$ 最大.

设 $E$ 为域 $K = \mathbb{F}_q(C)$ 上常数椭圆曲线如前, $L = E(K)/E(K)_{\mathrm{tors}}$ 配二次型 $q = 2\hat{h}$. 若知道数 $\omega_i$ 与 $\alpha_j$, 则定理 \ref{thm:4.11.1} 给出 $r$, 而定理 \ref{thm:4.11.2} 给出 $\disc L$ 的上界:
\[
\disc L = q^g\cdot\prod_{\alpha_i\neq\omega_j}\Bigl(1 - \frac{\omega_j}{\alpha_i}\Bigr)\cdot[\Sha]^{-1}\le q^g\cdot\prod_{\alpha_i\neq\omega_j}\Bigl(1 - \frac{\omega_j}{\alpha_i}\Bigr).
\]
最后, 一个容易但非初等的论证\footnote{$E(K)$ 的元素 $P$ 定义映射 $u\colon C\to E$, 且 $\hat{h}(P)$ 与 $u$ 的次数有关. 故得到用 $\omega_i$ 与 $\alpha_j$ 表示的 $m(L)$ 下界.} 表明对所有有限 $k\supset\mathbb{F}_q$,
\[
m(L)\ge 2[C(k)]/[E(k)].
\]

\subsection*{例子}

考虑曲线
\[
C\colon X^{q+1} + Y^{q+1} + Z^{q+1} = 0
\]
在 $\mathbb{F}_{q^2}$ 上 (注意, 不是在 $\mathbb{F}_q$ 上).

\begin{lemma} \label{thm:4.11.3} 
	\begin{enumerate}[label=\textup{(\alph*)}]
		\item 曲线 $C$ 非奇异, 亏格 $g = \frac{q(q-1)}{2}$.
		\item $\#C(\mathbb{F}_{q^2}) = q^3 + 1$.
		\item $Z(C, T) = \frac{(1+qT)^{q(q-1)}}{(1-T)(1-q^2T)}$.
	\end{enumerate}
\end{lemma}

\begin{proof}
	\begin{enumerate}[label=\textup{(\alph*)}]
		\item 定义方程对 $X, Y, Z$ 的偏导数是 $X^q, Y^q, Z^q$, 而它们在 $\mathbb{P}^2$ 中无公共零点. 故曲线非奇异, 公式 \eqref{eq:6}, p.~\pageref{eq:6} 表明其亏格为 $q(q-1)/2$.
		\item 群 $\mathbb{F}_{q^2}^{\times}$ 循环, 阶为 $q^2-1 = (q+1)(q-1)$, 而 $\mathbb{F}_q^{\times}$ 是其 $q-1$ 阶子群. 故 $x^{q+1}\in\mathbb{F}_q^{\times}$, 且 $x\mapsto x^{q+1}\colon\mathbb{F}_{q^2}^{\times}\to\mathbb{F}_q^{\times}$ 是核为 $q+1$ 阶循环群的满射同态. 当 $x$ 遍历 $\mathbb{F}_{q^2}$ 时, $x^{q+1}$ 取值 $0$ 一次, 而 $\mathbb{F}_q^{\times}$ 中每个非零值 $q+1$ 次. 对 $y^{q+1}$ 与 $z^{q+1}$ 有类似说明. 可缩放 $X^{q+1} + Y^{q+1} + Z^{q+1} = 0$ 的每个解, 使 $x = 0$ 或 $1$.
		\begin{enumerate}[label=\textup{情形\arabic*:}]
			\item $x = 1$, $1 + y^{q+1}\neq 0$. $y$ 有 $q^2 - q - 1$ 种可能, 然后 $z$ 有 $q+1$ 种可能. 故 $(q^2 - q - 1)(q+1) = q^3 - 2q - 1$ 个解.
			\item $x = 1$, $1 + y^{q+1} = 0$. $y$ 有 $q+1$ 种可能, 然后 $z$ 有一种. 故 $q+1$ 个解.
			\item $x = 0$. 可取 $y = 1$, 然后 $z$ 有 $q+1$ 种可能.
		总计 $q^3 + 1$ 个解.
		\end{enumerate}
		\item 我们知道 \[\#C(\mathbb{F}_{q^2}) = 1 + q^2 - \sum_{i=1}^{2g}\omega_i.\] \footnote{校对注：原书此处写为 $\#C(\mathbb{F}_q)$，但曲线按本例在 $\mathbb{F}_{q^2}$ 上考虑，且右侧的 $1+q^2$ 以及 (b) 的点数均对应 $\mathbb{F}_{q^2}$.} 故 $\sum_{i=1}^{2g}\omega_i = q^2 - q^3 = -2gq$. 因 $|\omega_i| = q$, 这迫使 $\omega_i = -q$.
	\end{enumerate}
\end{proof}

对所有 $q$, 已知存在 $\mathbb{F}_{q^2}$ 上椭圆曲线 $E$ 使 $E(\mathbb{F}_{q^2})$ 有 $q^2 + 2q + 1$ 个元素 (Riemann 假设允许的最大值). 对这样的曲线,
\[
Z(E, T) = \frac{(1+qT)^2}{(1-T)(1-q^2T)}.
\]

\begin{proposition} \label{thm:4.11.4} 
设 $L = E(K)/E(K)_{\mathrm{tors}}$, 其中 $E$ 与 $K = \mathbb{F}_{q^2}(C)$ 如前.\footnote{校对注：原书此处的 $K$ 写为 $\mathbb{F}_q(C)$，但本例及证明 (c) 明确以 $\mathbb{F}_{q^2}$ 为基域，故相应函数域应为 $\mathbb{F}_{q^2}(C)$.} 则:
\begin{enumerate}[label=\textup{(\alph*)}]
	\item $L$ 的秩 $r$ 为 $2q(q-1)$;
	\item $m(L)\ge 2(q-1)$;
	\item $[\Sha(E/K)]\disc (L) = q^{q(q-1)}$;
	\item $\delta(L)\ge 2(q-1)/\sqrt{q}$.
\end{enumerate}
\end{proposition}

\begin{proof}
	\begin{enumerate}[label=\textup{(\alph*)}]
		\item 因所有 $\alpha_i$ 与 $\omega_j$ 都等于 $-q$, 由 (\ref{thm:4.11.2}a) 秩为 $2\cdot 2g = 2q(q-1)$.
		\item 我们有\footnote{校对注：原书上式漏了前文点数下界中的因子 $2$，并把已知的 $[E(\mathbb{F}_{q^2})]=(q+1)^2$ 写为 $(q^2+1)^2$. 此处按前文更正；在常数曲线的情形，高度来自映射次数，格范数取偶整数，故由严格下界 $m(L)>2(q-2)$ 得 $m(L)\ge 2(q-1)$.} \[m(L)\ge 2\frac{[C(\mathbb{F}_{q^2})]}{[E(\mathbb{F}_{q^2})]} = \frac{2(q^3 + 1)}{(q + 1)^2} > 2(q - 2).\]

		\item 这是 (\ref{thm:4.11.2}) 的特殊情形, 注意到我们的域是 $\mathbb{F}_{q^2}$ (不是 $\mathbb{F}_q$) 且 $g = q(q-1)/2$.
		\item 立即由前推得.
	\end{enumerate}
\end{proof}

\begin{remark} \label{thm:4.11.5} 
	\begin{enumerate}[label=\textup{(\alph*)}]
		\item \cite{dummigan1995} 在上面及密切相关的情形获得 Tate--Shafarevich 群的信息. 例如, 若 $q = p$ 或 $p^2$, 则 $\Sha(E/K)$ 为零; 若 $q = p^3$, 则其基数为至少 $p^{p^3(p-1)^3/2}$.
		\item 对 $q = 2$, $L$ 同构于记为 $D_4$ 的格; 对 $q = 3$, 同构于 Coxeter--Todd 格 $K_{12}$; 对 $q = 3$ 它与 Leech 格类似.
	\end{enumerate}
	椭圆曲线对格的更详细应用见 \cite{oesterle1990}.
\end{remark}

\begin{exercise} \label{thm:4.11.6} 
考虑曲线 $E\colon Y^2Z + YZ^2 = X^3$.
\begin{enumerate}[label=\textup{(\alph*)}]
	\item 证明 $E$ 是 $\mathbb{F}_2$ 上非奇异曲线.
	\item 计算 $\#E(\mathbb{F}_4)$, 其中 $\mathbb{F}_4$ 是含 $4$ 个元素的域.
	\item 设 $K$ 为整环 $\mathbb{F}_4[X, Y]/(X^3 + Y^3 + 1)$ 的分式域. 令 $L = E(K)/E(K)_{\mathrm{tors}}$, 把它视为配高度配对的格 $V = L\otimes\mathbb{R}$ 中的格. 计算 $L$ 的秩, $m(L)$ 与 $\delta(L)$.
\end{enumerate}
\end{exercise}
